% !TeX program = lualatex
\documentclass[a4paper,12pt,titlepage,oneside]{book}
\setlength{\parindent}{1.5em}  % Sangría del primer renglón
\setlength{\parskip}{2.5pt}      % Sin espacio extra entre párrafos
%\setstretch{1.5}               % Interlineado 1.5 (si no lo tienes en otro sitio)
% ------------ Paquetes básicos ------------
\usepackage[utf8]{inputenc}
\usepackage[table]{xcolor}
\usepackage{longtable}
\usepackage{booktabs}
\usepackage[ruled,vlined,spanish]{algorithm2e}
\SetKwInput{KwEntrada}{Entrada}
\SetKwInput{KwSalida}{Salida}
\SetKwFor{ParaCada}{para cada}{:}{fin}
\SetKwIF{lSi}{lSiNo}{lSiNoSi}{si}{entonces}{si no}{si no si}{fin si}
\SetKw{Devolver}{devolver}

\usepackage[T1]{fontenc}
\usepackage[spanish,es-noshorthands]{babel}
\usepackage{ragged2e}
\usepackage{pdflscape}
\usepackage{fontspec}
\setmainfont{Times New Roman}
\usepackage{float}
\usepackage{enumitem}
\setlist[itemize]{label=-}
\usepackage{graphicx}
\usepackage{caption}
\usepackage{etoolbox}
\usepackage{titling}
\captionsetup{justification=centering, singlelinecheck=false}
\usepackage{amsmath,amsthm,amssymb,amsfonts}
\usepackage{multirow}
\usepackage{rotating}
\usepackage{tabularx}
\usepackage{textcomp}
\usepackage{listings}
\usepackage{xcolor}
\usepackage{comment}
\usepackage{booktabs} % Para \toprule, \midrule, \bottomrule
\usepackage{caption}  % Para mejorar el formato de los subtítulo
\usepackage{unicode-math}
\usepackage{lipsum}
\usepackage{setspace}
\usepackage{sectsty}
\usepackage{titlesec}
\usepackage{tablefootnote}
\usepackage{tocloft}
\usepackage{adjustbox} % para ajustar figuras a márgenes
\usepackage[none]{hyphenat} % para no cortar en sílabas
\sloppy % permite a LaTeX espaciar un poco más libremente para evitar huecos


\renewcommand{\cftchapnumwidth}{5em} % Ajusta si se descuadra el índice
\renewcommand{\cftchappresnum}{Capítulo~}
\renewcommand{\cftchapaftersnum}{.}
\renewcommand{\cftchapfont}{\normalfont}

\usepackage{smartdiagram}
\usepackage{xcolor}
\usepackage{tikz}
\usetikzlibrary{arrows.meta,positioning,shapes.geometric,shapes.symbols,fit,backgrounds}
\usepackage{fancyhdr}
\usepackage{hyperref}
\hypersetup{
    colorlinks=true,
    linkcolor=black,
    urlcolor=blue,
    citecolor=black,
}
\newlistof{listequations}{loe}{Lista de fórmulas}
% ---- Estilo visual para lista de fórmulas
\renewcommand{\cftloetitlefont}{\bfseries\fontsize{14}{16}\selectfont}
\renewcommand{\cftafterloetitle}{\par\vskip 0.5ex}
\renewcommand{\cftlistequationspresnum}{Fórmula~}
\renewcommand{\cftlistequationsaftersnum}{.}
\renewcommand{\cftlistequationsfont}{\fontsize{12}{14}\selectfont}
\setlength{\cftlistequationsnumwidth}{4.5em}
\AtBeginEnvironment{table}{\fontsize{10}{12}\selectfont}
% ---------- Comando personalizado para insertar fórmulas con título ----------
\newcommand{\equationcaption}[2]{%
  \refstepcounter{equation}%
  \addcontentsline{loe}{listequations}{\protect\numberline{\theequation}#1}%
  \begin{equation}
    #2
  \end{equation}
}

% ========= FLAG para apéndices =========
\newif\ifinappendix
\inappendixfalse
\usepackage{etoolbox}
\pretocmd{\appendix}{\inappendixtrue}{}{}


\urlstyle{same}
\usepackage{chngcntr} % para notas a pie
\counterwithout{footnote}{chapter}
\tikzstyle{etapa} = [rectangle, rounded corners, minimum width=4cm, minimum height=1.2cm, text centered, draw=black, fill=blue!10]
\tikzstyle{proceso} = [ellipse, minimum width=3.5cm, minimum height=1.2cm, text centered, draw=black, fill=green!15]
\tikzstyle{flecha} = [thick,->,>=stealth]
\newcommand{\firstparagraph}{\begingroup\noindent}

% ------------ Bibliografía APA ------------
\usepackage[style=apa, sorting=nyt]{biblatex}
\DeclareLanguageMapping{spanish}{spanish-apa}
\addbibresource{references.bib}

% Make "et al" italic using xpatch
\usepackage{xpatch}
\xpatchbibmacro{name:andothers}{\bibstring{andothers}}{\emph{\bibstring{andothers}}}{}{}

% ------------ Márgenes ------------
\usepackage[a4paper,
left=1.5in,
right=1in,
top=1in,
bottom=1in]{geometry}

% ------------ Encabezado y pie de página ------------
\pagestyle{fancy}
\fancyhf{}
\fancyfoot[R]{\thepage}
\renewcommand{\headrulewidth}{0pt}
\renewcommand{\footrulewidth}{0pt}
\fancypagestyle{plain}{
    \fancyhf{}
    \fancyfoot[C]{\thepage}
}

% ------------ Sangrado de párrafos ------------
\setlength{\parindent}{1.5em}
\setstretch{1.5}

% ------------ Título de capítulo: "Capítulo X" encima del título ------------
% El label [display] aparece automáticamente para \chapter numerado
% y se suprime para \chapter* — sin flags, sin splash, sin duplicación
\titleformat{\chapter}[display]
  {\normalfont\bfseries\Huge}
  {\normalfont\bfseries\Huge Capítulo~\thechapter}
  {0.6em}
  {}

% ------------ Numeración hasta subsubsecciones ------------
\setcounter{secnumdepth}{3}
\setcounter{tocdepth}{3}

% ------------ Personalización listas ToC/LoF/LoT ------------
\renewcommand{\cfttoctitlefont}{\bfseries\fontsize{14}{16}\selectfont}
\renewcommand{\cftloftitlefont}{\bfseries\fontsize{14}{16}\selectfont}
\renewcommand{\cftlottitlefont}{\bfseries\fontsize{14}{16}\selectfont}
\renewcommand{\cftdotsep}{1}
\renewcommand{\cftchapdotsep}{\cftdotsep}
\renewcommand{\cftchapfont}{\normalfont}
\renewcommand{\cftchappagefont}{\normalfont}
\renewcommand{\cftchapleader}{\cftdotfill{\cftdotsep}}

\setlength{\cftchapindent}{0pt}
\setlength{\cftsecindent}{1em}
\setlength{\cftsubsecindent}{2.4em}
\setlength{\cftsubsubsecindent}{3.4em}

% ------------ Mostrar prefijos Figura/Tabla ------------
\renewcommand{\figurename}{Figura}
\renewcommand{\tablename}{Tabla}
\renewcommand{\listfigurename}{Lista de figuras}
\renewcommand{\listtablename}{Lista de tablas}
\renewcommand{\cftfigpresnum}{Figura~}
\renewcommand{\cftfigaftersnum}{.}
\renewcommand{\cfttabpresnum}{Tabla~}
\renewcommand{\cfttabaftersnum}{.}
\renewcommand{\cftfigfont}{\fontsize{12}{14}\selectfont}
\renewcommand{\cfttabfont}{\fontsize{12}{14}\selectfont}
\setlength{\cftfignumwidth}{4.8em}
\setlength{\cfttabnumwidth}{4.5em}
\renewcommand{\cftlistequationsfont}{\fontsize{12}{14}\selectfont}
\setlength{\cftlistequationsnumwidth}{5.5em}

% ------------ quote más pequeño ------------
\makeatletter
\renewenvironment{quote}
  {\list{}{\rightmargin\leftmargin
           \parsep\z@ \@plus\p@
           \leftmargin=1.5em
           \topsep=0.5em
           \partopsep=0pt}
   \item\relax
   \fontsize{\dimexpr\f@size pt - 2}{\dimexpr\f@baselineskip - 1pt}\selectfont}
  {\endlist}
\makeatother

% ------------ Documento ------------
\begin{document}

% ---------- Evitar numeración en portada ----------
\pagenumbering{gobble}

% ---------- Página de título ----------
\input{title page}

% ---------- Páginas preliminares ----------
\pagenumbering{roman}
\input{Front Matter Pages/acknowledgments}
\input{Front Matter Pages/acronyms}
\input{Front Matter Pages/abstract}
\input{Front Matter Pages/TOC_LOF_LOT}


% ---------- Activar numeración arábiga aquí ----------
\mainmatter
\pagenumbering{arabic}

% ---------- Capítulos ----------
\cleardoublepage
\pagenumbering{arabic}
\pagestyle{plain}

\chapter*{Introducción}
\addcontentsline{toc}{chapter}{\textnormal{Introducción}}
\label{chap:introduccion}

\setstretch{1.5}
\fontsize{12}{14}\selectfont
\vspace{\baselineskip}



No cabe duda de que, en el ámbito de la ingeniería lingüística, la terminología es una materia de gran importancia, tanto en su vertiente descriptiva como aplicada (\cite{cabre2004termintegral}). Bajo esta premisa, la disciplina se ha consolidado dentro de las ciencias del lenguaje contemporáneas como una materia transversal, cuya relevancia no se limita a la descripción de los lenguajes especializados, sino que ha avanzado a pasos agigantados hacia el desarrollo de herramientas digitales y computacionales para el tratamiento, gestión y sistematización del léxico de especialidad. A medida que los lenguajes técnicos se diversifican y evolucionan, se vuelve necesario no solo registrar las unidades que los constituyen, sino también comprender y dar cuenta de aquellas propiedades que hacen que una palabra sea término. En este contexto, la intersección entre teoría lingüística, lingüística de corpus y tecnologías del lenguaje ofrece un terreno fértil para explorar nuevas formas de detección, formalización y modelización de unidades terminológicas.

En esta tesis doctoral proponemos la extracción automática de unidades terminológicas mediante la combinación de enfoques lingüísticos, estadísticos y computacionales, con el objetivo de avanzar hacia una modelización teórica y computacional de la noción de 'unidad terminológica' o 'término'. Para ello, trabajamos sobre un corpus especializado en lingüística de corpus, por tratarse de un subdominio que presenta una elevada densidad terminológica, una rápida evolución conceptual y una notable interdisciplinariedad. Estas características lo convierten en un entorno idóneo tanto para la observación empírica de unidades terminológicas en contexto como para la validación de estrategias de extracción automática y abstracción formal. Además, la elección de este dominio facilita la integración de métodos computacionales con principios teóricos procedentes de la propia disciplina lingüística, lo que permite articular un enfoque genuinamente interfacial entre la teoría terminológica y las tecnologías del lenguaje.

El dinamismo y dispersión que caracterizan a la terminología lingüística han sido siempre un problema constante y de larga tradición para los lingüistas (Perea, 2000), parejo a su interrelación con distintos campos del saber y su amplia variedad de líneas de investigación, especialmente en los últimos años. La consolidación de la lingüística como ciencia ha desembocado en una pluralidad de términos al que ya se apunta desde finales del siglo pasado:

\begin{quote}
    El desarrollo de la terminología lingüística ha ido de la mano del desarrollo del pensamiento lingüístico; a medida que se crean y se especifican ideas, hipótesis y teorías, la terminología se pluraliza, aumenta; la pluralidad terminológica refleja la diversidad teórica o de teorías dentro de las concepciones lingüísticas en un doble sentido, por un lado, en el de la relación entre teorías y terminologías individuales o particulares y, por otro lado, en el de la relación entre uso de términos diferentes para hacer referencia a las mismas realidades y nociones. (\cite{caballero1994reflexiones})
\end{quote}

Desde su establecimiento como ciencia a principios del siglo XX, la lingüística ha contado con un gran número de diccionarios y glosarios de consulta que tratan de consignar su extensa terminología (\cite{laborda2017diccionarios}). Estas obras lexicográficas, sin embargo, presentan una serie de cuestiones problemáticas o que no llegan a satisfacer nuestros intereses desde un punto de vista terminológico (i.e. los métodos de adición de unidades terminológicas al diccionario, sin tener en cuenta enfoques con evidencias de corpus; la concepción de la lingüística en un sentido amplio, sin tener en cuenta ramas de la lingüística aplicada o léxicos especializados centrados en dichas ramas). La ausencia de algunos términos en la compilación de diccionarios es una cuestión que ya se señala como una de las razones por las que el trabajo terminográfico se base actualmente en corpus: “[t]he usual consultation of sources [...] do not always compile all terms that we need” (\cite{oliver2007}).

El planteamiento central de esta tesis doctoral sostiene que la combinación de enfoques lingüísticos, estadísticos y computacionales —incluyendo métodos de aprendizaje automático y aprendizaje profundo— permite mejorar significativamente la extracción automática de términos en ámbito especializados, tanto en recuperación como en precisión, frente a métodos tradicionales. Previo al desarrollo de la investigación, se plantean las siguientes hipótesis:

\begin{enumerate}
    \item  La integración sinérgica de enfoques lingüísticos y estadísticos en los sistemas automáticos de extracción de términos permite una mejora sustancial tanto en la recuperación como en la precisión de los términos identificados en corpus técnicos. A diferencia de los métodos más rígidos —ya sean puramente lingüísticos, centrados en patrones morfosintácticos y estructuras gramaticales, o puramente estadísticos, enfocados en medidas de asociación o frecuencia—, los enfoques híbridos permiten captar tanto las regularidades formales como las distribuciones empíricas de los términos en contexto. Esta combinación metodológica resulta especialmente eficaz en dominios técnicos y científicos donde la variación terminológica y la alta densidad informativa dificultan la extracción mediante técnicas unidimensionales. Asimismo, se argumenta que este enfoque integrador se alinea con las propiedades  de los términos, cuya identificación computacional requiere modelos que capturen simultáneamente su forma, uso y función discursiva.
    \item La caracterización automática de unidades terminológicas requiere incorporar información más allá de propiedades léxicas y morfosintácticas aisladas. Se sostiene que la integración de información relacional —derivada del comportamiento de la unidad en su entorno discursivo, sus asociaciones recurrentes y su funcionalidad dentro del texto especializado— mejora significativamente la detección y clasificación de términos frente a enfoques basados únicamente en patrones formales o frecuencias locales. Esta información relacional es capturada más eficazmente por modelos que operan a nivel de secuencia y contexto que por sistemas que procesan tokens de forma aislada.
    \item  Los modelos computacionales basados en aprendizaje automático y aprendizaje profundo, particularmente aquellos que incorporan representaciones distribucionales de palabras (como \textit{word embeddings}, \textit{contextual embeddings} o modelos \textit{Transformer}), superan en eficacia a los métodos tradicionales de extracción de términos, tanto en tareas de identificación como en clasificación y normalización terminológica. Esta superioridad se manifiesta especialmente cuando los modelos se entrenan sobre datos lingüísticamente enriquecidos, es decir, textos anotados con etiquetas morfosintácticas, funciones discursivas, límites terminológicos y otras propiedades lingüísticas relevantes. Se argumenta que estas técnicas permiten modelar computacionalmente la noción de término como una construcción situada, sensible al contexto y potencialmente polimorfa, lo cual no puede lograrse mediante métodos puramente formales o estadísticos. Además, el aprendizaje automático facilita la incorporación dinámica de nuevas evidencias lingüísticas durante el entrenamiento, lo que resulta crucial en dominios altamente especializados y en evolución.
    \item  Los resultados obtenidos mediante técnicas avanzadas de extracción automática de términos van más allá de su valor instrumental o auxiliar —al facilitar la creación de glosarios, índices o bases de datos terminológicas—, sino que también permiten generar inferencias de orden teórico sobre la naturaleza, el comportamiento y la organización de los términos en uso real. En particular, se argumenta que la modelización computacional de patrones terminológicos basada en datos empíricos posibilita la formulación de hipótesis lingüísticas más sólidas sobre los mecanismos de formación, variación, especialización y difusión de los términos en contextos técnicos y científicos. Esta retroalimentación entre el análisis computacional y la reflexión teórica permite el desarrollo de una terminología computacional con fundamento lingüístico, en la que los modelos automáticos no solo operan como herramientas de detección, sino como instrumentos epistémicos para la comprensión profunda del fenómeno terminológico. De este modo, se propicia una integración entre teoría y práctica que fortalece el carácter científico y aplicable de la disciplina terminológica.
    \item  El grado de representatividad terminológica de los corpus empleados en la extracción automática de términos constituye un factor determinante en la calidad y fiabilidad de los resultados obtenidos, hecho fundamentado en la premisa de que la distribución y densidad de los términos en un corpus está directamente condicionada por su capacidad de reflejar de forma equilibrada y exhaustiva los distintos subdominios, géneros discursivos y niveles de especialización del ámbito temático en cuestión. Se sostiene que corpus deficientemente representativos —ya sea por sesgo temático, escasez de fuentes, baja diversidad textual o insuficiente cobertura léxica— generan un sesgo en la identificación automática de términos, afectando tanto a la variedad como a la pertinencia de las unidades extraídas. En consecuencia, se argumenta la necesidad de aplicar métricas compuestas de representatividad terminológica (como el índice de cobertura terminológica o el β de Heaps) que permitan evaluar y garantizar la calidad del corpus antes de su explotación computacional.
    \item  La incorporación de mecanismos de explicabilidad en los modelos computacionales empleados para la extracción de términos es esencial para garantizar la interpretabilidad lingüística de los resultados y, por tanto, su integración en una terminología computacional con fundamento teórico. A diferencia de los modelos puramente “caja negra”, los enfoques explicables permiten rastrear los factores que conducen a la identificación de una unidad como término, ya sea a través de rasgos morfosintácticos, medidas estadísticas o patrones contextuales aprendidos. Esta capacidad explicativa no solo favorece la validación manual y la confianza en el sistema, sino que también facilita la retroalimentación entre los resultados computacionales y la teoría terminológica, permitiendo afinar conceptos como variación, saliencia o grado de terminologización. Por ello, se defiende que los sistemas de extracción deben diseñarse con criterios de transparencia, trazabilidad y replicabilidad.
\end{enumerate}

Para la comprobación de las hipótesis expuestas, se proponen los siguientes objetivos generales y específicos, estructurados en torno a los tres ejes fundamentales que articulan esta investigación: la formalización teórico-computacional del concepto de término, el desarrollo y evaluación de métodos automáticos para su detección, y la sistematización de los resultados en un recurso terminológico estructurado. Los objetivos propuestos requieren el diseño y la aplicación de un flujo metodológico secuencial, tal y como se describe en la Figura \ref{fig:flujo_simple_colores}.

Cada objetivo general se acompaña de un conjunto de objetivos específicos que orientan las etapas del trabajo, desde la construcción del corpus hasta la validación de los modelos y la explotación de los resultados, con el fin de garantizar la coherencia entre el planteamiento teórico, la metodología empleada y los resultados obtenidos.

\begin{itemize}
    \item \textbf{O. G. 1.}: Formalizar computacionalmente la noción de término especializado a partir de evidencias empíricas observables en corpus representativos y lingüísticamente anotados.
    \begin{itemize}
        \item \textit{O. E. 1. 1}: Diseñar y construir un corpus especializado representativo del dominio de la lingüística, segmentado y enriquecido con anotaciones morfosintácticas, discursivas y terminológicas, y evaluar su representatividad terminológica y adecuación a la tarea terminográfica propuesta.
        \item \textit{O. E. 1. 2.}: Identificar, sistematizar y operacionalizar un conjunto de rasgos lingüísticos, sintáctico-semánticos, contextuales y pragmáticos relevantes para la caracterización de las unidades terminológicas.
        \item \textit{O. E. 1. 3.}: Elaborar un modelo teórico-computacional que integre los rasgos seleccionados en una abstracción formal del término, compatible con tareas de extracción y análisis automatizado.
    \end{itemize}
    \item \textbf{O. G. 2.}: Desarrollar, entrenar y evaluar sistemas automáticos de detección de términos especializados, basados en metodologías lingüísticas, estadísticas, híbridas y de aprendizaje supervisado y profundo.
    \begin{itemize}
        \item \textit{O. E. 2. 1.}: Implementar una batería de métodos de extracción terminológica que incluya enfoques basados en reglas, medidas estadísticas clásicas, métodos híbridos y modelos de aprendizaje automático y profundo.
        \item \textit{O. E. 2. 2}.: Diseñar un protocolo de validación que contemple la exclusión sistemática de términos del conjunto de entrenamiento, con el fin de medir la capacidad de generalización y la detección de términos no vistos.
        \item \textit{O. E. 2. 3.}: Evaluar el rendimiento de los modelos mediante métricas estándar en el ámbito del procesamiento del lenguaje natural, complementadas con indicadores de coherencia terminológica y explicabilidad del modelo.
        \end{itemize}
    \item \textbf{O. G. 3.}: Construir un glosario digital de términos especializados contextualizados, estructurado a partir de los resultados del sistema de extracción y orientado a aplicaciones en investigación terminológica.
    \begin{itemize}
        \item \textit{O.E . 3. 1.}: Diseñar un modelo de entrada terminológica que integre definición, variantes, ejemplos en contexto, relaciones semánticas y metadatos lingüísticos derivados del análisis computacional.
        \item \textit{O. E. 3. 2.}: Generar el glosario a partir de las unidades terminológicas extraídas y validadas, incorporando mecanismos de trazabilidad hacia el corpus fuente y el modelo utilizado.
        \item \textit{O. E. 3. 3.}:  Vincular el glosario con los resultados del modelo y con los rasgos lingüístico-computacionales que motivaron su identificación, garantizando así la trazabilidad terminológica, la replicabilidad metodológica y la posibilidad de ampliar el recurso.
    \end{itemize}
\end{itemize}

\begin{figure}[H]
\centering
\begin{tikzpicture}[node distance=4.2cm, every node/.style={font=\small}]

% Etapas principales con colores diferentes
\node (modelo) [rounded corners=3pt, minimum width=4cm, minimum height=1.2cm, align=center, fill=violet!20] {Modelo teórico-\\computacional};
\node (extraccion) [right of=modelo, rounded corners=3pt, minimum width=4cm, minimum height=1.2cm, align=center, fill=teal!20] {Extracción automática\\de términos};
\node (glosario) [right of=extraccion, rounded corners=3pt, minimum width=4cm, minimum height=1.2cm, align=center, fill=orange!20] {Glosario computacional\\contextualizado};

% Flechas principales
\draw[flecha] (modelo) -- (extraccion);
\draw[flecha] (extraccion) -- (glosario);

\end{tikzpicture}
\caption{Flujo metodológico secuencial de la investigación}
\label{fig:flujo_simple_colores}
\end{figure}

De este modo, la tesis articula un recorrido metodológico que no solo avanza linealmente desde la teoría hacia la aplicación, sino que incorpora mecanismos de retroalimentación entre los resultados obtenidos y la formalización conceptual, permitiendo una revisión progresiva del modelo a partir de los datos empíricos. Este enfoque integrador pone de manifiesto la estrecha relación entre teoría y datos, al defender que, si bien una base teórica sólida es condición necesaria para la eficacia de los modelos computacionales, los datos empíricos también desempeñan un papel activo en la reformulación de hipótesis y en la comprensión profunda del lenguaje y las lenguas. En última instancia, esta interacción dinámica responde a uno de los fines centrales de la investigación lingüística: observar, describir y explicar el comportamiento real del lenguaje y las lenguas a partir de su uso, combinando intuición teórica y evidencia cuantitativa. Esta articulación resulta aún más pertinente en el desarrollo de modelos computacionales, donde el objetivo no es únicamente detectar regularidades, sino también formalizar, explicitar y modelar de forma sistemática determinados fenómenos lingüísticos a partir de datos observables.

La tesis se estructura en seis capítulos, además la introducción y las conclusiones finales. El Capítulo \ref{chap:1_terminologia} está dedicado al estudio de la terminología desde una orientación lingüística, con especial atención a la noción de término o unidad terminológica, su variabilidad, su comportamiento discursivo en textos especializados y su representación en productos lexicográficos y terminográficos. El Capítulo \ref{chap:2_tecnologias} contextualiza la investigación en el marco de las tecnologías del lenguaje, abordando tres áreas clave —lingüística computacional, lingüística de corpus e ingeniería lingüística— cuya intersección resulta esencial para el desarrollo de modelos automáticos de extracción terminológica. En el caso de la lingüística de corpus, este capítulo permite, además, delimitar conceptualmente el dominio sobre el que se construye el corpus de trabajo. El Capítulo \ref{chap:3_extraccion} expone la extracción automática de términos, presentando sus fundamentos, los principales enfoques existentes en la literatura y su relación con los desafíos terminológicos descritos previamente. La metodología de trabajo se detalla en el Capítulo \ref{chap:4_metodologia}, incluyendo la construcción y anotación del corpus, la definición de rasgos, la implementación de modelos y los criterios de evaluación empleados. En el Capítulo \ref{chap:5_resultados} se presentan y analizan los resultados obtenidos, tanto en términos de rendimiento de los sistemas como de coherencia terminológica de las unidades detectadas. Finalmente, se recogen las conclusiones generales, se discuten las implicaciones teóricas y metodológicas derivadas del presente estudio y se proponen líneas futuras de investigación.   % ← Usa \chapter* + \addcontentsline dentro
\input{chapters/chapter 1}
\input{chapters/chapter 2}
% ============================
% CAPÍTULO 3: METODOLOGÍA
% ============================

\phantomsection
\renewcommand{\tablename}{Tabla}
\setstretch{1.5}
\fontsize{14}{16}\selectfont

\chapter{Metodología de la investigación}
\label{chap:4_metodologia}

\makeatletter
\@afterindentfalse\@afterheading
\makeatother

- quitar el S de párrafo y poner véase lo discutido
- explícito- tiene que ver con el contexto de explícito (definitorio, metaling.)
- decir que las fuentes están en el anexo
- actualidad > mejor delimitación sincrónica
- explicar figura 4.3
- escáner -> mencionar al ILA
- separar términos y entidades

\setcounter{section}{0}
\renewcommand{\thesection}{\thechapter.\arabic{section}}

\fontsize{12}{14}\selectfont
\vspace{\baselineskip}

\phantomsection
\setstretch{1.5}
\begin{quote}

    [L]a lingüística es una ciencia empírica cultural y, en consecuencia, debe respetar todas las características del conocimiento científico y, al tiempo, reconocer las peculiaridades que derivan de la naturaleza cultural y social de su objeto de estudio, que son el lenguaje y las lenguas. (\cite{rojo2021introduccion})
\end{quote}

% Chapter 4.1
\phantomsection
\setstretch{1.5}
\fontsize{12}{14}\selectfont

\section{Resumen}

Este capítulo describe de manera detallada el diseño metodológico de la investigación, estableciendo los procedimientos, decisiones y criterios que han guiado la construcción del corpus, su anotación terminológica, la evaluación de su representatividad y la aplicación de distintos métodos de extracción automática de términos. La metodología se articula en torno a un corpus técnico-documental concebido como eje central que conecta la modelización teórica del término con su identificación automática y su posterior explotación terminográfica.

En primer lugar, se presenta el corpus como base empírica del estudio, situándolo como la fuente de datos más natural y menos monitorizada para el análisis del léxico de especialidad. A partir de una definición operativa de \textit{corpus} y de los principios clásicos de representatividad, equilibrio y adecuación al objetivo de investigación, se diseña un corpus especializado, metalingüístico, bilingüe (inglés–español) y comparable, orientado específicamente al dominio de la lingüística de corpus. El proceso de construcción del corpus se basa en una recopilación controlada de textos académico-científicos procedentes de diversas fuentes especializadas y repositorios bibliográficos. Este proceso combina materiales en acceso abierto, acceso institucional y documentos digitalizados, garantizando tanto la calidad científica de las fuentes como la trazabilidad y reproducibilidad del corpus resultante. 

El capítulo ofrece asimismo una descripción cuantitativa exhaustiva del corpus, atendiendo a la distribución de documentos, tokens y tipologías textuales en cada subcorpus. Este análisis permite identificar asimetrías estructurales entre el inglés y el español y evaluar sus posibles implicaciones para las tareas de extracción terminológica y aprendizaje automático, distinguiendo entre escenarios exploratorios y experimentos supervisados. Lejos de entender estas diferencias como limitaciones, se argumenta que reflejan la producción científica real del dominio y que, siempre que estén debidamente documentadas, no comprometen la validez del corpus para los fines propuestos. En lo que respecta a la anotación, el capítulo describe un proceso en dos fases que combina una primera anotación automática basada en términos semilla con una anotación manual exhaustiva realizada mediante una herramienta especializada. La anotación terminológica se apoya en un conjunto explícito de criterios lingüísticos y extralingüísticos que permiten identificar unidades terminológicas, dando lugar a un corpus anotado que asimismo proporciona el conjunto de datos necesario para el entrenamiento y evaluación de modelos supervisados de extracción automática de términos.

Una parte central del capítulo se dedica a la evaluación de la representatividad terminológica del corpus. Desde una concepción operativa de la representatividad, se propone un marco multimétrico que integra medidas de cobertura y saturación léxica, aplicadas específicamente a la terminología anotada. A partir de estas métricas se define y formaliza un Índice Compuesto de Representatividad Terminológica (ITR), que sintetiza el grado de adecuación del corpus en una puntuación única, interpretable y adaptable a distintos objetivos analíticos. Los resultados obtenidos permiten afirmar que el corpus es representativo desde un punto de vista terminológico, lo que respalda la generalización de los análisis y experimentos realizados sobre él.

Finalmente, el capítulo introduce la fase de extracción automática de términos, describiendo los distintos enfoques metodológicos considerados —lingüísticos, estadísticos, híbridos y basados en aprendizaje automático— y deteniéndose especialmente en los métodos supervisados. En este contexto, se detalla el proceso de ingeniería de características y el diseño de un conjunto de rasgos lingüísticos y estadísticos que permiten modelizar computacionalmente el comportamiento de los términos en el corpus. Esta sección sienta las bases técnicas y metodológicas para los experimentos de extracción automática y evaluación que se desarrollarán en los capítulos posteriores, así como para la construcción del glosario digital bilingüe.



Complementariamente, el capítulo también aborda la extracción automática de términos como un procedimiento orientado a obtener listas de candidatos a término mediante técnicas computacionales, subrayando su función de prerrequisito en múltiples aplicaciones (p. ej., compilación de recursos terminográficos, indexación, traducción especializada o minería de textos, entre otras). Desde el punto de vista conceptual, se describen las nociones de 	extit{unithood} y 	extit{termhood}, adaptadas como 	extit{grado de estabilidad} y 	extit{grado de especialización}, respectivamente, y se expone su relevancia para la delimitación de unidades terminológicas y para la construcción de criterios de filtrado y ponderación. Se sistematizan asimismo los principales enfoques metodológicos en la literatura: métodos estadísticos, métodos lingüísticos, métodos híbridos y métodos basados en aprendizaje automático, incluyendo las más recientes aproximaciones de aprendizaje profundo. Se introduce el uso de un 	extit{Gold Standard} como referencia para la evaluación objetiva y se formalizan las métricas habituales en el PLN (precisión, recuperación/exhaustividad y F1), junto con los conceptos de 	extit{ruido} y 	extit{silencio} como manifestaciones operativas de los falsos positivos y falsos negativos.
\section{El corpus técnico-documental}
\vspace{-1.0em}

Uno de los aspectos fundamentales en cualquier tipo de investigación lingüística contemporánea es la recolección y análisis de datos. El uso de datos empíricos y observables, independientemente de la fuente de la que provengan, permite a los investigadores validar hipótesis, contrastar teorías y generar nuevo conocimiento sobre el lenguaje y las lenguas (\cite{moyano2025antecedentes,moyano2025corpusexp}). Todos estos procedimientos y métodos para la obtención de datos pueden situarse a lo largo de un \textit{continuum} (\cite{gilquin2009corpora,krug2013research}) dependiendo del carácter ``más o menos natural y más o menos monitorizado del proceso'' (\cite{rojo2021introduccion}), partiendo de la base de que el corpus es el más natural y menos monitorizado, llegando hasta la experimentación como fuente de datos menos natural pero más monitorizada:

\begin{figure}[H]
    \centering
    \includegraphics[width=0.8\linewidth]{images/recogida_datos_guilquin_gries_2009.png}
    \caption{Métodos para la recogida de datos en lingüística. (Fuente: \cite{krug2013research} a partir de \cite{gilquin2009corpora}, tomado de \cite{rojo2021introduccion})}.
    \label{fig:recogida-datos}
\end{figure}

Considerando que el análisis de corpus es una práctica realizada por investigadores de distintas disciplinas y subdisciplinas, el término \emph{corpus} ha recibido diferentes conceptualizaciones y definiciones (§ lo discutido en \ref{sec:ling_corpus}). De manera generalizada, los atributos asignados al término corpus suelen ser (\cite{Egbert_Biber_Gray_2022}):

\begin{itemize}
    \renewcommand\labelitemi{$\rightarrow$}
    \item colección/muestra/cuerpo/\textit{dataset}\footnote{Listado adaptado y traducido de (\cite{Egbert_Biber_Gray_2022}, Figure 1.1. Key concepts included in the definition of a “corpus”.) }
    \item textos
    \item representativo
    \item electrónico
    \item amplio
    \item colección representativa de una lengua o dominio
    \item colección que permite investigaciones sobre fenómenos lingüísticos
    \item natural/auténtico
\end{itemize}

Teniendo en cuenta estas apreciaciones, establecemos que un corpus es una colección de textos, orales o escritos, producidos de manera natural y almacenados en formato electrónico (\cite{Sinclair1991, McEnery2006, Gries2009, rojo2021introduccion}), que han sido seleccionados en función de criterios explícitos y con un propósito específico en mente (\cite{Atkins1992}). Un corpus ha de suponer una muestra representativa de una lengua o de una variedad lingüística, con el fin de ser un modelo de la realidad lingüística y del uso que sus hablantes hacen de ella (\cite{Pitkowski2009}).

Desde el establecimiento de una terminología con orientación lingüística, el uso de corpus textuales como fuente documental se ha consolidado como metodología de trabajo (§ Capítulo \ref{chap:3_extraccion}) por dos motivos fundamentales (\cite{duran2011trabajo}): por un lado, por la necesidad de identificar y recopilar aquellos términos que los especialistas usan realmente (\cite{cabre1999terminology}), y, por otro, por la posibilidad de procesar toda esta información especializada de manera automática o semiautomática (\cite{duran2011trabajo}). En el contexto de la extracción terminológica, así como para la modelización teórico-computacional del léxico especializado y la elaboración de material lexicográfico, resulta imprescindible trabajar con un corpus especializado del dominio (§ Figura \ref{fig:objetivos-corpus}). A diferencia de los corpus generales, un corpus especializado reúne textos propios del discurso de especialidad, proporcionando, entre otros aspectos lingüísticos, contextos de uso reales para las unidades léxicas especializadas que conforman el vocabulario técnico de un ámbito de especialización determinado, hecho alineado con las premisas de la terminología contemporánea, que enfatiza el análisis de los términos en su uso comunicativo real (\cite{cabre1999terminology}). De este modo, la compilación de un corpus técnico-documental de lingüística de corpus se ve justificada por la necesidad de disponer de evidencia textual auténtica del dominio, asegurando que la extracción de términos y la elaboración del glosario bilingüe se fundamenten en usos lingüísticos reales, contribuyendo a la calidad y validez del trabajo terminográfico. En este sentido, la bibliografía científica señala que la calidad de los resultados en terminología está directamente ligada a la calidad y exhaustividad de un proceso previo de documentación (\cite{picht1985terminology, rondeau1984,bowker1996towards, meyer1996corpus, vargas2005aproximacion}).

\begin{figure}[H]
\centering
\begin{tikzpicture}[
    node distance=0.3cm and 0.3cm,
    every node/.style={font=\small},
    objective/.style={rectangle, rounded corners=5pt, thick, text width=4.5cm, align=center, minimum height=1.2cm, fill=gray!10},
    core/.style={ellipse, thick, minimum height=1cm, minimum width=1cm, align=center, fill=blue!20}
]

% Objetivos
\node[objective] (extraccion) {Extracción \\ automática de términos};
\node[objective, below left=of extraccion] (modelo) {Modelo teórico-computacional};
\node[objective, below right=of extraccion] (glosario) {Glosario digital bilingüe};

% Núcleo
\node[core, below=3.2cm of extraccion] (corpus) {Corpus técnico-documental};

% Flechas
\draw[flecha] (extraccion) -- (corpus);
\draw[flecha] (modelo) -- (corpus);
\draw[flecha] (glosario) -- (corpus);

\end{tikzpicture}
\caption{Convergencia de los objetivos de investigación en el corpus técnico-documental.}
\label{fig:objetivos-corpus}
\end{figure}

Las cuestiones metodológicas referentes al corpus técnico-documental se estructuran en cinco apartados estructurados. En primer lugar, se expone el diseño tipológico del corpus, en el que se detallan las decisiones relativas a su naturaleza documental y especializada, así como los criterios de delimitación temática, lingüística, temporal y genérica adoptados para su configuración (§ \ref{diseño}). A continuación, se describe el proceso de compilación, centrado en una metodología basada en la búsqueda en bases de datos académicas y motores especializados, con el fin de garantizar que las fuentes seleccionadas cumpliesen con ciertos criterios de calidad científica (i.e., estar indexadas en repertorios, contar con procesos de revisión por pares y/o haber sido publicadas en editoriales o revistas reconocidas (§ \ref{construccion}). Posteriormente, se ofrece una descripción del corpus, incluyendo estadísticas básicas relativas a su volumen, distribución por lenguas y tipos de texto (§ \ref{descripcion}). El cuarto apartado detalla el proceso de procesamiento lingüístico y anotación terminológica, con especial atención a los criterios de selección, segmentación y etiquetado utilizados para identificar las unidades terminológicas relevantes en el corpus (§ \ref{anotacion}). Finalmente, se presenta una evaluación del corpus desde perspectivas cuantitativas y cualitativas, atendiendo a criterios como la representatividad terminológica, la actualidad de los textos, la riqueza contextual de los términos y la calidad formal de los datos recogidos (§ \ref{evaluacion}).

\subsection{Diseño} \label{diseño}

El diseño del corpus parte de la premisa metodológica de que todo análisis terminológico \emph{in vivo} exige una fase documental previa de localización, selección y vaciado de fuentes especializadas. Esta etapa resulta fundamental en cualquier trabajo terminológico basado en el corpus como base empírica, cuya finalidad es representar el uso lingüístico de un dominio técnico concreto. Ahora bien, la calidad de estas fuentes debe ser cuidadosamente evaluada, atendiendo específicamente a varios criterios:

\begin{itemize}

\item Deben ser \textit{representativas} del ámbito temático, en consonancia con los objetivos de la tarea y con la delimitación del área temática, de modo que permitan elaborar una lista inicial de unidades lo suficientemente amplia como para cubrir el repertorio terminológico completo del dominio pero lo suficientemente acotada para evitar la inclusión de términos ajenos, colindantes o periféricos.

\item Deben estar \textit{actualizadas}, tanto en lo que respecta a las designaciones efectivamente utilizadas por los expertos como al contenido temático tratado, ya que ``ad-hoc \emph{corpora have been shown to support more responsive and up-to-date terminological descriptions}'' (\cite{laursen2014using}), en tanto que se aspira a poder ofrecer una representación fidedigna del estado actual del dominio.

\item Deben ser lo suficientemente \textit{explícitas} como para garantizar la recuperación e identificación precisa de la información en todas las fases del proceso de trabajo terminológico (\cite{cabre1999terminology}). Esto implica que las fuentes seleccionadas incluyan datos bibliográficos completos y normalizados (autoría, título, fecha de publicación, editorial o revista, identificadores como ISBN/ISSN o DOI, si los hubiere), así como metadatos que permitan su trazabilidad y recuperación posterior, facilitando no solo el control interno durante el proceso de recopilación, sino también la transparencia y reproducibilidad de los resultados en una fase posterior de difusión de los resultados.

\item Deben encontrarse preferiblemente en \textit{formato electrónico}, de manera que sean directamente interpretables mediante métodos informáticos y computacionales (\cite{negrillo2010corpus}), salvo en aquellos casos en los que, por criterios de diseño del corpus, resulte necesario incorporar fuentes que no estén disponibles en formato digital. En este caso, estas se someterán a un proceso de digitalización mediante escáner cenital con el fin de garantizar su posterior tratamiento automático y su integración homogénea en el corpus.

\end{itemize}

Las cuestiones principales a tener en cuenta a la hora de diseñar y construir cualquier tipo de corpus son tamaño, equilibrio y representatividad. Estas propiedades están directamente relacionadas con la/s pregunta/s de investigación y el fenómeno lingüístico que se va a investigar (\cite{randi2010building,lavid2005lenguaje}), en nuestro caso:

\begin{table}[H]
\centering
\fontsize{10}{12}\selectfont
\begin{tabular}{@{}p{4.5cm}p{9cm}@{}}
\toprule
\textbf {Fenómeno lingüístico} & \textbf {Pregunta de investigación} \\
\midrule
{Unidades terminológicas} & ¿Cuáles y cómo son los términos que componen el léxico especializado de la lingüística de corpus? \\
\bottomrule
\end{tabular}
\caption{Definición del objeto de estudio y pregunta de investigación.}
\label{tab:fenomeno-pregunta}
\end{table}

Como parte del diseño metodológico del corpus, se han definido una serie de características  tipológicas que orientan su configuración. El corpus debe ser de tipo documental, con una función claramente orientada a un análisis terminológico descriptivo. Debido a su naturaleza restringida al dominio de la lingüística de corpus, se concibe como un corpus especializado y técnico, pero también metalingüístico, pues pretende dar cuenta del uso del léxico por parte de los especialistas del ámbito de especialización de la lingüística de corpus. Además, se plantea como un corpus bilingüe (inglés-español) que permita el contraste de usos terminológicos en ambas lenguas, bajo una relación comparable (i.e., sin que exista una correspondencia directa entre textos traducidos, sino con textos independientes pero equivalentes en cuanto a temática y tipología textual). Desde el punto de vista temporal, el corpus se define como sincrónico, ya que se limita a una franja temporal concreta y reciente con el fin de observar el uso actual de los términos. En términos de modalidad, se trata de un corpus de textos escritos, en el que se recopilan exclusivamente textos producidos en el ámbito académico. Las fuentes principales del corpus serán textos publicados en acceso abierto y en soporte digital. Esta decisión metodológica implica trabajar con documentos que no solo cumplen criterios de calidad científica, sino que además se adecúan a principios de accesibilidad de los datos y transparencia de los resultados, fundamentales tanto para la replicabilidad de los análisis como para su eventual reutilización por parte de otros investigadores. Finalmente, se prevé que el corpus tenga un tamaño mediano, lo suficientemente amplio como para permitir análisis cuantitativos representativos, pero manejable para tareas de anotación manual y su ulterior evaluación cualitativa.

\begin{table}[H]
\centering
\fontsize{12}{12}\selectfont
\begin{tabular}{@{}p{5cm}p{8cm}@{}}
\toprule
\textbf{Atributo tipológico} & \textbf{Valor} \\
\midrule
Temática & Metalingüístico \\
Función & {Documental, de vaciado y de entrenamiento} \\
Nivel de especialización & {Especializado-técnico} \\
Lenguas & {Bilingüe (inglés-español)} \\
Relación entre lenguas & {Comparable} \\
Temporalidad & {Sincrónico} \\
Modalidad & {Escrito} \\
Género & {Multigénero académico} \\
Origen & {Natural} \\
Formato & {Digital} \\
Acceso & {Abierto} \\
Fuente & {Repositorios bibliográficos} \\
Tamaño & {Mediano} \\
\bottomrule
\end{tabular}
\caption{Diseño preliminar del corpus.}
\label{tab:tipologia-corpus-simple}
\end{table}

\subsubsection{Delimitación conceptual del dominio}\label{subsec:conceptualization}

Establecer los límites conceptuales de un dominio implica tener en cuenta características cualitativas y extralingüísticas que determinan qué textos van a ser incluidos en el corpus, de modo que, previo al proceso de muestreo, es necesario analizar y describir el dominio (\cite{Egbert_Biber_Gray_2022}). Una vez definido este, se consideran tanto los límites externos (i.e., hasta dónde abarca el dominio de análisis y qué relación puede tener con otros dominios cercanos) como internos (i.e., qué tipos de texto predominan en este). Sin embargo, los límites de los ámbitos especializados tienden a ser difusos (\cite{meyer1996corpus}) y la mayoría de áreas de interés de los terminógrafos son interdisciplinares -como es el caso que nos ocupa-, por lo que resulta complejo establecer una delimitación conceptual del dominio desde un punto de vista tanto externo o ``lateral'' (i.e., su relación con otros campos de especialidad fronterizos) como interno o ``superior'' (i.e., cuán generales pueden ser los términos que aparecen en el dominio y su relación con los niveles de especialización que va a incluir el corpus).

En línea con el marco teórico establecido, (§ \ref{sec:ling_corpus}), para la confección del corpus metalingüístico, concebimos la lingüística de corpus en términos estrictos, es decir, según su primera acepción o uso terminológico, como el estudio de las propiedades de los corpus (\cite{Gries2011}). Esta delimitación excluye, por tanto, el uso instrumental del corpus para estudiar fenómenos lingüísticos generales (p. ej., la variación sintáctica o la frecuencia léxica en sí mismas), y se centra exclusivamente en las investigaciones metalingüísticas que giran en torno al corpus como objeto de estudio. Esta decisión conlleva una serie de implicaciones directas en cuanto al nivel de especialización del corpus y a los tipos textuales que deben ser incluidos (§ \ref{subsec:calidad} y \ref{subsec:genres}). Dado que el objetivo es construir un corpus que refleje la terminología y discursos propios de la reflexión metalingüística sobre corpus, los textos seleccionados deben proceder de fuentes en las que el corpus sea el eje central. Esto incluye, por ejemplo, artículos científicos centrados en la elaboración, tipología, etiquetado, explotación o evaluación de corpus lingüísticos; capítulos de manuales especializados sobre diseño y metodologías de corpus, etc. Serán incluidos en el corpus, pues, los textos cuya temática se centre en algunas de estas líneas:

\begin{enumerate}
    \item Historiografía y antecedentes de la lingüística de corpus;
    \item Bases teóricas y fundamentos epistemológicos de la lingüística de corpus;
    \item Principios de muestreo y representatividad;
    \item Diseño y construcción de corpus lingüísticos;
    \item Exposición, documentación y presentación de corpus lingüísticos;
    \item Tipología de corpus (generales, especializados, sincrónicos, diacrónicos, paralelos, comparables, etc.);
    \item Problemas (y soluciones) metodológicos derivados del uso de corpus, y
    \item Herramientas y procesos de anotación, etiquetado y procesamiento de corpus.  

\end{enumerate}

En consecuencia, el nivel de especialización de los textos seleccionados condiciona tanto el grado de tecnicidad de los términos presentes como el perfil discursivo del corpus. Así, se excluyen conscientemente textos divulgativos o de nivel introductorio que no desarrollen los conceptos con profundidad terminológica, así como también se evita incluir textos pertenecientes a dominios colindantes salvo que el foco temático recaiga específicamente en el uso, diseño o evaluación de corpus lingüísticos. Este criterio selectivo es fundamental para asegurar que el corpus resultante mantenga coherencia temática y relevancia terminológica. De este modo, la conceptualización precisa del dominio no solo orienta la selección textual, sino que también sustenta la coherencia terminológico-conceptual del corpus.

\subsubsection{Criterios de calidad y adecuación de las fuentes}\label{subsec:calidad}

La calidad del corpus de trabajo está estrechamente vinculada con los tipos de texto que lo componen, ya que estos determinan en gran medida su adecuación al objeto de estudio, su representatividad tanto en términos lingüísticos como temáticos, así como la fiabilidad de los resultados y las conclusiones que puedan derivar del análisis del corpus (\cite{perez2002explotacion}). En este caso se han seleccionado diversos géneros académicos dentro del dominio (artículos de investigación publicados en revistas especializadas, manuales, capítulos de libro, ponencias o actas de congresos, tesis doctorales, etc.). La inclusión de múltiples subgéneros del discurso académico obedece a la necesidad de garantizar una cobertura contextual amplia de los términos, más que al interés por estudiar variaciones discursivas \textit{per se}, es decir, no se incorporan diferentes tipos de textos con el fin de comparar su retórica, sino para asegurar que cada término de interés aparezca en contextos variados que faciliten su comprensión y análisis. Al abarcar distintos géneros, el corpus proporciona para cada unidad terminológica una gama de contextos en los que se emplea el término —definiciones, ejemplos, discusiones metodológicas, entre otros—, lo cual reforzando la extracción terminológica como la posterior redacción de definiciones en el glosario (ver \ref{chap:5_resultados}). Desde la teoría de corpus, se reconoce que la representatividad de un corpus especializado depende de que las muestras textuales cubran la variación relevante dentro de la “población” de textos objetivo (\cite{McEnery2006}). En este sentido, incluir únicamente un tipo de texto podría dejar fuera parte del léxico especializado. Siguiendo este principio, es necesario, pues, extraer muestras que permitan alcanzar el equilibrio y la representatividad que requiera la investigación. Nuestro objeto de estudio abarca todo un subdominio del discurso científico; por tanto, representar adecuadamente dicho subdominio implica incorporar sus diversos medios de comunicación escritos. Una selección variada pretende garantizar que el corpus no esté sesgado hacia un único tipo de texto.

\subsubsection{Cobertura documental y selección de géneros académicos}\label{subsec:genres}

Para la compilación del corpus se ha adoptado una estrategia de búsqueda y descarga a partir de múltiples fuentes académicas especializadas, en lugar de centralizar la recuperación en un único catálogo. Esta decisión metodológica responde a la necesidad de cubrir de manera lo más exhaustiva posible la producción científica sobre lingüística de corpus en ambas lenguas. Por ello, el diseño del corpus establece de manera explícita qué subgéneros académicos se incluyen y cuáles se excluyen, en función de su relevancia para la representación terminológica de la lingüística de corpus. En este sentido, la representatividad del corpus se vincula tanto a la inclusión de distintos tipos de texto como a la distribución equilibrada de estos, de manera que reflejen fielmente la heterogeneidad de la lingüística de corpus en su dimensión académica. Asimismo, se atiende a la diversidad de distribuciones lingüísticas presentes en dichos textos, con el fin de garantizar que el corpus recoja un espectro amplio de, en el caso que nos ocupa, fenómenos léxicos característicos del dominio.

Los subgéneros incluidos en el corpus son artículos científicos procedentes de revistas indexadas como especializadas en lingüística de corpus (p. ej., \textit{International Journal of Corpus Linguistics}, \textit{Corpora} o \textit{Corpus Linguistics and Linguistic Theory} para el inglés, y \textit{Procesamiento del Lenguaje Natural} (SEPLN), \textit{Revista de Lingüística Teórica y Aplicada} u \textit{Onomázein} para el español); libros y manuales de referencia y capítulos de libro extraídos de monografías especializadas. El diseño del corpus tiene en cuenta cuestiones de acceso abierto y de replicabilidad, en consonancia con las buenas prácticas de la investigación reproducible. La gran parte de los textos proceden de publicaciones disponibles en acceso abierto o acceso institucional disponibles en formato digital, lo que permite compartir el diseño del corpus corpus, las referencias y el procedimiento seguido para su compilación.

\subsubsection{Comparabilidad interlingüística del corpus bilingüe}
Dado que uno de los objetivos generales de esta investigación es la elaboración de un glosario bilingüe inglés-español, el corpus se ha diseñado como bilingüe y comparable, de acuerdo con la terminología comúnmente aceptada (McEnery y Hardie, 2012). La consecución de este objetivo implica que el corpus de trabajo se estructure en dos subcorpus monolingües –uno en inglés y otro en español– integrados por textos similares en cuanto a tema, género y nivel de especialización, pero que no sean traducciones entre sí. La decisión de utilizar un corpus comparable obedece a la posibilidad de  examinar la terminología en cada lengua en su forma auténtica de uso, evitando las posibles distorsiones lingüísticas introducidas por la traducción (\cite{negrillo2010corpus,duran2011trabajo,hourani2017unidades}). Los corpus comparables se han utilizado con éxito en diversos estudios terminológicos recientes, especialmente en el ámbito de la traducción:

\begin{quote}
    [Los corpus paralelos] constituyen una herramienta interesante para solucionar problemas de diferente naturaleza (temáticos, \textit{terminológicos}\footnote{La cursiva es mía.}, textuales, estilísticos, etc.) que surgen durante la realización de una traducción determinada, ya que permiten compilar una documentación fiable y específica de manera económica (en tiempo y coste) y muy eficaz. (\cite{sancheztrigo2005})
\end{quote}

A diferencia de un corpus paralelo (formado por textos traducidos), donde los términos en la lengua meta podrían verse influidos por la estructura de la lengua origen, en un corpus comparable cada término se extrae de textos originalmente escritos en esa lengua, reflejando así la idiosincrasia terminológica propia de cada comunidad científica. En este sentido, la noción de corpus comparable ha sido definida por Baker (\citeyear{baker1995corpora}) y otros autores como aquella colección de textos originales en dos o más idiomas con características similares en cuanto a contenido y condiciones de producción. En nuestro caso, ambos subcorpus deben reunir, idealmente, textos del ámbito de la lingüística de corpus, escritos por y para especialistas, con criterios equivalentes de selección (géneros académicos, periodo temporal, nivel técnico, etc.). Este paralelismo en el diseño garantiza la comparabilidad: podemos contrastar las dos lenguas asumiendo que las diferencias encontradas responden a genuinas diferencias terminológicas o discursivas entre el inglés y el español, y no a desequilibrios en el muestreo de textos. Ahora bien, aunque en principio se aspira a una simetría en cuanto a la selección y el muestreo de textos, podemos anticipar que la disponibilidad efectiva de materiales en acceso abierto no es igual para ambas lenguas; en particular, la producción científica en inglés sobre lingüística de corpus es notablemente más abundante, derivando bien en un subcorpus en inglés de mayor tamaño o bien en un submuestreo de dicha producción con el fin de mantener la comparabilidad entre ambos conjuntos.


\subsubsection{Delimitación temporal y vigencia terminológica}
Otro aspecto fundamental del diseño es la delimitación sincrónica del corpus, basada en los cambios recientes de la disciplina. En lugar de compilar textos de un amplio rango temporal o incluir documentos históricos de la lingüística de corpus, se ha optado por circunscribir el corpus a un período contemporáneo específico, aproximadamente la última década. Esta decisión responde a que la lingüística de corpus ha experimentado desarrollos significativos en años recientes –por ejemplo, la incorporación de grandes volúmenes de datos, nuevas técnicas de PLN, etc.– que han podido derivar en la inclusión de nuevos términos. Con el fin de reflejar el estado actual del área de especialidad, resulta necesario que los textos recopilados respondan a un criterio de actualidad, es decir, que el corpus sea contemporáneo. Siguiendo esta premisa, el corpus se limita sincrónicamente a publicaciones del año 2015 en adelante, periodo en el cual la disciplina ha incorporado nuevos términos (p. ej., \textit{big data}, \textit{embeddings}, \textit{data-driven approaches}, etc.) que no aparecerían en textos más antiguos. Es más, de acuerdo con Kyto et al. (\citeyear{kyto2000building}) y Durán Muñoz (\citeyear{duran2011trabajo}), una horquilla temporal demasiado amplia podría comprometer la coherencia terminología, al incluir muestras potencialmente desactualizadas que podrían no reflejar el uso actual del léxico de especialidad (p. ej., términos que pudieron tener acepciones diferentes o menor frecuencia décadas atrás podrían confundir el proceso de extracción y compilación de un recurso lexicográfico actual). Un corpus sincrónico garantiza capturar un corte temporal representativo del lenguaje especializado tal y como lo emplean los especialistas actualmente. Según la definición clásica, un corpus sincrónico reúne textos publicados en un periodo determinado de tiempo y, en nuestro diseño, esto se traduce en incorporar producciones científicas de los últimos años relevantes para la investigación. Esta acotación temporal está asimismo en línea con la noción de corpus con finalidad terminológica en contextos dinámicos: únicamente a partir de la observación de textos recientes se pueden detectar y documentar los cambios terminológicos de una disciplina en constante evolución, y únicamente a partir de dicha base empírica es posible formular generalizaciones fiables acerca de las tendencias de variación, fijación y posible obsolescencia terminológica en el dominio.

\subsubsection{Tamaño, dimensiones y viabilidad}
La construcción de un corpus adecuado es fundamental en estudios de terminología basada en corpus, puesto que los textos recopilados del dominio constituyen la base documental indispensable para la descripción del lenguaje especializado y la elaboración de recursos especializados, proporcionando contextos y datos empíricos fiables (\cite{cabre1999terminology,picht1985terminology}). Por ello, el tamaño del corpus resulta de vital importancia: debe ser lo suficientemente amplio para garantizar la representatividad terminológica, pero a la vez, en el caso que nos ocupa, manejable para permitir una anotación manual exhaustiva (§ \ref{anotacion}). No existe un tamaño específico ni estandarizado para la creación de corpus especializados —ni tampoco para los generales—; sin embargo, el tamaño continúa siendo una de las principales preocupaciones de los investigadores en lingüística de corpus (\cite{Egbert_Biber_Gray_2022}). En este sentido, tanto el tamaño como el resto de las propiedades del corpus dependen de su configuración y de los objetivos establecidos para su exploración (\cite{rojo2021introduccion}), tal y como se resume en la Tabla \ref{tab:fenomeno-pregunta}.
  

Ahora bien, el tamaño óptimo no es un valor absoluto, sino una variable dependiente de múltiples factores: el tipo de tarea terminológica, la lengua o lenguas implicadas, el enfoque utilizado (estadístico, lingüístico o híbrido), y la naturaleza del dominio. En el caso de la extracción automática bilingüe, por ejemplo, se ha observado que corpus más pequeños pero cuidadosamente seleccionados pueden producir resultados de mayor calidad que grandes colecciones de textos ruidosos o poco representativos (Morin y Daille, 2006; Morin, Daille y Kageura, 2010). Asimismo, en contextos jurídicos, se ha propuesto el uso de herramientas específicas para calcular tamaños mínimos funcionales en corpus \textit{ad hoc} (\cite{seghiri2014toobig}). En tareas de clasificación de términos mediante modelos distribucionales, el rendimiento depende tanto del tamaño como de parámetros como la ventana de contexto, la granularidad sintáctica o la estrategia de desambiguación (Zadeh, 2016). Además, en vistas a su procesamiento computacional, un corpus mediano permite aplicar eficientemente técnicas semiautomáticas como los cálculos de frecuencia, dispersión o asociación, y facilita una posterior filtración y anotación manual. El tamaño también impacta en la posibilidad de realizar doble anotación o revisión cruzada, prácticas metodológicamente deseables pero inviable en corpus excesivamente extensos. Como señalan Kit (2002) o Wermter y Hahn \citeyear{wermter2005finding}, los corpus de gran escala son útiles para detectar neología mediante variación distribucional, pero su uso requiere herramientas especializadas y criterios de validación exigentes, lo que no siempre resulta compatible con proyectos que buscan un control humano más directo.

En consonancia con estas apreciaciones, para este estudio se ha optado inicialmente por un corpus de tamaño mediano, que garantice la cobertura suficiente del dominio de la lingüística de corpus y que a su vez resulte asumible desde el punto de vista del costo computacional, la exploración manual y la anotación lingüística, una decisión que no solo viene dada por cuestiones metodológicas, sino que también asegura la calidad del glosario bilingüe resultante, permitiendo revisar contextos, extraer concordancias, y evaluar cada término candidato de forma contextualizada.


\subsection{Procedimiento de compilación y selección de fuentes del corpus} \label{construccion}

La compilación del corpus terminológico bilingüe se ha llevado a cabo, tal y como se ha establecido en una fase previa de diseño, mediante un proceso manual y controlado, a fin de asegurar la relevancia y calidad de los textos recopilados. El corpus resultante está compuesto exclusivamente por publicaciones científico-académicas, concretamente artículos de investigación, capítulos de libros y libros (monografías especializadas y manuales). Estas obras se obtuvieron a través de las búsquedas controladas mencionadas, abarcando tanto revistas científicas especializadas en el ámbito de la lingüística de corpus como colecciones de monografías publicadas por editoriales de reconocido prestigio. Para garantizar el acceso a textos representativos, se incluyeron fuentes de diversas procedencias: materiales en acceso abierto (\textit{open access}) y otros disponibles mediante suscripción institucional, aprovechando los recursos bibliotecarios y bases de datos académicas disponibles.

\begin{figure}[H]
\centering
\begin{adjustbox}{max width=\textwidth}
\begin{tikzpicture}[
  font=\small,
  node distance=1.2cm and 2.6cm,
  startstop/.style={ellipse, draw, thick, fill=green!15, text width=5cm, align=center, inner sep=6pt, minimum height=1.1cm},
  source/.style={rectangle, rounded corners=3pt, thick, fill=blue!15, text width=4.0cm, align=center, inner sep=6pt},
  process/.style={rectangle, rounded corners=3pt, thick, fill=gray!15, text width=4.6cm, align=center, inner sep=6pt},
  decision/.style={diamond, aspect=2, draw, thick, fill=yellow!25, text width=4.3cm, align=center, inner sep=2pt},
  ovalnode/.style={ellipse, draw, thick, fill=purple!12, text width=5.2cm, align=center, inner sep=6pt, minimum height=1.1cm},
  % --- Estilos para las dos extracciones:
  extractpre/.style={
    trapezium, trapezium stretches=true, trapezium left angle=70, trapezium right angle=110,
    draw, thick, fill=teal!18,
    text width=5.2cm, align=center, inner sep=6pt
  },
  % Hexágono de tamaño fijo (el texto no cambia el tamaño)
  extractml/.style={
    regular polygon, regular polygon sides=6,
    shape border uses incircle,
    draw, thick, fill=orange!22,
    minimum size=4.5cm,  % ajusta 5.0--5.4cm si quieres fino
    inner sep=0pt
  },
  flecha/.style={-{Stealth}, thick, shorten >=1pt}
]

% --- Inicio: criterios de diseño
\node[startstop] (criteria) {Búsqueda y selección de fuentes};

% --- Fila superior: TRES PATITAS (fuentes)
\node[source, below left=1.6cm and 5cm of criteria] (open) {Acceso abierto};
\node[source, below=1.6cm of criteria] (inst) {Acceso institucional};
\node[source, below right=1.6cm and 5cm of criteria] (scan) {Digitalización};

% Flechas desde criterios hacia cada fuente
\draw[flecha] (criteria.south west) -- (open.north);
\draw[flecha] (criteria.south) -- (inst.north);
\draw[flecha] (criteria.south east) -- (scan.north);

% --- Fusión hacia PDFs
\node[process, below=2.2cm of inst] (pdfs) {Recopilación de PDFs};
\draw[flecha] (open.south) |- (pdfs.west);
\draw[flecha] (inst.south) -- (pdfs.north);
\draw[flecha] (scan.south) |- (pdfs.east);

% --- Óvalo: Corpus piloto
\node[ovalnode, below=of pdfs] (pilot) {Balance inicial 50--50\\ \textbf{Corpus piloto (100 textos)}};
\draw[flecha] (pdfs.south) -- (pilot.north);

% --- Conversión y estadísticas en paralelo
\node[process, below left=1.2cm and 1.8cm of pilot] (conv) {Conversión PDF $\rightarrow$ TXT\\ (\emph{ABBYY FineReader})};
\node[process, below right=1.2cm and 1.8cm of pilot] (stats) {Estadísticas iniciales del corpus};
\draw[flecha] (pilot.south) -| (conv.north);
\draw[flecha] (pilot.south) -| (stats.north);

% --- Decisión submuestreo
\node[decision, below=1.6cm of pilot] (sub) {¿Submuestreo?};
\draw[flecha] (conv.south) |- (sub.west);
\draw[flecha] (stats.south) |- (sub.east);

% === RAMA NO: Extracción preliminar (SIN submuestreo) ===
\node[extractpre, below left=2.0cm and 5.0cm of sub] (extract) {Extracción preliminar de candidatos};
\draw[flecha] (sub.south west) -- node[above, sloped]{No} (extract.north);

\node[process, below=of extract] (rank) {Filtrado y selección de unidades terminológicas};
\draw[flecha] (extract.south) -- (rank.north);

\node[process, below=of rank] (eval_extr) {Evaluación cualitativa \\ \footnotesize (pertinencia terminológica, coherencia, análisis de errores)};
\draw[flecha] (rank.south) -- (eval_extr.north);

% === RAMA SÍ: Extracción basada en aprendizaje automático (CON submuestreo) ===
\node[ovalnode, below right=2.0cm and 5.0cm of sub] (balanced) {\textbf{Corpus balanceado en tokens}\\ \footnotesize (sólo para aprendizaje automático)};
\draw[flecha] (sub.south east) -- node[above, sloped]{Sí} (balanced.north);

\node[process, below=of balanced] (ann) {Anotación terminológica};
\draw[flecha] (balanced.south) -- (ann.north);

% --- Flecha transversal de pre-anotación
\draw[flecha] (rank.east) to[out=0, in=180, looseness=1.1]
  node[above, sloped]{\footnotesize Pre-anotación} (ann.west);

% --- Hexágono (forma fija) + texto superpuesto
\node[extractml, below=1.5cm of ann] (ml_extract) {};
\node[align=center, text width=5.0cm] at (ml_extract.center) {Extracción basada en\\ aprendizaje automático};

\draw[flecha] (ann.south) -- (ml_extract.north);

\node[process, below=1.5cm of ml_extract] (eval_ml) {Evaluación cuantitativa \\ \footnotesize (F1 macro/por clase, precisión y recuperación)};
\draw[flecha] (ml_extract.south) -- (eval_ml.north);

\end{tikzpicture}
\end{adjustbox}
\caption{Flujo de construcción del corpus desde su diseño hasta la anotación final}
\label{fig:flujo-diseno-sub-geom}
\end{figure}

La selección de textos obedece a criterios estrictos previamente definidos en el diseño tipológico del corpus, cubriendo dimensiones temáticas, genéricas, lingüísticas y temporales. Por el criterio temático, se han incluido únicamente trabajos claramente enmarcados en el dominio de la lingüística de corpus y pertinentes para el estudio terminológico bilingüe, asegurando así la relevancia conceptual de los términos recopilados (§ \ref{subsec:conceptualization}). En cuanto al criterio genérico, el corpus se circunscribe a subgéneros académicos específicos: artículos de investigación, capítulos de libro y libros, manuales o monografías especializadas, excluyendo otros tipos de textos no científicos o divulgativos para mantener la homogeneidad formal. El criterio lingüístico implicó incorporar textos en las dos lenguas de trabajo (español e inglés, en este caso). Por último, se ha reducido el corpus a un criterio temporal que delimita el período de publicación de los textos, logrando una representación actualizada y equilibrada del dominio. Por otra parte, puesto que algunas obras que consideramos clave solo se encontraban disponibles en formato físico (particularmente ciertos libros y manuales fundamentales), se ha procedido a a su digitalización mediante un escáner cenital profesional (modelo \textit{Zeutschel OS I2002 Advanced Plus})\footnote{<https://ila.uca.es/digitalizacion-de-alta-resolucion/>} para incorporarlos al corpus en formato electrónico.

\subsection{Descripción cuantitativa y distribución interna del corpus} \label{descripcion}
Tras el proceso de recopilación de fuentes atendiendo a los criterios descritos en la fase de diseño, el corpus resultante consta de 100 fuentes distribuidas de manera equitativa en las lenguas de trabajo. Ahora bien, aunque en términos numéricos el corpus presenta un equilibrio aparente (50–50), resulta pertinente analizar con mayor detalle los valores descriptivos por lengua, con el fin de comprobar hasta qué punto la distribución es verdaderamente homogénea en aspectos como extensión, número de palabras o tipología textual.

La Tabla \ref{tab:en_corpus} muestra, por un lado, la distribución del corpus en inglés según el tipo de texto. En este caso, los libros y manuales, aunque solo suman siete documentos, concentran con diferencia la mayor parte de los \textit{tokens} (más del 70\%). Este peso cuantitativo resulta esperable, ya que se trata de textos de gran extensión media ($\approx 132,200$ \textit{tokens}), con una desviación típica elevada ($\approx 30,200$), lo que refleja la presencia tanto de manuales más sintéticos como de monografías más extensas. Los capítulos de libro representan un número reducido de documentos (cinco), pero con un tamaño relativamente homogéneo ($\approx 11,500$ \textit{tokens} de media y una desviación típica moderada de $\approx 4,500$). Finalmente, los artículos de revista constituyen el grupo mayoritario en términos de número de documentos (38), pero con un peso proporcional en \textit{tokens} notablemente menor ($\approx 25$\%). Su extensión media ronda los $8,800$ \textit{tokens} y se caracteriza por una desviación típica también en torno a los $4,000$, lo que confirma que este subgénero es más uniforme y homogéneo en su longitud. En conjunto, la configuración del subcorpus en inglés reproduce el patrón habitual de la producción académica: pocos libros con gran impacto cuantitativo, numerosos artículos más breves y algunos capítulos intermedios, garantizando así un equilibrio representativo entre cantidad de documentos y volumen textual.


\begin{table}[H]
\centering
\fontsize{10}{12}\selectfont
\begin{tabular}{lrrrrr}
\toprule
Tipo de texto & Nº doc. & Tokens & Media & Desv. típ. & \% tokens \\
\midrule
Capítulos de libro & 5&  57\,439& 11\,487,8&  4\,518,7&  4,3\%\\
Libros/manuales    & 7  & 925\,591  & 132\,227,3 & 30\,166,0 & 70,3\% \\
Artículos de revista & 38& 333\,860&  8\,785,8&  4\,037,5& 25,4\%\\
\rowcolor{gray!20}\textbf{Total} & \textbf{50} & 1\,316\,890& 26\,337,8& — & \textbf{100,0\%} \\
\bottomrule
\end{tabular}
\caption{Distribución y características del corpus en inglés por tipo de texto.}
\label{tab:en_corpus}
\end{table}

Por otro lado, la Tabla \ref{tab:es_corpus} muestra los valores cuantitativos del corpus en español. En ella se aprecia que los libros/manuales, pese a ser únicamente cuatro, aportan casi el 40\% del total de tokens, algo que no debe entenderse directamente como un sesgo no deseado, sino como un rasgo inherente a su naturaleza textual (i.e., un libro, de manera generalizada, contiene una extensión muy superior a la de un artículo o un capítulo, de modo que su peso cuantitativo en \textit{tokens} es esperable). La desviación típica elevada ($\approx91 600$ \textit{tokens}) refleja, además, la heterogeneidad habitual entre obras de distinta envergadura (desde manuales breves hasta monografías de gran extensión). Los capítulos de libro, por su parte, constituyen la mayor parte de los documentos incluidos en el corpus (hasta 28), y también concentran una proporción importante de \textit{tokens} ($\approx40$\%). La desviación típica igualmente alta ($\approx17 100$) confirma que se trata de un subgénero diverso, que contiene desde capítulos cortos hasta otros mucho más extensos. En contraste, los artículos de revista representan una quinta parte de los \textit{tokens} ($\approx20$\%) y se caracterizan por una mayor homogeneidad en su extensión, con una media en torno a 9 500 tokens y una desviación típica relativamente baja ($\approx3 000$), siendo este el tipo de texto más estable en cuanto a extensión. Por lo tanto, la configuración del corpus en español no debe interpretarse como un desequilibrio, sino como la consecuencia natural de la tipología textual incluida: que los libros ocupen un lugar central en términos de \textit{tokens} es no solo esperable, sino también coherente con el objetivo de reflejar la producción académico-científica real en lingüística de corpus.

\begin{table}[H]
\centering
\fontsize{10}{12}\selectfont
\begin{tabular}{lrrrrr}
\toprule
Tipo de texto & Nº doc. & Tokens & Media & Desv. típ. & \% tokens \\
\midrule
Capítulos de libro & 28 & 333\,653 & 11\,916,2 & 17\,173,0 & 39,8\% \\
Libros/manuales    &  4 & 332\,330 & 83\,082,5 & 91\,654,5 & 39,7\% \\
Artículos de revista & 18 & 171\,655 &  9\,536,4 &  3\,008,5 & 20,5\% \\
\rowcolor{gray!20}\textbf{Total} & \textbf{50} & \textbf{837\,638} & \textbf{16\,752,8} & — & \textbf{100,0\%} \\
\bottomrule
\end{tabular}
\caption{Distribución y características del corpus en español por tipo de texto.}
\label{tab:es_corpus}
\end{table}

En términos cuantitativos, ambos subcorpus presentan una estructura general similar, en tanto que ambos incluyen las tres tipologías textuales (libros/manuales, capítulos de libro y artículos de revista) y mantienen un equilibrio formal en el número de documentos (50 en cada lengua). Sin embargo, al considerar los valores internos y distribuidos, surgen algunas diferencias que merecen ser cuestionadas. En cuanto al volumen total de \textit{tokens}, los dos subcorpus no presentan un equilibrio perfecto: mientras que el bloque en inglés alcanza aproximadamente $1\,316\,900$ \textit{tokens}, el bloque en español suma unos $837\,600$. Esto hace al subcorpus en inglés alrededor de un 57\% más grande. Por tipologías textuales, las diferencias son todavía más acusadas, pues en cuanto a libros y manuales el subcorpus inglés concentra cerca de $925\,600$ \textit{tokens} (70\%), frente a los $332\,300$ del español (40\%). Los capítulos de libro en español suman aproximadamente $333\,700$ \textit{tokens} (40\%), mientras que en inglés apenas llegan a $70\,100$ (5\%). El bloque español, por tanto, se apoya mucho más en este tipo textual, lo que aumenta su heterogeneidad interna. En cuanto a los artículos de revista, los valores son comparables en términos proporcionales (24\% en inglés, 20\% en español), aunque el inglés prácticamente duplica en número absoluto ($321\,200$ frente a $171\,700$). Este desbalance tiene consecuencias directas para la extracción terminológica y, especialmente, para el entrenamiento de modelos de aprendizaje automático. En un análisis exploratorio (extracción lingüística, estadística o híbrida), mantener los volúmenes tal cual es válido, pues permite reflejar la producción textual real en cada lengua; sin embargo, para experimentos supervisados, el mayor número de \textit{tokens} en inglés puede inducir sesgos en el entrenamiento, al ofrecer más contextos de aparición. 

No obstante, este desequilibrio no debe interpretarse necesariamente como una limitación, puesto que tal y como han señalado estudios clásicos (\cite{Zweigenbaum2001}) y recientes en torno al diseño de corpus \textit{ad-hoc} (\cite{Gupta2021, Yang2023}), el balance en este tipo de recursos no tiene por qué ser absoluto, sino funcional: lo relevante es que el corpus sea representativo de la variedad lingüística que se desea modelar y que se documenten de manera transparente las decisiones de muestreo. En este sentido, los corpus generales o referencia requieren un balance amplio en géneros y registros para garantizar inferencias generalizables, mientras que los corpora científicos, técnicos o ad hoc deben priorizar la representatividad del dominio y la densidad de fenómenos relevantes para la tarea (\cite{VandenAuken2014, Nastou2024}). Para métodos de extracción terminológica clásica (i.e., métodos estadísticos o lingüísticos sin aprendizaje automático), la bibliografía especializada señala que no resulta necesario un equilibrio estricto entre tipologías textuales. Lo fundamental, en cambio, es disponer de un corpus lo suficientemente grande y representativo del dominio, ya que el rendimiento de las medidas de frecuencia, especificidad y asociación depende más de la densidad terminológica interna que de la proporcionalidad de los géneros incluidos (\cite{Yamamoto2001,Fuentes2001}). De hecho, incluir textos extensos —como manuales o monografías— puede aportar una mayor estabilidad a las estimaciones de frecuencia y coocurrencia, siempre y cuando se controle la documentación de las fuentes  integradas en el corpus (\cite{Kruger2012, Braun2007}).

\subsection{Anotación terminológica del corpus: fases, recursos y criterios} \label{anotacion}

La anotación de corpus consiste en hacer explícita y accesible información lingüística interpretable mediante su incorporación estructurada al corpus (\cite{lavid2005lenguaje,leech2004,hajicova2010treebank}). Existen tres grandes tipos de anotación: (i) la anotación extratextual, que generalmente atiende a los metadatos relevantes de los textos (p. ej., fecha, autoría, género, variedad lingüística, etc.); (ii) la anotación textual, que estructura los textos en cabeceras, párrafos, oraciones, etc., y, por último, y de nuestro interés en este caso, (iii) la anotación lingüística. Este tipo de anotación incluye información sobre rasgos lingüísticos específicos que se encuentran implícitos en los textos (\cite{lavid2005lenguaje}), y que el lingüista ha de explicitar mediante la anotación manual o automática.

Independientemente de su naturaleza o nivel de intervención, toda anotación debe cumplir con ciertos principios metodológicos fundamentales. En primer lugar, la capa de anotación debe ser separable y opcional, es decir, que esta no debe alterar irreversiblemente el texto base y debe poder añadirse o retirarse sin comprometer la integridad del corpus. En segundo lugar, ha de ser detallada y adecuadamente documentada, de modo que otros investigadores puedan interpretar su contenido, reproducirla o adaptarla a otros fines; de ahí que se recurra al desarrollo previo de un esquema o unas directrices de anotación, especialmente en casos donde se recurre a más de un anotador. Por último, aunque no contemplado en este trabajo, en aquellos casos en que intervienen varios anotadores, es esencial que la anotación sea consensuada, con criterios previamente acordados que garanticen la coherencia y la fiabilidad de los datos (\cite{leech2004}), generalmente recurriendo a métricas de acuerdo interanotador, como el índice Kappa de Cohen, el coeficiente de correlación de Pearson o la medida de concordancia porcentual, entre otras (\cite{artstein2017inter}).

Una vez cumplimentada una primera fase de anotación extratextual, en la que se registraron metadatos relevantes como la autoría, la fecha de publicación, la lengua o el tipo de texto (§ \hyperref[appendix:registro]{Apéndice~I} para el registro completo), se procedió a la anotación lingüística, que se ha dividido en dos fases consecutivas: una automática, basada en diversas fuentes de autoridad, y otra manual, llevada a cabo mediante la lectura exhaustiva de los textos que conforman el corpus. La anotación automática se ha realizado a partir de un glosario de 39 términos clave o términos ``semilla'' (del inglés, {\textit{seed terms}}) (\cite{heylen2015automatic}) de lingüística de corpus. El procedimiento técnico consiste en la elaboración de dos listas —una por cada lengua del corpus— compuestas por expresiones especializadas que describen nociones, metodologías y herramientas propias del ámbito de especialización de la lingüística de corpus. Estas listas han sido extraídas directamente del glosario incluido en \emph{CASS: Briefings – Corpus Linguistics: Some Key Terms} (2013)\footnote{El \href{https://cass.lancs.ac.uk/wp-content/uploads/2013/12/CASS-Gloss-final1.pdf}
{\textit{CASS: Briefings – Corpus Linguistics: Some Key Terms}} es un documento elaborado por la Universidad de Lancaster que reúne una selección representativa de términos clave en lingüística de corpus y cuya finalidad es facilitar la comprensión de conceptos fundamentales en el área tanto a investigadores como a estudiantes.}, de la Universidad de Lancaster, que reúne una selección de términos clave en lingüística de corpus. Entre las designaciones seleccionadas se encuentran, por ejemplo, \textit{corpus linguistics}, \textit{representative}, \textit{monitor corpus} o \textit{corpus-driven}. 
La Tabla \ref{tab:glosario_cass} recoge el listado completo de términos semilla en ambas lenguas. Adicionalmente, con el objetivo de elaborar un listado equivalente en español que permitiera la construcción de un corpus comparable bilingüe (inglés-español), se procedió a la búsqueda sistemática de los términos originales del glosario en varios bancos de datos terminológicos reconocidos. Esta consulta permitió la identificación de equivalentes normalizados o ampliamente aceptados en el ámbito hispanoparlante para una parte significativa de los términos. No obstante, en aquellos casos en los que no se encontró una forma documentada, se propuso una traducción funcional basada en el uso frecuente en publicaciones académicas especializadas. Conviene subrayar que este listado preliminar no pretende constituir un inventario exhaustivo del vocabulario especializado del dominio, ni funcionar como un estándar de referencia cerrado, sino que cumple una función estrictamente heurística, cuyo objetivo es facilitar una primera identificación automática de contextos potencialmente especializados que sirva de apoyo a la fase posterior de anotación manual. La Tabla \ref{tab:glosario_cass} recoge el listado de términos semilla en ambas lenguas. El asterisco marca aquellas unidades que cuentan con equivalente en español y que también aparece con su forma en inglés.

\begin{table}[H]
\centering
\fontsize{10}{12}\selectfont
\begin{tabular}{lll}
\toprule
\textbf{Término (\texttt{en})} & \textbf{Equivalente (\texttt{es})} & \textbf{Fuente del equivalente}\\
\midrule
\textit{Annotation} & \textit{Anotación} & Gestor ILA \\
\textit{Balance} & \textit{Equilibrio} & -\\
\textit{Colligation} & \textit{Coligación} & - \\
\textit{Collocation} & \textit{Colocación} & CercaTerm \\
\textit{Comparability} & \textit{Comparabilidad} & Gestor ILA\\
\textit{Concordance} & \textit{Concordancia} & Gestor ILA\\
\textit{Concordancer} & \textit{Programa de concordancias} & TERMIUM Plus\\
\textit{Corpus} & \textit{Corpus (textual)} & TERMIUM Plus\\
\textit{Corpus construction} & \textit{Construcción de corpus} & Gestor ILA \\
\textit{Corpus-based} & \textit{Basado en corpus} & Gestor ILA \\
\textit{Corpus-driven} & \textit{Guiado por el corpus}* & Gestor ILA \\
\textit{Diachronic} & \textit{Diacrónico} & CercaTerm  \\
\textit{Encoding} & \textit{Codificación} & CercaTerm \\
\textit{Frequency list} & \textit{Lista de frecuencias} & Gestor ILA \\
\textit{Key word in context (KWIC)} & \textit{Palabra clave en contexto}* & Gestor ILA\\
\textit{Keyword} & \textit{Palabra clave} & CercaTerm \\
\textit{Lemma} & \textit{Lema} & CercaTerm \\
\textit{Lemmatisation} & \textit{Lematización} & CercaTerm \\
\textit{Lexis} & \textit{Léxico} & - \\
\textit{Markup} & \textit{Marcado} & CercaTerm \\
\textit{Markup language} & \textit{Lenguaje de marcado} & TERMIUM Plus \\
\textit{Metadata} & \textit{Metadato/s} & CercaTerm \\
\textit{Monitor corpus} & \textit{Corpus monitor} & Gestor ILA \\
\textit{Node} & \textit{Nodo} & Gestor ILA \\
\textit{Reference corpus} & \textit{Corpus de referencia} & Gestor ILA \\
\textit{Representativeness} & \textit{Representatividad} & -\\
\textit{Sample} & \textit{Muestra} & CercaTerm \\
\textit{Sample corpus} & \textit{Corpus de muestra/muestral} (G. Rojo?) & - \\
\textit{Sampling frame} & \textit{Marco de muestreo} & -\\
\textit{Significance test} & \textit{Prueba estadística} & TERMIUM Plus \\
\textit{Statistical significance} & \textit{Significación estadística} & TERMIUM Plus\\
\textit{Synchronic} & \textit{Sincrónico} & CercaTerm \\
\textit{Tagging} & \textit{Etiquetado} & (uso=ing. c,\\pero no viene en español) \\
\textit{Text} & \textit{Texto} & CercaTerm \\
\textit{Token} & \textit{Token} & Gestor ILA \\
\textit{Type} & \textit{Tipo} & Gestor ILA \\
\textit{Type-token ratio} & \textit{Type-token ratio} & Gestor ILA\\
\textit{XML} & \textit{XML} & TERMIUM Plus, CercaTerm \\
\bottomrule
\end{tabular}
\caption{Listado de términos ``semilla''.}
\label{tab:glosario_cass}
\end{table}

La Tabla \ref{tab:proporcion_equivalencias} resume la proporción de términos cuya equivalencia cuenta con respaldo terminológico frente a aquellos que no lo tienen. Este problema de documentación respecto a los equivalentes terminológicos anticipa varias implicaciones metodológicas. La necesidad de construir un corpus bilingüe comparable pone de manifiesto cierta desigualidad con respecto a la disponibilidad de términos documentados entre ambas lenguas, ya que poco más de la mitad de las equivalencias cuentan con respaldo en bancos de datos terminológicos consolidados. Este hecho evidencia una asimetría en los recursos terminológicos existentes, incluso en una disciplina con alto grado de internacionalización como la lingüística de corpus. En consecuencia, ha sido necesario recurrir a traducciones funcionales ampliamente utilizadas en la literatura académica hispanohablante, o bien mantener ciertos términos en su forma original, especialmente cuando su uso está lexicalizado en el discurso especializado (por ejemplo, \textit{token} o\textit{ type-token ratio}). Asimismo, la falta de documentación oficial para parte del vocabulario técnico puede influir en la calidad y la cobertura de la descarga de documentos técnicos y la primera anotación automática basada en coincidencias léxicas, generando posibles sesgos entre las lenguas.

\begin{table}[H]
\centering
\fontsize{10}{12}\selectfont
\begin{tabular}{lrr}
\toprule
\textbf{Categoría} & \textbf{Cantidad} & \textbf{Proporción} \\
\midrule
Total de términos & 39 & 100\,\% \\
\quad Equivalentes con fuente documentada & 20 & 51.3\,\% \\
\quad Equivalentes sin fuente documentada & 19 & 48.7\,\% \\
\bottomrule
\end{tabular}
\caption{Proporción de equivalencias en español y documentación terminológica.}
\label{tab:proporcion_equivalencias}
\end{table}

Mediante un proceso de \textit{matching} exacto implementado en Python que identifica los términos mediante coincidencias léxicas exactas con el texto, y tras un preprocesamiento que incluye normalización a minúsculas y tokenización, se han anotado aproximadamente X ejemplos de términos bajo la etiqueta \texttt{term}. A continuación, se ha llevado a cabo una lectura exhaustiva de los X textos que componen el corpus para añadir una segunda capa de etiquetado manual mediante Doccano, una herramienta de anotación de texto de código abierto, diseñada para facilitar la creación de conjuntos de datos etiquetados utilizados en diversas tareas de PLN y aprendizaje automático.

\begin{figure}[H]
    \centering
    \includegraphics[width=0.8\linewidth]{images/sample-doccano.png}
    \caption{Texto anotado en Doccano}
    \label{fig:sample-doccano}
\end{figure}

La anotación del corpus se ha implementado en Doccano como un proyecto de etiquetado secuencial, donde cada unidad terminológica ha sido barcada bajo la etiqueta \texttt{term}. Esta etiqueta se asigna a cualquier unidad léxica que tenga un significado especializado dentro del dominio de la lingüística de corpus. A continuación se describen los criterios para el marcado de unidades como términos.

\begin{enumerate}
    \item Criterios lingüísticos
    \begin{enumerate}
        \item Criterio gramatical. Se consideran términos, principalmente y de acuerdo con lo que establece la bibliografía al respecto, sustantivos o sintagmas nominales que nombra entidades (p. ej., \textit{corpus)}), procesos (p. ej., \textit{anotación}) o propiedades del dominio (p. ej., \textit{representatividad}). También se anotan verbos (p. ej., \textit{tokenizar}) o adjetivos (p. ej., \textit{representativo}) cuando refieran directamente a procedimientos o cualidades, respectivamente. En este sentido, se omiten, inicialmente, palabras funcionales como artículos, preposiciones, etc., a excepción de aquellas que formen parte de un sintagma candidato a término (p. ej., \textit{alineación de corpus paralelos}).
        \item Criterio de dependencia terminológica. Se ha establecido este criterio para la detección de nuevos términos en relación con términos previamente identificados y anotados, estableciendo dependencias tanto paradigmáticas y sintagmáticas como morfológicas. Según este criterio, un candidato se anota como término:
        \begin{itemize}
            \item Si guarda una relación paradigmática con otro término previamente marcado (p. ej., \textit{corpus especializado}/\textit{general}/\textit{de referencia});
            \item Si se da una relación sintagmática con un término que ya ha sido previamente anotado (p. ej., \textit{compilar} un \textit{corpus});
            \item Si es una unidad léxica que mantiene una relación morfológica con un término ya identificado (p. ej., \textit{tokenización}/\textit{tokenizar}).
            \item Si es una unidad léxica que se encuentra incrustrada en una cadena más larga que ha sido previamente marcada como término.
        \end{itemize}
                \item Criterio de variación formal. Se anotan las variantes ortográficas, morfológicas y sintácticas de un mismo término cuando se reconoce su equivalencia conceptual:
        \begin{itemize}
            \item Variación ortográfica: con o sin guion, en una o dos palabras (p. ej., \textit{wordlist}/\textit{word list});
            \item Variación morfológica: singular/plural, adjetival/verbal (p. ej., \textit{lemma}/\textit{lemmas}, \textit{lemmatize}/\textit{lemmatization});
            \item Variación denominativa: reordenamientos que mantienen el significado (p. ej., \textit{corpus linguistics} $\sim$ \textit{linguistics of corpora}).
        \end{itemize}

        \item Criterio de incrustación y jerarquía. Cuando un término compuesto contiene dentro de sí otros términos reconocibles, se anotan tanto la cadena completa como sus subunidades si estas también son relevantes como términos:
        \begin{itemize}
            \item Ejemplo: en \textit{part-of-speech tagging} se anotan tanto \textit{part-of-speech tagging} como \textit{part-of-speech}.
            \item Ejemplo: en \textit{lexical bundle extraction} se anotan tanto \textit{lexical bundle} como \textit{extraction}.
        \end{itemize}
    \end{enumerate}
    \item{Criterios extralingüísticos}
    \begin{enumerate}
        \item Criterio de especificidad o pertenencia al dominio. Se marcará una unidad léxica como unidad terminológica si cumple con los siguientes requisitos en función de su grado de pertenencia al dominio:
        \begin{itemize}
            \item Si es un término específico del dominio (p. ej., \textit{corpus paralelo});
            \item Si es un término importado de otra disciplina periférica pero se usa en el dominio (p. ej., \textit{expresión regular}).
        \end{itemize}

        \item Criterio de frecuencia y estabilidad de uso. Se considerará término anotable si aparece de manera recurrente en la literatura especializada y no constituye un uso aislado o idiosincrático:
        \begin{itemize}
            \item Ejemplo: \textit{keyword} (término estable y recurrente);
            \item No se marcaría una metáfora ocasional como \textit{corpus zoo}.
        \end{itemize}

        \item Criterio de institucionalización. Se anotarán aquellas unidades que, aunque en principio sean descriptivas o metafóricas, se hayan lexicalizado en el discurso disciplinar y funcionen como etiquetas terminológicas:
        \begin{itemize}
            \item Ejemplo: \textit{bag of words}, \textit{data-driven learning}.
        \end{itemize}
    \end{enumerate}
\end{enumerate}


\subsection{Evaluación de la representatividad terminológica del corpus} \label{evaluacion}

Una vez completado el proceso de anotación terminológica del corpus, se procede a su evaluación desde la perspectiva de su grado de representatividad. Esta evaluación no constituye un fin en sí mismo, sino una fase analítica complementaria que permite valorar en qué medida el corpus anotado refleja de forma adecuada el repertorio terminológico característico del dominio de estudio. De forma general, los corpus lingüísticos suelen evaluarse en función de su representatividad, aun tratándose de un concepto escurridizo y de difícil delimitación teórica (\cite{Egbert_Biber_Gray_2022}). No obstante, en el ámbito de los corpus especializados y de las aplicaciones del PLN orientadas al dominio, la evaluación de la representatividad constituye un paso metodológico esencial, en tanto condiciona la validez de los análisis terminológicos posteriores y el rendimiento de los sistemas computacionales entrenados a partir de dichos datos. Más allá de consideraciones tradicionales como el tamaño del corpus, su composición o el equilibrio entre géneros textuales, en este trabajo se adopta una concepción funcional y operativa de la representatividad. Desde esta perspectiva, la representatividad se entiende como la adecuación terminológica del corpus, es decir, como su capacidad para reflejar de forma fiable tanto la terminología central del dominio como el comportamiento léxico estable de dicho repertorio a lo largo del conjunto de textos:

\begin{quote}
    La clave está en lograr un diseño congruente con los objetivos perseguidos en la construcción del corpus. Estos factores son los que establecen el ámbito de validez del corpus y los que, por tanto, pueden permitirnos juzgar su adecuación. (\cite{rojo2021introduccion})
\end{quote}

En coherencia con este planteamiento, la representatividad terminológica se concibe como un constructo multidimensional que no puede reducirse a una única medida. En este sentido, se consideran cuatro dimensiones complementarias que permiten caracterizar de manera empírica y replicable la adecuación del corpus al dominio de estudio: (i) la cobertura terminológica, entendida como el grado en que el corpus incorpora los términos clave del dominio; (ii) la saturación terminológica, que da cuenta del nivel de madurez del inventario especializado conforme se incorporan nuevos textos; (iii) la variedad o consistencia interna del léxico especializado, que permite evaluar la homogeneidad discursiva del corpus; y (iv) la estabilidad terminológica, concebida como la convergencia del inventario de términos independientemente del orden de incorporación de los datos.

\begin{figure}[H]
\centering
\begin{tikzpicture}[
    node distance=0.5cm and 1.2cm,
    every node/.style={font=\footnotesize, align=center},
    central/.style={
        ellipse,
        draw=none,
        fill=gray!15,
        thick,
        minimum width=4cm,
        minimum height=1.5cm,
        font=\small\bfseries
    },
    dim_rect/.style={
        rectangle,
        rounded corners=0pt,
        thick,
        text width=3.8cm,
        minimum height=1.2cm
    }
]

% Nodo Central
\node[central] (centro) {Representatividad};

% Nodos Periféricos (Dimensiones)
\node[dim_rect, fill=blue!15, above left=of centro] (cobertura) {
    \textbf{Cobertura} \\ 
    Presencia del léxico esencial del dominio
};

\node[dim_rect, fill=green!15, above right=of centro] (saturacion) {
    \textbf{Saturación} \\ 
    Madurez del inventario terminológico
};

\node[dim_rect, fill=yellow!25, below left=of centro] (variedad) {
    \textbf{Variedad} \\ 
    Consistencia interna del léxico especializado
};

\node[dim_rect, fill=pink!20, below right=of centro] (estabilidad) {
    \textbf{Estabilidad} \\ 
    Convergencia del inventario terminológico
};

% Flechas de conexión
\draw[thick, ->, >=stealth, gray!60] (cobertura) -- (centro);
\draw[thick, ->, >=stealth, gray!60] (saturacion) -- (centro);
\draw[thick, ->, >=stealth, gray!60] (variedad) -- (centro);
\draw[thick, ->, >=stealth, gray!60] (estabilidad) -- (centro);

\end{tikzpicture}
\caption{Marco conceptual de la representatividad entendida como adecuación terminológica.}
\label{fig:representatividad_conceptual}
\end{figure}

Sobre la base de estas cuatro dimensiones, se propone un marco de evaluación multimétrica que permite cuantificar la representatividad del corpus y su adecuación para tareas terminográficas mediante indicadores empíricos específicos. La integración de estas en un índice compuesto proporciona una medida continua e interpretable de la adecuación del corpus, alineada tanto con objetivos teóricos de la terminología como con requisitos metodológicos propios de aplicaciones computacionales y de PLN orientadas al dominio.


\subsubsection{Tests empíricos}\label{subsubsec:tests-representatividad}

La evaluación de la representatividad terminológica del corpus se articula
mediante cuatro tests empíricos, cada uno asociado a una de las dimensiones
conceptuales definidas en la sección anterior: cobertura, saturación,
dispersión y estabilidad. Su aplicación conjunta permite operacionalizar
de forma complementaria la adecuación terminológica del corpus al dominio
de análisis. Los tests se aplican sobre el corpus crudo procesado
lingüísticamente, con independencia de los umbrales de filtrado propios
de la extracción automática, de forma que los resultados reflejen
propiedades del corpus y no del extractor.

\begin{enumerate}

    \item \textbf{Test de cobertura terminológica.}
    Evalúa en qué medida el corpus incorpora el léxico considerado esencial
    para el dominio, tomando como referencia un estándar externo independiente.
    La cobertura se define como la proporción de términos del conjunto de
    referencia que están presentes en el corpus:
    %
    \begin{equation}
        \mathrm{Cobertura} = \frac{|T_C \cap T_G|}{|T_G|}
        \label{eq:cobertura}
    \end{equation}
    %
    donde $T_C$ representa el conjunto de términos identificados en el corpus
    y $T_G$ el conjunto de términos del estándar externo. Esta formulación
    es equivalente al \textit{recall} en evaluación de sistemas de extracción
    automática y evita el razonamiento circular que supondría evaluar el corpus
    únicamente a partir de sus propiedades internas. Como estándar de referencia
    se utiliza el glosario semilla de 38 unidades terminológicas (Tabla
    \ref{tab:glosario_cass}).

    \item \textbf{Test de saturación terminológica.}
    Analiza la evolución del vocabulario conforme aumenta el tamaño del corpus
    mediante el ajuste de la Ley de Heaps \autocite{heaps1978information}:
    %
    \begin{equation}
        V(n) = K \cdot n^{\beta}
        \label{eq:heaps}
    \end{equation}
    %
    donde $n$ es el número de tokens procesados, $V(n)$ el número de tipos
    léxicos distintos, $K$ una constante de escala y $\beta$ el exponente de
    crecimiento. En corpus especializados se espera que $\beta$ adopte valores
    moderados ($0{,}4 \lesssim \beta \lesssim 0{,}7$): valores bajos indican
    saturación léxica —el vocabulario se estabiliza al añadir nuevos textos—,
    mientras que valores altos indican que el inventario sigue creciendo de
    forma notable. Para su integración en el índice compuesto, $\beta$ se
    normaliza de forma inversa sobre un intervalo empírico
    $[\beta_{\min}, \beta_{\max}] = [0{,}20, 0{,}80]$:
    %
    \begin{equation}
        S_{\beta} = \max\!\left(0,\;\min\!\left(1,\;
        \frac{\beta_{\max} - \beta}{\beta_{\max} - \beta_{\min}}
        \right)\right)
        \label{eq:norm-beta}
    \end{equation}
    %
    de modo que valores altos de $S_\beta$ corresponden a mayor saturación.

    \item \textbf{Test de dispersión terminológica.}
    Evalúa la distribución de los términos principales a lo largo de los
    documentos del corpus mediante la medida $\mathrm{DP}$ de
    \textcite{gries2021dispersion}:
    %
    \begin{equation}
        \mathrm{DP}(w) = \frac{1}{2}\sum_{i=1}^{n}
        \left|\frac{f_i}{F} - s_i\right|
        \label{eq:dp}
    \end{equation}
    %
    donde $f_i$ es la frecuencia del término $w$ en el documento $i$,
    $F$ su frecuencia total en el corpus y $s_i$ la proporción de tokens
    del corpus que corresponde al documento $i$. $\mathrm{DP} \in [0,1]$:
    el valor 0 indica dispersión perfecta (el término se distribuye
    proporcionalmente al tamaño de cada documento) y 1 concentración máxima
    (el término aparece en un único documento). Para la evaluación se
    calcula el $\mathrm{DP}$ medio de los 50 términos con mayor puntuación
    terminológica.

    \item \textbf{Test de estabilidad terminológica.}
    Evalúa la robustez del inventario de términos frente a variaciones en la
    composición del corpus mediante \textit{bootstrapping} de documentos
    \autocite{efron1994introduction}. En cada iteración se genera una muestra con
    reemplazamiento del conjunto de documentos (mismo tamaño que el corpus
    original), se reordena el inventario de términos por frecuencia y se
    registra qué términos se mantienen en el top-$N$. El índice de estabilidad
    se define como la proporción media de términos del top-$N$ original
    que permanecen en el top-$N$ de las muestras bootstrap:
    %
    \begin{equation}
        S_{\mathrm{est}} = \frac{1}{B}\sum_{b=1}^{B}
        \frac{|\hat{T}^{(b)}_N \cap T_N|}{N}
        \label{eq:estabilidad}
    \end{equation}
    %
    donde $T_N$ es el conjunto de los $N$ términos más frecuentes en el corpus
    original, $\hat{T}^{(b)}_N$ el top-$N$ en la $b$-ésima muestra bootstrap
    y $B$ el número total de iteraciones. En este estudio se utilizan
    $B = 200$ iteraciones y $N = 30$.

\end{enumerate}

\begin{table}[H]
\centering
\fontsize{10}{12}\selectfont
\begin{tabular}{>{\raggedright\arraybackslash}p{4.2cm}
                >{\raggedright\arraybackslash}p{4cm}
                >{\raggedright\arraybackslash}p{5.8cm}}
\toprule
\textbf{Test} & \textbf{Métrica} & \textbf{Indica} \\
\midrule
Cobertura terminológica
    & Recall vs.\ estándar externo
    & Proporción de términos de referencia presentes en el corpus \\
\addlinespace
Saturación terminológica
    & Exponente $\beta$ (Ley de Heaps)
    & Velocidad de crecimiento léxico; $\beta$ bajo = corpus saturado \\
\addlinespace
Dispersión terminológica
    & DP de Gries (top-50 términos)
    & Distribución de los términos principales entre documentos \\
\addlinespace
Estabilidad terminológica
    & Índice bootstrap ($B=200$, top-30)
    & Robustez de los rankings terminológicos ante variaciones del corpus \\
\bottomrule
\end{tabular}
\caption{Tests aplicados para la evaluación de la representatividad
         terminológica del corpus.}
\label{tab:tests-representatividad}
\end{table}


\subsubsection{Índice Compuesto de Representatividad Terminológica}
\label{subsubsec:ITR}

A partir de los cuatro tests descritos se construye un índice compuesto que
sintetiza la representatividad terminológica en una única puntuación
interpretable. El \textit{Índice de Representatividad Terminológica} (ITR)
combina versiones normalizadas de cada componente, de modo que valores altos
en todos ellos se traducen en puntuaciones más altas:

\begin{equation}
    \mathrm{ITR} = \frac{S_{\beta} + S_{\mathrm{hápax}}
                        + S_{\mathrm{DP}} + S_{\mathrm{est}}}{4}
                  \times 100
    \label{eq:itr}
\end{equation}

Las transformaciones aplicadas a cada componente garantizan consistencia
direccional: $S_\beta$ se calcula según la ecuación \eqref{eq:norm-beta};
$S_{\mathrm{hápax}}$ penaliza ratios de hápax superiores al 40\,\%
(umbral a partir del cual la falta de masa textual compromete la
confirmación estadística de los términos); $S_{\mathrm{DP}}$ transforma
el DP medio de forma que una mayor dispersión produzca una puntuación
más alta; y $S_{\mathrm{est}}$ corresponde directamente al índice
bootstrap \eqref{eq:estabilidad}. La puntuación resultante oscila entre
0 y 100; los umbrales interpretativos se recogen en la Tabla
\ref{tab:escala_itr}.

\begin{table}[H]
\centering
\fontsize{10}{12}\selectfont
\begin{tabularx}{\textwidth}{c c X}
\toprule
\textbf{ITR} & \textbf{Nivel} & \textbf{Implicaciones} \\
\midrule
$\geq 75$ & Muy representativo
    & Vocabulario saturado, términos bien distribuidos y rankings
      estables; idóneo para terminografía y PLN de dominio. \\
60–74 & Representativo
    & Buen equilibrio entre cobertura y estabilidad; adecuado para
      la mayoría de aplicaciones especializadas. \\
40–59 & Moderadamente representativo
    & Corpus en desarrollo o con léxico disperso; puede requerir
      ampliación o depuración. \\
$< 40$ & Poco representativo
    & Masa textual o estabilidad insuficientes para fines
      terminológicos robustos. \\
\bottomrule
\end{tabularx}
\caption{Escala interpretativa del ITR.}
\label{tab:escala_itr}
\end{table}

Los resultados obtenidos para el subcorpus en inglés se recogen en la Tabla
\ref{tab:itr_summary}. El corpus consta de 50 artículos especializados
(306.288 tokens), de los cuales TECLA extrae 10.012 candidatos terminológicos.

\begin{table}[H]
\centering
\fontsize{10}{12}\selectfont
\begin{tabular}{lcc}
\toprule
\textbf{Métrica} & \textbf{Valor} & \textbf{Interpretación} \\
\midrule
\multicolumn{3}{l}{\textit{Subcorpus en inglés}} \\
\hspace{1em}Saturación — $\beta$ (Heaps) & 0,61
    & Vocabulario en crecimiento moderado-alto \\
\hspace{1em}Masa textual — ratio de hápax & 53,8\,\%
    & Moderada; esperable en 50 artículos académicos \\
\hspace{1em}Dispersión — DP medio (top-50) & 0,56
    & Concentración moderada; homogeneidad temática \\
\hspace{1em}Estabilidad bootstrap (top-30) & 0,80
    & Rankings terminológicos robustos \\
\hspace{1em}Cobertura externa (recall) & 48,7\,\%
    & 18--19 de 38 términos de referencia identificados \\
\rowcolor{gray!20}\hspace{1em}\textbf{ITR} & \textbf{63/100}
    & \textbf{Representativo} \\
\midrule
\multicolumn{3}{l}{\textit{Subcorpus en español}} \\
\hspace{1em}Saturación — $\beta$ (Heaps) & 0,68 & Vocabulario en crecimiento moderado-alto \\
\hspace{1em}Masa textual — ratio de hápax & 41,0\,\% & Moderada; típica en corpus académico \\
\hspace{1em}Dispersión — DP medio (top-50) & 1,00 & Concentración moderada; homogeneidad temática \\
\hspace{1em}Estabilidad bootstrap (top-30) & 1,00 & Rankings terminológicos robustos \\
\hspace{1em}Cobertura externa (recall) & 63,2\,\% & 24--25 de 38 términos de referencia identificados \\
\rowcolor{gray!20}\hspace{1em}\textbf{ITR} & \textbf{73/100} & \textbf{Representativo} \\
\bottomrule
\end{tabular}
\caption{Métricas de representatividad terminológica por subcorpus.}
\label{tab:itr_summary}
\end{table}

Con una puntuación de 63 sobre 100, el subcorpus en inglés se sitúa en el
nivel \textit{representativo}. El valor $\beta = 0{,}61$ indica que el
inventario léxico aún no ha alcanzado la meseta de saturación, lo que sugiere
que una ampliación del corpus mejoraría la cobertura terminológica. No
obstante, el índice de estabilidad bootstrap ($0{,}80$) evidencia que los
términos centrales del dominio se identifican de forma consistente
independientemente de la composición concreta del corpus, lo que proporciona
una base terminológica fiable para los análisis que se presentan en las
secciones siguientes. La cobertura externa del 48,7\,\% frente al glosario
semilla confirma que el corpus recoge aproximadamente la mitad del repertorio
terminológico nuclear del dominio; los términos no identificados corresponden
mayoritariamente a unidades de baja frecuencia o alto grado de generalidad
que el extractor filtra por diseño.

\section{Extracción automática de términos}\label{sec:ate}

En este dinámico terreno dinámico en el que surgen nuevos términos constantemente, resulta primordial la automatización de tareas mediante la integración de recursos tecnológicos en el proceso de identificación y vaciado de términos en textos especializados. A finales del siglo XX, con la creciente disponibilidad de corpus en formato electrónico y el crecimiento exponencial de la cantidad de información textual en ámbitos especializados, se comienzan a desarrollar enfoques más sofisticados basados en el procesamiento del lenguaje natural. En este contexto, surge la extracción automática de términos (del inglés, ATE o \emph{Automatic Term Extraction}) como consecuencia de la automatización del vaciado de unidades terminológicas a partir del análisis computacional de corpus especializados, con estudios pioneros (Bourigault, 1994; Jacquemin y Royauté, 1994; Estopà, 1999) y el desarrollo de los primeros extractores automáticos de base lingüística (Estopà, 1999).

Bajo términos anglófonos como \emph{Automatic Term Extraction}, \emph{Automatic Term Recognition} (\cite{acosta2015reconocimiento,astrakhantsev2015methods}) o \emph{Terminology Mining} (\cite{daille2003conceptual}) se engloba esta tarea enmarcada en el procesamiento del lenguaje natural desde la década de los noventa (\cite{heylen2015automatic}) cuyo objetivo principal es la extracción de una lista de candidatos a término mediante el vaciado de un corpus especializado a través de soportes informáticos, de manera que tiene una base cuantitativa, objetiva, empírica y automatizada que no encontramos en los métodos manuales tradicionales. La automatización y sistematización de esta tarea surge como respuesta a la necesidad de acceder a la gran cantidad de datos contenidos en los textos digitales de las áreas especializadas con el fin de agilizar el trabajo de identificación de términos y la construcción del léxico especializado de un dominio (\cite{acosta2015reconocimiento,heylen2015automatic})), puesto que la extracción manual, además de ser más laboriosa, se limitaba a la selección de candidatos a término bajo el juicio del terminológo y la ulterior consulta con expertos (Heylen y de Hertog, 2015; Rodríguez-Tapia, 2024). La relevancia de la extracción automática de terminología radica, además de en su base cuantitativa y empírica -en oposición a un fundamento manual y perceptivo-, en su función de punto de partida, prerrequisito (Acosta, Aguilar, et al., 2015) o punto central de la cadena de trabajo en terminología (Cabré, 2008a: 169) con respecto a otras grandes aplicaciones de la terminología moderna como la elaboración de recursos lexicográficos (p. ej., glosarios, tesauros, fichas terminológicas y bancos de datos terminológicos) o el desarrollo de sistemas y aplicaciones del procesamiento del lenguaje natural (p. ej., traducción especializada, indexación automática de documentos especializados, etc.).

La detección de términos en corpus especializados se ha bifurcado en dos etapas (Cabré, 2008a: 169) o, mejor dicho, direcciones que se corresponden con la distinción que ya hemos establecido entre \textit{terminótica} y \textit{terminología computacional}: por un lado, (i) el uso de herramientas de gestión de corpus para la detección de candidatos a término y, por otro, (ii) el desarrollo de sistemas de extracción automática de terminología, respectivamente. Es de nuestro interés la segunda dirección orientada a la creación de recursos y sistemas terminológicos. Sin embargo, a pesar de haber sido esta abordada desde la informática y otras disciplinas cercanas, se le ha dedicado escasa atención desde una perspectiva lingüística, especialmente en lengua española, hecho que resulta desconcertante considerando el papel central de la lingüística en esta tarea, pues aquellos métodos de extracción terminológica que no se basan en la formalización de reglas lingüísticas tienden a ser limitados en términos de precisión y eficiencia, evidenciando aún más el hecho de que estas aplicaciones propias del procesamiento del lenguaje natural requieran de presupuestos lingüísticos teóricos para una descripción y representación adecuadas de la terminología. Este descuido por parte de la lingüística moderna puede deberse, en parte, a la fluctuante falta de comunicación entre lingüística computacional y lingüística teórica, de larga tradición en el panorama lingüístico (Moreno Sandoval, 2016). Precisamente, la extracción automática de terminología se encuentra en una posición singular al estar en la intersección de estas dos disciplinas: por un lado, resulta teórica, ya que requiere de fundamentos y caracterizaciones lingüísticas sólidos sobre las unidades terminológicas; mientras que, por otro lado, forma parte de la lingüística computacional, dado que sus modelos, métodos y sistemas resultan extremadamente útiles para abordar problemas concretos en la extracción, gestión y representación de la terminología.

\subsection{Bases conceptuales y metodológicas} \label{subsec:ate_bases}

En el ámbito de la extracción automática de terminología, es importante destacar la definición de dos términos clave (Kageura y Umino 1996): \textit{unithood} y \textit{termhood}, propiedades inherentes a los términos ampliamente exploradas en la literatura científica centrada en la extracción automática de terminología (Vivaldi et al., 2011; Ananiadou y Macnaught, 2005; Acosta, Aguilar et al., 2015; Kockaert y Steurs, 2015). Para ambos términos proponemos una traducción adaptada a la lengua española: \emph{grado de estabilidad} y \emph{grado de especialización}.

\begin{enumerate}[label=(\arabic*)]
    \item El grado de estabilidad sintagmática o \emph{unithood} atiende al grado de fuerza o estabilidad de las combinaciones sintagmáticas y colocaciones (\cite{kageura1996}), es decir, al grado en el que una secuencia de palabras constituye una unidad lingüística cohesionada y estable\footnote{El hecho de que una cadena de texto tenga cierta estabilidad y unidad no explica que esta tenga un alto grado de tecnicidad: si consideramos la cadena ‘por ejemplo’ a lo largo de un corpus especializado, observaremos que esta tiene un alto nivel de estabilidad puesto que ‘por’ y ‘ejemplo’ son unidades léxicas que suelen coocurrir juntas; sin embargo, no tiene un alto grado de especialización porque no representa ningún concepto específico del dominio especializado.}. Este es un criterio sumamente significativo a la hora de establecer un primer filtrado de candidatos a término, ya que es posible encontrar de manera automática cadenas de texto que supongan una combinación estable que conformen unidades léxicas candidatas a término. Normalmente, esta propiedad de los términos se trabaja a partir de heurísticas lingüísticas (\cite{acosta2015reconocimiento}) basadas en el conocimiento previo del dominio y de la lengua de trabajo
    \item El grado de especialización, la terminologicidad o \emph{termhood} se corresponde con el grado en que una unidad léxica estable está relacionada con algunos conceptos específicos de un dominio (\cite{kageura1996methods}), o el proceso por el cual una palabra o grupo de palabras adquieren un sentido técnico o científico en contextos específicos (\cite{depecker2003saussure}). Se asume que hay marcadores cuantitativos para los términos, y que estos son ‘inusuales’, en el sentido de que tienen una propiedad estadística que puede ayudar a distinguirlos de las palabras de la lengua común. Esta propiedad se mide, por lo general, a partir de criterios estadísticos, partiendo de la premisa de que aquellos términos específicos de un dominio tendrán una frecuencia más alta en dicho dominio que en otros dominios o en un corpus de la lengua común. Esto puede determinarse de dos maneras (\cite{nazar2016distributional}):
        \begin{itemize}
            \item mediante una clasificación binaria, separando las unidades léxicas del corpus en términos y no términos, o
            \item mediante un \textit{continuum} de candidatos a término, ordenados según un puntaje o ponderación que indica la probabilidad de que un candidato sea un término.
        \end{itemize}

\begin{table}[H]
\small
\centering
\renewcommand{\arraystretch}{1.4}
\begin{tabular}{>{\centering\arraybackslash}m{3cm} m{3.5cm} m{5.5cm}}
\hline
\textbf{Tipo de clasificación} & \textbf{Descripción} & \textbf{Métodos más representativos} \\
\hline
\textbf{Binaria} & Asigna a una unidad una de estas clases: \emph{término} \textbar{} \emph{no término} & Clasificadores supervisados (SVM, Random Forest), redes neuronales, etiquetas BIO en NER \\
\textbf{Ponderada} & Asigna a una unidad un valor numérico que representa su grado de tecnicidad & \emph{C-value}, \emph{NC-value}, \emph{Weirdness Index}, \emph{TF-IDF} \\
\hline
\end{tabular}
\caption{Tipos de clasificación en la extracción automática de términos}
\label{tab:tipos-clasificacion}
\end{table}


El análisis de estas propiedades resulta fundamental para el proceso de extracción automática de terminología, ya que permite identificar de manera precisa las unidades léxicas con potencial para convertirse en términos dentro de un dominio de especialidad: mientras que el grado de estabilidad facilita la identificación de combinaciones de palabras estables, el grado de especialización distingue aquellas unidades léxicas que adquieren un significado especializado en un contexto específico. Las diversas estrategias y métodos para detectar y medir ambas propiedades se detallarán a continuación.

\end{enumerate}

\subsection{Métodos} \label{subsec:ate_metodos}

Un extractor de terminología es una herramienta diseñada específicamente para identificar términos en un corpus especializado (Bernier-Colborne y Drouin 2014). No existe actualmente consenso ni estándar en torno a las directrices para construir un extractor automático de terminología óptimo, pues sus criterios, ajustes y funciones son susceptibles a variar en función de la aplicación específica para la que se pretende realizar la extracción de términos (Heylen y De Hertog 2015). Lo que sí podemos afirmar es que la extracción automática de términos es una línea de investigación que atrae tanto a terminológicos más vinculados a corrientes de lingüística teórica como a aquellos enfocados en la práctica terminográfica y estandarización de terminología (Nazar y Cabré 2011), trabajando arduamente en el diseño de métodos automáticos o, al menos, semi-automáticos, para asistir a los expertos en conocimiento en el proceso de catalogación de los términos y conceptos de un dominio en forma de vocabularios, tesauros, terminologías y ontologías.

Los métodos para llevar a cabo la extracción automática de terminología son variados (Astrakhantev, Fedorenko et al., 2015) y pueden categorizarse en cuatro enfoques principales (Lossio-Ventura et al., 2016): métodos estadísticos, métodos lingüísticos, métodos híbridos y, más recientemente, métodos de aprendizaje supervisado, aunque dichos enfoques no son excluyentes y, en muchos casos, se combinan para obtener resultados más precisos y completos, puesto que ninguno de ellos, considerado de forma independiente, ha conseguido resultar del todo exitoso (Acosta, Aguilar, et al., 2015).

\subsubsection{Métodos estadísticos} \label{subsec:stat_methods}

El enfoque estadístico para la extracción de términos no requiere, \textit{a priori}, información lingüística para la identificación de candidatos a términos, pues se basa en la idea de que las unidades especializadas a menudo muestran una alta frecuencia en textos relacionados con un dominio específico. Estos métodos se centran, generalmente, en los análisis de frecuencias bien de manera interna con respecto al corpus o bien mediante el contraste con un corpus de referencia, que puede ser de carácter general o de otro dominio distinto (\cite{cabre2008como}).

Estas mediciones estadísticas pueden aplicarse a la extracción de candidatos a término con estabilidad sintagmática. Es el caso de métricas de asociación como el estadístico chi cuadrado (χ²) o la Pointwise Mutual Information (PMI), que permiten detectar combinaciones de palabras cuya coocurrencia es significativamente superior a la esperable por azar. En el ámbito de la extracción terminológica, el estadístico χ² se emplea para evaluar si la coocurrencia de dos unidades léxicas no puede explicarse por independencia estadística, lo que sugiere que forman una unidad relativamente estable. De manera complementaria, la PMI mide el grado de dependencia entre dos palabras comparando la probabilidad observada de su coocurrencia con la probabilidad esperada bajo una hipótesis de independencia; valores elevados de PMI indican una asociación fuerte entre los elementos, lo que puede interpretarse como un indicio de fijación sintagmática en contextos especializados. Aunque efectivas para detectar asociaciones entre palabras que coocurren con frecuencia en un corpus y podrían indicar la formación de unidades terminológicas estables, estas métricas se limitan a medir la relación entre las palabras sin considerar su relevancia o grado de especialización (\emph{termhood}). Para abordar esta limitación, existen numerosas técnicas (Andrade et al., 2023) para la interpretación estadística de los términos. Una de ellas es el enfoque distribucional o TF-IDF, un método clásico en el ámbito de la extracción de información (\cite{sparck1972statistical, tran2023recent}). Este método permite ponderar la relevancia de una unidad léxica —frecuentemente, un candidato a término— en un documento específico, en función de su distribución en el conjunto total de documentos del corpus. Así, se otorga mayor peso a las unidades que son frecuentes en un documento particular pero escasamente frecuentes en el resto del corpus, lo cual suele ser indicativo de su especialización temática.

\equationcaption{TF-IDF}{
\mathrm{TF\text{-}IDF}(t,d,D) = \mathrm{tf}(t,d) \cdot \mathrm{idf}(t,D)
}


El valor TF-IDF de un término $t$ en un documento $d$ respecto a un corpus $D$ se obtiene multiplicando dos componentes:


\begin{itemize}
    \item La frecuencia del término (\emph{TF}), que mide cuántas veces aparece $t$ en el documento $d$, en proporción al número total de palabras de ese documento:
    \equationcaption{TF}{
\mathrm{tf}(t,d) = \frac{f_{t,d}}{n_d}
}
\label{eq:tf}

    donde $f_{t,d}$ es la frecuencia absoluta del término $t$ en el documento $d$ y $n_d$ es el número total de palabras en ese documento.

    \item La frecuencia inversa de documentos (\emph{IDF}), que mide la dispersión del término $t$ a lo largo del corpus:
    \equationcaption{IDF}{
\mathrm{idf}(t,D) = \log \left( \frac{N}{1 + df_t} \right)
}
\label{eq:idf}

    donde $N$ es el número total de documentos del corpus y $df_t$ es la cantidad de documentos en los que aparece al menos una vez el término $t$.
\end{itemize}
La lógica subyacente es que una unidad léxica muy frecuente en todos los documentos (por ejemplo, preposiciones o conjunciones) no aporta información distintiva sobre el contenido temático de un texto y, por tanto, no debe ser considerada término. En cambio, las unidades que aparecen reiteradamente en documentos concretos pero no en otros se consideran potencialmente relevantes para el dominio de especialidad. En este sentido, un valor alto de TF-IDF sugiere un mayor grado de especialización o \emph{terminologicidad} (\cite{rodtapia2024}) del término en cuestión.

\subsubsection{Métodos lingüísticos}

Los extractores automáticos basados en una metodología exclusivamente lingüística aprovechan el empleo del conocimiento explícito acerca de la lengua y el dominio temático hacia los que está enfocada la extracción (Nazar y Cabré, 2011), utilizando filtros de información lingüística en forma de reglas y gramáticas elaboradas, generalmente, de manera manual (Lossio Ventura et al., 2016) para restringir lo que se reconoce como término y lo que no. Un método bien asentado para la detección de términos poliléxicos o \emph{multi-word terms} es el uso de plantillas morfosintácticas específicas para identificar unidades terminológicas en estructuras recurrentes que son frecuentes en la terminología de diversas disciplinas.

\begin{table}[H]
\centering
\begin{tabular}{@{}p{7cm}p{4cm}@{}}
\toprule
\textbf{Plantilla morfosintáctica} & \textbf{Ejemplo} \\
\midrule
Sustantivo + Sustantivo & \textit{corpus linguistics} \\
Adjetivo + Sustantivo & \textit{raw data} \\
Sustantivo + Sustantivo + Sustantivo & \textit{type token ratio} \\
Sustantivo + Preposición + Sustantivo & \textit{part of speech} \\
\bottomrule
\end{tabular}
\caption{Ejemplo de plantillas morfosintácticas.}
\label{tab:plantillas-morfosintacticas}
\end{table}

Otro componente esencial en la extracción de términos mediante métodos lingüísticos son las listas de exclusión o stoplists. Estas listas contienen palabras funcionales o genéricas que no aportan ningún tipo de información significativa para el dominio específico. Por lo general, este tipo de listados incluyen artículos, preposiciones, pronombres, y otros elementos que, aunque frecuentes, no tienen un valor especializado. Estas listas son esenciales para reducir el ruido en el proceso de extracción, ayudando a filtrar palabras que podrían generar resultados erróneos al identificar estas unidades como candidatos a término. A medida que se trabaja con el sistema, las stoplists se refinan y actualizan continuamente para mejorar la precisión de la extracción.

Tanto la detección automática de patrones morfosintácticos como la supresión de determinadas unidades léxicas a través de la generación y/o implementación de una stoplist pueden llevarse a cabo mediante algoritmos del procesamiento del lenguaje natural (PLN) y diversas técnicas de análisis estadístico, que aprovechan herramientas y librerías especializadas para facilitar y optimizar este tipo de procesos. Algunas de las librerías más utilizadas para implementar estos algoritmos son NLTK\footnote{<\hyperlink{https://www.nltk.org/}{https://www.nltk.org}>} o spaCy\footnote{<\hyperlink{https://spacy.io/}{https://spacy.io/}>}, ambas con una amplia oferta de funciones completas para el análisis de corpus como tokenización, etiquetado de partes del discurso, análisis sintáctico y creación de stoplists.

Si bien los métodos lingüísticos ofrecen un enfoque preciso en dominios bien definidos, su principal desventaja radica en su limitada generalización. Al depender de reglas específicas para una lengua y un dominio particular, no son domain-independent ni language-independent, lo que significa que su implementación en otros contextos requiere un esfuerzo considerable de adaptación. Además, estos métodos tienden a generar bastante ruido debido a las limitaciones en la cobertura de las reglas lingüísticas establecidas. A pesar de estas limitaciones, los métodos lingüísticos son más que útiles para escenarios donde se dispone de conocimiento detallado del dominio.

\subsubsection{Métodos híbridos}\label{subsec:hybrid_methods}

En vista a las limitaciones de los extractores anteriores, especialmente aquellos puramente lingüísticos o estadísticos, este tipo de extractor combina métodos lingüísticos y estadísticos para proporcionar mejores resultados, siendo los más extendidos actualmente en el panorama de la extracción automática de términos (Heylen y De Hertog, 2015; Andrade et al., 2023). Por lo general, sigue dos fases consecutivas: la extracción de cadenas de texto utilizando filtros lingüísticos y el posterior filtrado de estas cadenas candidatas a término utilizando métodos estadísticos. Se incorporan, además, métodos adicionales para mejorar la precisión del extractor (Zhang et al. 2018).

\begin{itemize}
    \item C-value (Frantzi et al., 1999) es el método híbrido preferente por los investigadores (Andrade et al., 2023). Este método aborda la extracción de unidades terminológicas poliléxicas o multi-word terms (Heylen \& De Hertog, 2015) mediante la determinación del grado de tecnicidad de estas unidades a través de mediciones estadísticas (p. ej., la frecuencia de aparición de la cadena completa de texto, sus componentes por separado, la longitud de la cadena de texto, etc.). Este enfoque es contextual, es tanto que tiene en cuenta el contexto lingüístico, aplica un filtro lingüístico -etiquetado gramatical, uso de una stoplist (Frantzi et al., 1998)- para delimitar y reducir las unidades que van a procesarse como candidatos a término y luego analiza cómo las unidades del corpus se relacionan entre sí. Este algoritmo asume, entonces, las siguientes premisas:
    \begin{itemize}
        \item que los términos son unidades poliléxicas (\emph{multi-word terms}), y
        \item que los términos pueden aparecer dentro de cadenas candidatas a término más largas, es decir, que están anidados (\emph{nested terms})
    \end{itemize}

A partir de estas apreciaciones, se expone una relación aleatoria de las siguientes variables (Rodríguez-Tapia, 2024: 143):

\begin{enumerate}
    \item la frecuencia absoluta del candidato a término a lo largo de todo el corpus;
    \item la frecuencia del candidato a término como parte o constituyente de un una cadena candidata más larga;
    \item el número de candidatos a término más largos, y
    \item la longitud del candidato a término, por número de caracteres o de tokens que lo componen.
\end{enumerate}

Este algoritmo se recoge en la siguiente fórmula, donde:

\begin{itemize}
    \item \emph{a} es la cadena candidata a término;
    \item \emph{f} es la frecuencia de aparición de la cadena candidata en el corpus;
    \item \emph{T\textsubscript{a}} es el conjunto de candidatos a término extraídos que contienen la cadena \emph{a};
    \item \emph{P(T\textsubscript{a})} es el número de candidatos a término en el conjunto \emph{T\textsubscript{a}}, y
    \item $|a|$ representa la longitud de la cadena \emph{a} en número de palabras.
\end{itemize}

\equationcaption{C-value}{
\mathrm{C\text{-}value}(a) =
\begin{cases}
\log_2 |a| \cdot f(a) & \text{si } a \text{ no está anidada}, \\\\[6pt]
\log_2 |a| \cdot \left( f(a) - \frac{1}{P(T_a)} \sum_{b \in T_a} f(b) \right) & \text{si } a \text{ está anidada}.
\end{cases}
}
\label{eq:cvalue}



El algoritmo se bifurca en dos direcciones, en función de dos condiciones distintas: a) la cadena candidata a no está anidada (primera línea) o b) si está anidada (segunda línea):

\begin{itemize}
    \item Para las cadenas no anidadas, es decir, aquellas que no aparecen como parte de otros términos candidatos, el C-value se calcula simplemente multiplicando la longitud de la cadena por su frecuencia en el corpus, asumiendo que las cadenas más largas y frecuentes tienen una mayor probabilidad de ser términos válidos.
    \item Sin embargo, para las cadenas anidadas, que sí aparecen como parte de otros términos candidatos, se tiene en cuenta tanto la frecuencia de la cadena como parte de otros términos como el número de esos términos. Se resta la frecuencia de la cadena como parte de otros términos de su frecuencia total para penalizar las cadenas que solo aparecen como parte de otras cadenas. Al mismo tiempo, se divide la frecuencia de la cadena como parte de otros términos por el número de esos términos para considerar la independencia de la cadena. Una cadena que aparece como parte de muchas cadenas diferentes tiene más probabilidad de ser un término válido, incluso si no aparece por sí sola.
\end{itemize}

Esto quiere decir que, si un candidato a término aparece como parte de un término más largo, se considera una secuencia candidata a término incompleta y recibe, por lo tanto, una ponderación de C-value bajo. A su vez, cuanto mayor sea el número de candidatos a término distintos en los que el candidato a término aparezca anidado, más probable es que sea un término. Por ejemplo, en el dominio de la lingüística, 'linguistic' aparece repetidamente en varias combinaciones como '\textit{linguistic analysis}', '\textit{linguistic structure}' o '\textit{linguistic patterns}'. En este caso, '\textit{linguistic}' no se considera un término independiente, ya que generalmente aparece como parte de términos más largos y no ocurre de forma autónoma en el corpus. Es por esta relación de dependencia que por la que '\textit{linguistic}' recibe un C-value bajo. Por otro lado, términos como '\textit{collocational patterns}', aunque también puedan aparecer dentro de secuencias más largas (como '\textit{collocational patterns of verbs}' o '\textit{collocational patterns in academic texts}'), se consideran términos completos porque ocurren con frecuencia como unidades independientes. Este comportamiento, en el que el término aparece de forma recurrente en diferentes contextos sin depender de otros términos para su validez, le otorga un C-value alto. En resumen,

\begin{quote}
\textit{the more frequently a candidate term appears as a substring of other candidates, the less likely it is to be an actual term. At the same time, the larger the number of distinct term candidates in which the target term candidate oceurs as nested, the more likely it is to be an actual term. }(Korkontzelos y Ananiadou, 2022: 996-997)

\end{quote}

\item NC-value, por su parte, es una extensión de C-value que mejora la extracción automática de términos al incorporar información contextual sobre los términos poliléxicos (Frantzi et al., 2000). Sigue la siguiente lógica:
\begin{itemize}
    \item El algoritmo parte de una lista ya generada de candidatos a término ordenados según su puntaje o ponderación de C-value. 
    \item De los términos que componen el listado se extraen palabras de su cotexto asignándoles un peso que refleja su tendencia a aparecer junto con potenciales términos. Este peso se calcula teniendo en cuenta la frecuencia de aparición de cada palabra de contexto junto a los términos candidatos, permitiendo resaltar aquellas palabras que son indicativas de la relevancia de un término dentro de un determinado contexto.
    \item Las palabras que rodean al candidato a término y sus respectivos pesos se utilizan para reordenar la lista de términos candidatos, integrando tanto la puntuación de C-value como la nueva información contextual. Este reordenamiento asegura que los términos realmente relevantes se posicionen en las primeras posiciones de la lista, mejorando así el rendimiento de la extracción de términos poliléxicos.
\end{itemize}

\equationcaption{NC-value}{
\mathrm{NC\text{-}value}(a) = 0.8 \cdot \mathrm{C\text{-}value}(a) + 0.2 \sum_{b \in C_a} f_a(b) \cdot \mathrm{weight}(b)
}
\label{eq:ncvalue}


donde:

\begin{itemize}
    \item $\mathrm{C\text{-}value}(a)$ es el valor \emph{C-value} del término candidato \emph{a};
    \item $C_a$ es el conjunto de palabras de contexto asociadas al candidato a término \emph{a};
    \item $f_a(b)$ es la frecuencia con la que la palabra de contexto \emph{b} aparece junto al candidato \emph{a};
    \item $\mathrm{weight}(b)$ es el peso asignado a la palabra de contexto \emph{b}.
\end{itemize}

El primer término\footnote{En este caso utilizamos término en su uso técnico referido a un componente o elemento de una expresión matemática, no como sinónimo de unidad terminológica.
} de la fórmula representa la ponderación del valor original del término, mientras que el segundo término corresponde al factor contextual que ajusta la relevancia del término según aquellas palabras que lo rodean. Estos dos factores, \textit{C-value} y el factor contextual, se ponderan con los valores 0.8 y 0.2, respectivamente (Frantzi et al., 2000: 126). Este proceso permite que los candidatos a término más relevantes queden reordenados, favoreciendo aquellos con una mayor concentración de palabras de contexto significativas.


\end{itemize}

\subsubsection{Métodos de aprendizaje automático}

En el ámbito del PLN se ha establecido de manera generalizada que muchas tareas de análisis lingüístico pueden formularse como problemas de clasificación supervisada, paradigma que se ha trasladado, de manera natural, a la extracción automática de términos. Los métodos de aprendizaje automático abordan la extracción automática de términos formulándola explícitamente como una tarea de clasificación, generalmente supervisada, en la que el sistema aprende a distinguir entre unidades terminológicas y unidades léxicas no especializadas a partir de ejemplos previamente anotados (\cite{nazar2016distributional,lossio2016biomedical}). A diferencia de los métodos estadísticos e híbridos clásicos, estos enfoques no se limitan a asignar una puntuación de relevancia a los candidatos, sino que aprenden patrones lingüísticos, distribucionales y contextuales asociados a la terminologicidad. Desde un punto de vista formal, la extracción automática de términos mediante aprendizaje automático puede plantearse de dos maneras principales. Por un lado, puede abordarse como una tarea de clasificación, ya sea binaria —en la que se distingue entre término y no término— o gradual, en la que las unidades se ordenan a lo largo de un \textit{continuum} según su grado de \textit{termhood} o terminologicidad (\cite{nazar2011}). Por otro lado, cuando el objetivo es identificar términos en texto continuo y, especialmente, delimitar unidades poliléxicas, la tarea suele operacionalizarse como un problema de etiquetado secuencial a nivel de token, mediante esquemas como BIO (\textit{Beginning–Inside–Outside}). Este tipo de codificación no define el criterio de terminologicidad en sí mismo, sino que proporciona un mecanismo formal para representar los límites y la estructura interna de las unidades terminológicas y facilitar su reconstrucción automática.

En los enfoques supervisados basados en rasgos o \textit{feature-based}, el aprendizaje automático se apoya en una representación explícita de las unidades lingüísticas mediante vectores de características diseñadas manualmente. Estas características suelen incluir información lingüística (categoría gramatical, lema, estructura morfosintáctica), ortográfica (capitalización, longitud, forma del token), contextual (tokens vecinos, pertenencia a sintagmas nominales) y estadística (frecuencia, TF-IDF, medidas de dispersión). Esta representación permite al modelo capturar propiedades observables que, de forma conjunta, se consideran indicativas de terminologicidad. Sobre estos vectores de rasgos se entrenan clasificadores clásicos del aprendizaje automático, como Máquinas de Vectores de Soporte (SVM), Árboles de Decisión o Random Forest, que se han empleado principalmente para la validación o clasificación de candidatos a término, especialmente en el caso de términos uniléxicos o cuando la unidad de decisión es el candidato completo. En cambio, cuando la extracción se plantea como una tarea de etiquetado secuencial orientada a la detección de términos poliléxicos, los Campos Aleatorios Condicionales (CRF) se han consolidado como uno de los modelos más adecuados, al permitir modelar dependencias entre etiquetas consecutivas y capturar regularidades secuenciales propias de las unidades terminológicas (Peng et al., 2004; Zhang et al., 2010).

No obstante, estos métodos presentan limitaciones importantes. Requieren corpus anotados de alta calidad, cuyo diseño y validación suponen un coste considerable, y su rendimiento depende en gran medida de la selección y formalización de las características empleadas. A ello se suma que la ingeniería manual de rasgos puede introducir sesgos teóricos o limitar la portabilidad del sistema a otros dominios o lenguas. Pese a estas restricciones, los métodos supervisados basados en rasgos continúan siendo una referencia sólida en terminología computacional, tanto para la evaluación controlada de la terminologicidad como para el análisis interpretativo de los factores lingüísticos que intervienen en la identificación de términos.

\subsubsection{Métodos de aprendizaje profundo}

Los métodos de aprendizaje profundo representan la evolución más reciente en la extracción automática de términos, caracterizándose por prescindir en gran medida de la ingeniería manual de rasgos y aprender representaciones distribuidas directamente a partir de los datos. Estos enfoques se apoyan en arquitecturas neuronales capaces de modelar relaciones complejas y dependencias de largo alcance en el texto, lo que resulta especialmente relevante para la identificación de unidades terminológicas en contextos especializados. Las primeras aproximaciones neuronales a la extracción terminológica emplearon redes neuronales recurrentes (RNN) y redes convolucionales (CNN), combinadas a menudo con capas CRF para el etiquetado secuencial. Estas arquitecturas permitieron mejorar la detección de términos poliléxicos y reducir la dependencia de reglas lingüísticas explícitas, aunque seguían requiriendo cantidades considerables de datos anotados y mostraban dificultades para generalizar entre dominios.

Un punto de inflexión lo constituye la introducción de modelos basados en \textit{Transformers} (§ \ref{subsec:neural_models}), como BERT y sus variantes, que han demostrado un rendimiento sobresaliente en tareas de etiquetado secuencial, incluido el Reconocimiento de Entidades Nombradas (NER). En este marco, la extracción automática de términos se reformula de manera natural como una tarea de NER, donde los términos especializados se tratan como una categoría de entidad definida por el investigador. Estudios recientes aplicados al dominio de la lingüística han mostrado que el ajuste fino (\textit{fine-tuning}) de modelos \textit{Transformer} sobre corpus terminológicos anotados permite alcanzar resultados competitivos, incluso en ámbitos tradicionalmente poco explorados desde el punto de vista terminográfico. Sin embargo, estos enfoques no están exentos de problemas metodológicos. El alto coste computacional, la necesidad de grandes volúmenes de datos anotados y el carácter opaco de las representaciones internas dificultan tanto la reproducibilidad como la interpretación lingüística de los resultados. Además, existe el riesgo de que los modelos neuronales aprendan patrones de superficie o dependan excesivamente de la frecuencia de los términos más recurrentes, sin capturar plenamente los principios que subyacen a la terminologicidad.

En este contexto emergen los grandes modelos de lenguaje (LLM) como una alternativa novedosa para la extracción automática de términos. A diferencia de los enfoques anteriores, los LLM permiten abordar la tarea sin entrenamiento específico, mediante estrategias de prompting en configuraciones \textit{zero-shot} o \textit{few-shot}. Investigaciones recientes aplicadas a la lingüística de corpus han demostrado que modelos como GPT-3.5 y GPT-4 pueden extraer listas de términos con un rendimiento competitivo, especialmente cuando se proporcionan ejemplos mínimos que orientan la noción de término buscada. Estos resultados sugieren que los LLM capturan de forma implícita patrones semánticos y contextuales asociados a la terminología especializada, probablemente adquiridos durante su entrenamiento masivo. No obstante, el uso de LLMs en extracción terminológica plantea desafíos específicos, como la variabilidad en las salidas, la falta de control sobre los criterios internos de selección y la necesidad de estrategias de evaluación adaptadas que tengan en cuenta la variación morfológica, la inclusión de subcadenas y la proximidad semántica. A pesar de ello, los enfoques basados en prompting abren una vía prometedora para la extracción terminológica asistida, especialmente en escenarios de bajo recurso o como complemento a métodos supervisados tradicionales.

\subsection{Flujo de trabajo y evaluación}

En la actualidad, la mayoría de los sistemas de extracción automática de términos son híbridos, combinando tanto información lingüística como estadística (Heylen y De Hertog, 2015) y generando mejores resultados que aquellos extractores únicamente estadísticos o lingüísticos (Macken et al., 2013). Por lo general, la extracción se divide en las siguientes subtareas:

\begin{figure}[H]
\centering
\begin{tikzpicture}[node distance=0.4cm and 0.4cm]
% Estilos de nodos
\tikzstyle{entrada} = [rounded corners, minimum width=2.4cm, minimum height=1.1cm, align=center, fill=gray!15, font=\scriptsize]
\tikzstyle{proceso} = [rectangle, minimum width=2.8cm, minimum height=1.1cm, align=center, fill=blue!10, font=\scriptsize]
\tikzstyle{evaluacion} = [rectangle, minimum width=2.8cm, minimum height=1.1cm, align=center, fill=green!10, font=\scriptsize]
\tikzstyle{arrow} = [-{Latex}, thick]

% Nodos
\node[entrada] (input) {Corpus};
\node[proceso, right=of input] (prepro) {Procesamiento\\
$\bullet$ Tokenización\\
$\bullet$ Caracteres alfabéticos\\
$\bullet$ Etiquetado gramatical};
\node[proceso, right=of prepro] (filtro) {Filtro lingüístico};
\node[evaluacion, right=of filtro] (medida) {Medición\\ del grado\\ de tecnicidad};
\node[entrada, right=of medida] (output) {Lista de\\ candidatos\\ a término};

% Flechas
\draw[arrow] (input) -- (prepro);
\draw[arrow] (prepro) -- (filtro);
\draw[arrow] (filtro) -- (medida);
\draw[arrow] (medida) -- (output);

\end{tikzpicture}
\caption{Arquitectura típica de un sistema de extracción automática de terminología.}
\label{fig:arquitectura-ate}
\end{figure}


Las últimas tareas del flujo de trabajo son la validación de los candidatos a término y la evaluación del sistema de extracción diseñado (\cite{heylen2015automatic}). La primera atiende a la revisión manual por parte del terminólogo o del usuario para garantizar que los resultados son, efectivamente, términos, mientras que la segunda consiste en la evaluación del rendimiento del sistema utilizando métricas como la precisión, la exhaustividad y la medida F1, que permiten cuantificar la eficacia del sistema de extracción automática de términos en comparación con un estándar de referencia o \textit{Gold Standard}.

Un \textit{Gold Standard} o lista de referencia, en el contexto de la extracción automática de términos y otras áreas de la lingüística computacional y el procesamiento del lenguaje natural, se refiere a un conjunto de datos que se considera como la ‘verdad absoluta’ o la asumida como solución correcta contra la cual se evalúan los resultados de un sistema o modelo. Este conjunto está cuidadosamente anotado y revisado por expertos en la materia para garantizar su precisión. En el caso de la extracción automática de términos, el Gold Standard consiste generalmente en un corpus o dataset que ha sido anotado manualmente por terminólogos o especialistas en el dominio\footnote{  Otra opción de evaluación es tomar diccionarios y otros recursos especializados como referencia, aunque esto supondría renunciar al uso real y contextual de los términos. Teniendo en cuenta que estos recursos suelen ser limitados en cuanto a su cobertura terminológica, consideramos que pueden resultar una fuente de información a valorar, pero deberían ser complementados con los datos extraídos de corpus especializados para garantizar una mayor cobertura.
}, quienes identifican y etiquetan los términos específicos dentro de un texto. Estos términos anotados sirven como referencia para medir de manera objetiva el rendimiento de un extractor de términos. Para evaluar el rendimiento del sistema de extracción en contraste con un Gold Standard se utilizan las métricas mencionadas anteriormente, calculadas a partir de las distintas categorías de resultados que se muestran en la Tabla \ref{tab:clasificacion_resultados}:

\begin{table}[ht]
\centering
\fontsize{10}{12}\selectfont
\begin{tabular}{l p{7cm} l}
\toprule
\textbf{Categoría} & \textbf{Descripción} & \textbf{Ejemplo} \\
\midrule
Verdadero Positivo (VP) & Términos relevantes que han sido correctamente extraídos por el sistema. & \textit{corpus linguistics} \\
Falso Positivo (FP) & Unidades no terminológicas que han sido erróneamente extraídas como términos. & \textit{according to} \\
Verdadero Negativo (VN) & Unidades no terminológicas que el sistema no ha extraído, correctamente. & \textit{for example} \\
Falso Negativo (FN) & Términos relevantes que el sistema no ha logrado extraer. & \textit{dependency parsing} \\
\bottomrule
\end{tabular}
\caption{Ejemplo de clasificación de resultados en un sistema de extracción automática de términos.}
\label{tab:clasificacion_resultados}
\end{table}

La presencia de falsos positivos (FP) y falsos negativos (FN) afecta directamente a la calidad de un sistema de extracción automática de términos y están directamente relacionados con los conceptos de \textit{ruido} y \textit{silencio}:

\begin{itemize}
    \item El \textit{ruido} se genera cuando el sistema extrae propone unidades como candidatos a término, afectando negativamente a la precisión, ya que introduce información errónea dentro de los resultados, reduciendo la proporción de términos correctos entre los propuestos por el sistema (p. ej., si el sistema extrae como término ‘\textit{according to}’ en un corpus técnico, está generando ruido).
    \item El \textit{silencio} ocurre cuando el sistema no extrae términos relevantes que deberían haber sido identificados. Estos términos relevantes no son recuperados y se interpretan como falsos negativos (FN). El silencio impacta negativamente en la exhaustividad, recuperación o \textit{recall}, ya que reduce la capacidad del sistema para identificar todos los términos que realmente deberían haber sido extraídos (p. ej., si el término \textit{dependency parsing} aparece en el \textit{Gold Standard}, pero el sistema no lo extrae, se está produciendo silencio).
\end{itemize}

Estas métricas se definen mediante las fórmulas que detallamos a continuación:

\begin{itemize}
    \item La precisión mide la proporción de términos correctamente identificados (VP) sobre el total de unidades propuestas como candidatos a término (VP + FP). Refleja cuántos de los elementos que el sistema propone efectivamente son términos, y puede evaluarse de forma automática o mediante revisión experta.

    \equationcaption{Precisión}{
    \textit{Precisión} = \frac{VP}{VP + FP}
    }
    \label{eq:precision}


    \item La recuperación (también denominada \textit{exhaustividad} o \textit{recall}) evalúa la proporción de términos relevantes correctamente identificados (VP) sobre el total de términos que deberían haberse recuperado (VP + FN). Indica la capacidad del sistema para no omitir términos importantes del dominio.

    \equationcaption{\text{\textit{Recuperación}}}{
    \text{\textit{Recuperación}} = \frac{VP}{VP + FN}
    }
    \label{eq:recall}

    \item La F1-\textit{score} es la media armónica entre la precisión y la exhaustividad, y proporciona una medida única y equilibrada del rendimiento general del sistema de extracción. Un valor alto de F1 indica un buen compromiso entre ambas métricas.

    \equationcaption{\text{\textit{F1-score}}}{
    \text{\textit{F1}} = 2 \cdot \frac{\text{Precisión} \cdot \text{Exhaustividad}}{\text{Precisión} + \text{Exhaustividad}}
    }
    \label{eq:f1score}

\end{itemize}

La metodología de extracción terminológica adoptada en este trabajo se articula como un proceso incremental y multifase, orientado a la construcción progresiva de un corpus técnico anotado sin recurrir, en una fase inicial, a una anotación manual exhaustiva. Este enfoque concibe la extracción automática de términos como un procedimiento gradual de construcción y refinamiento del conocimiento terminológico, en el que se combinan de manera complementaria recursos de referencia, técnicas lingüísticas basadas en reglas, métodos estadísticos y validación experta.

En una primera fase se llevó a cabo una preanotación automática basada en recursos terminológicos de referencia, tales como obras lexicográficas, glosarios especializados y repertorios terminológicos existentes. Esta preanotación se realizó mediante coincidencia léxica exacta, tanto a nivel superficial como lematizado, con el objetivo de identificar de forma automática unidades terminológicas previamente documentadas. Como resultado de esta etapa se obtuvo una primera capa de anotación que puede caracterizarse como un silver standard inicial, necesariamente parcial y sesgado hacia el conocimiento terminológico ya establecido, pero útil como punto de partida para fases posteriores.

A continuación, se aplicaron técnicas de extracción terminológica no supervisada de carácter híbrido, combinando métodos lingüísticos basados en reglas con medidas estadísticas y probabilísticas ampliamente consolidadas en la literatura. En particular, se emplearon métricas como C-value y TF-IDF con el fin de identificar candidatos terminológicos no cubiertos por las listas de referencia. Esta fase permitió ampliar el espacio terminológico del corpus mediante la detección de unidades emergentes o altamente específicas del dominio, que no forman parte de los recursos iniciales.

El conjunto resultante de candidatos fue sometido posteriormente a un proceso de validación y revisión experta. En esta etapa se confirmaron las unidades consideradas terminológicamente relevantes, se descartaron falsos positivos y se integraron los nuevos términos aceptados en la anotación existente. Este procedimiento dio lugar a un silver standard enriquecido, que combina conocimiento terminológico previo con evidencia empírica derivada del análisis del corpus.

Sobre la base de este silver standard consolidado, se llevó a cabo una segunda capa de anotación automática informada, consistente en la proyección sistemática del conjunto terminológico validado sobre el corpus completo. Esta fase tuvo como objetivo incrementar la coherencia y homogeneidad de la anotación, asegurando una aplicación consistente de los criterios terminológicos previamente definidos.

En una fase posterior se realizó una tercera capa de anotación manual experta, mediante una revisión exhaustiva del corpus utilizando la herramienta Doccano. Esta intervención manual permitió corregir falsos positivos y falsos negativos, así como resolver casos de ambigüedad estructural, semántica o discursiva que no podían ser abordados de forma fiable mediante procedimientos automáticos. El resultado de esta etapa es un corpus técnico especializado anotado con un alto grado de control lingüístico y coherencia interna.

Finalmente, a partir de este corpus anotado y consolidado, se desarrolló la ingeniería de características y se entrenaron modelos de aprendizaje automático supervisado para la identificación automática de términos. En este diseño metodológico, el aprendizaje automático se sitúa como una fase posterior que se apoya explícitamente en las decisiones y validaciones previas, lo que permite garantizar tanto la trazabilidad del proceso como la interpretabilidad lingüística de los modelos resultantes.


Los métodos supervisados de aprendizaje automático requieren, como paso previo fundamental, un proceso de ingeniería de características que permita representar computacionalmente las propiedades lingüísticas y estadísticas de las unidades analizadas. Este proceso implica la selección, extracción y normalización de un conjunto de rasgos relevantes que puedan ser utilizados como entrada por los algoritmos de clasificación. En este trabajo, se ha diseñado un conjunto de rasgos que abarcan tanto dimensiones lingüísticas como estadísticas (Tablas \ref{tab:features_linguistic} y \ref{tab:features_statistical}), con el objetivo de capturar patrones característicos del comportamiento de los términos en el corpus. Estos rasgos actúan como variables explicativas en modelos supervisados que aprenden a distinguir automáticamente entre unidades terminológicas y no terminológicas a partir de ejemplos anotados. La selección y combinación adecuada de estos rasgos resulta determinante para el rendimiento y la interpretabilidad del modelo, y constituye una de las etapas más críticas en la aplicación de técnicas de aprendizaje automático al ámbito de la terminología computacional.


Para la identificación de unidades terminológicas el corpus anotado se ha convertido, para su procesamiento secuencial, al esquema BIO (\emph{Beginning, Inside, Outside}), ampliamente utilizado en tareas de secuenciación lingüística como el reconocimiento de entidades nombradas. Este esquema permite representar tanto unidades terminológicas poliléxicas como términos simples (formados por un único token), codificando de manera explícita la estructura secuencial de cada término.

\begin{itemize}
    \item \texttt{B-T} indica el inicio (\emph{beginning}) de una unidad terminológica. En el caso de términos simples, se marca únicamente el token con esta etiqueta.
    \item \texttt{I-T} señala un token interno (\emph{inside}) dentro de una unidad terminológica ya iniciada.
    \item \texttt{O} se asigna a cualquier token fuera (\emph{outside}) de una unidad terminológica.
\end{itemize}

Este formato es especialmente útil para modelar secuencias lingüísticas en algoritmos supervisados, como los modelos de regresión logística o las arquitecturas basadas en redes neuronales, ya que permite reconstruir con precisión las unidades terminológicas a partir de predicciones token por token:

\begin{table}[H]
\centering
\begin{tabular}{@{}lll@{}}
\toprule
\textbf{Token} & \textbf{Lema} & \textbf{Etiqueta BIO} \\
\midrule
Lemmatization   & lemmatization         & B-T \\
is              & be                    & O   \\
a               & a                     & O   \\
common          & common                & O   \\
step            & step                  & O   \\
in              & in                    & O   \\
a               & a                     & O   \\
parallel        & parallel              & B-T \\
corpus          & corpus                & I-T \\
.               & .                     & O   \\
\bottomrule
\end{tabular}
\caption{Ejemplo de anotación en formato BIO.}
\label{tab:bio_example}
\end{table}

Con el objetivo de construir un modelo supervisado capaz de identificar términos especializados en textos del dominio de la lingüística de corpus, se ha diseñado un conjunto de rasgos (\emph{features}) que capturan propiedades léxicas, sintácticas, estadísticas y contextuales de cada token. Estos rasgos permiten caracterizar de manera detallada la presencia y comportamiento de las unidades léxicas en distintos niveles de análisis. Para facilitar su interpretación y posterior análisis, los rasgos se han organizado en dos grandes categorías:

\begin{enumerate}
    \item Rasgos lingüísticos, que incluyen información gramatical, estructural y discursiva sobre las unidades léxicas, y
    \item Rasgos estadísticos, basados en propiedades frecuentistas y distribucionales.
\end{enumerate}
    
A continuación se describen todos los rasgos considerados en el corpus, acompañados, cuando proceda, de las fórmulas pertinentes para su cálculo.

\begin{enumerate}
  \item \emph{Rasgos lingüísticos}, que incluyen información gramatical, estructural y discursiva sobre las unidades léxicas.

  \begin{enumerate}
    \item \emph{Rasgos formales}: incluyen la forma superficial del token (\emph{token}), su forma lematizada (\emph{lemma}) y un valor binario que indica si se trata de una palabra vacía (\emph{is\_stop}). \textbf{<- esto podría ser FUNCIONAL}

    \item \emph{Rasgos sintácticos}: abarcan las categorías morfosintácticas binarizadas del token (\emph{is\_noun}, \emph{is\_verb}, \emph{is\_adj}, \emph{is\_propn}) y la probabilidad de aparición del lema con cada una de estas categorías:

    % Probabilidad condicional de POS dado un lema
    \begin{equation}
    P(\text{POS}_x \mid \text{lemma}) = \frac{\text{freq}(\text{lemma}, \text{POS}_x)}{\sum_{y} \text{freq}(\text{lemma}, \text{POS}_y)}
    \end{equation}

    También se incluyen \emph{has\_adjective\_modifier} y \emph{has\_pre\_modifier}, que indican si el token tiene un modificador adjetival o un modificador antepuesto, respectivamente.

    \item \emph{Rasgos discursivos}: se considera el rasgo binario \emph{definition\_pattern}, que señala si el token aparece en un patrón definitorio (por ejemplo, “\emph{X is defined as}”).

    \begin{table}[H]
\centering
\begin{tabular}{ll}
\toprule
\textbf{Patrón} & \textbf{Ejemplo de uso} \\
\midrule
\emph{is a} & A corpus \textbf{is a} collection of texts. \\
\emph{is defined as} & A token \textbf{is defined as} a basic unit of analysis. \\
\emph{refers to} & The term \emph{NLP} \textbf{refers to} Natural Language Processing. \\
\emph{known as} & This process is \textbf{known as} tokenization. \\
\emph{also known as} & POS tagging is \textbf{also known as} grammatical tagging. \\
\emph{is a part of} & A lemma \textbf{is a part of} a lexical entry. \\
\emph{is a type of} & A bigram \textbf{is a type of} n-gram. \\
\bottomrule
\end{tabular}
\caption{Patrones empleados para la detección de contextos definitorios.}
\label{tab:definition_patterns}
\end{table}

    
  \end{enumerate}

  \item \emph{Rasgos estadísticos}, basados en propiedades frecuentistas y distribucionales.

  \begin{enumerate}
    \item \emph{Frecuenciales}: se incluye la frecuencia del término en el texto actual (\emph{freq\_in\_text}):

    % Frecuencia del término en el documento
    \begin{equation}
    f_{t,d} = \text{número de ocurrencias de } t \text{ en } d
    \end{equation}

    así como su frecuencia global en el corpus (\emph{freq\_global}):

    % Frecuencia global del término en el corpus
    \begin{equation}
    f_t = \sum_{d \in \text{corpus}} f_{t,d}
    \end{equation}

    y su frecuencia de documento (\emph{document\_frequency}):

    % Frecuencia de documento
    \begin{equation}
    df_t = \left| \{ d \in \text{corpus} : t \in d \} \right|
    \end{equation}

    \item \emph{Distribucionales}: se considera una medida de dispersión (\emph{dispersion}), que puede aproximarse mediante el coeficiente de Juilland:

    % Dispersión de Juilland
    \begin{equation}
    D_J = 1 - \frac{\sqrt{\sum_{i=1}^{n} (f_i - \bar{f})^2 / n}}{\bar{f}}
    \end{equation}

    donde \( f_i \) es la frecuencia del término en el subcorpus \( i \), y \( \bar{f} \) es la frecuencia media. Asimismo, se incluye la diversidad contextual (\emph{contextual\_diversity}):

    % Diversidad contextual
    \begin{equation}
    CD(t) = \left| \{ c : t \in c \} \right|
    \end{equation}

    donde \( c \) es un contexto o documento distinto.
  \end{enumerate}

  \item \emph{Rasgos mixtos}, que combinan información posicional, de relevancia léxica y estructural.

  \begin{enumerate}
    \item \emph{Posicionales}: se considera la posición del token en la oración (\emph{position\_in\_sentence}) y la media de posiciones dentro de la oración (\emph{avg\_position\_in\_sentence}):

    % Posición media del token en la oración
    \begin{equation}
    \text{avg\_pos} = \frac{1}{n} \sum_{j=1}^{n} \text{posición}_j
    \end{equation}

    \item \emph{Relevancia léxica}: se calcula el valor \emph{tf\_idf} del término en el documento como:

    % Valor TF-IDF
    \begin{equation}
    \text{tf-idf}_{t,d} = \text{tf}_{t,d} \cdot \log\left(\frac{N}{df_t}\right)
    \end{equation}

    donde \( \text{tf}_{t,d} \) es la frecuencia relativa del término en el documento, \( df_t \) es su frecuencia de documento, y \( N \) el número total de documentos del corpus. También se considera la variedad léxica local (\emph{lexical\_variety}), que puede aproximarse mediante el índice tipo-token:

    % TTR (Type-Token Ratio)
    \begin{equation}
    \text{TTR} = \frac{V}{N}
    \end{equation}

    donde \( V \) es el número de tipos únicos y \( N \) el número total de tokens.

    \item \emph{Estructurales}: el rasgo \emph{is\_multiword\_term} indica si el token forma parte de una unidad léxica multipalabra.

    \item \emph{Identificadores}: finalmente, se incluyen \emph{doc\_id} y \emph{sentence\_id} para conservar la trazabilidad del token en su contexto de origen.
  \end{enumerate}
\end{enumerate}

\begin{table}[H]
\centering
\begingroup
\small
\begin{tabular}{@{}p{5cm}p{6cm}p{3cm}@{}}
\toprule
\multicolumn{3}{c}{\textbf{Rasgos lingüísticos en el modelado terminológico}} \\
\midrule
\textbf{Rasgo} & \textbf{Descripción} & \textbf{Tipo de dato} \\
\midrule
\multicolumn{3}{l}{\textbf{\textit{Morfosintáctico}s}} \\
\texttt{pos} & Categoría gramatical del token & Categórico \\
\texttt{tag} & Etiqueta morfosintáctica específica del token & Categórico \\
\texttt{dep} & Relación de dependencia sintáctica del token & Categórico \\
\texttt{is\_stop} & El token es una palabra funcional & Binario \\
\texttt{is\_punct} & El token es un signo de puntuación & Binario \\
\texttt{is\_alpha} & El token es alfabético & Binario \\
\texttt{is\_propn} & El token es un nombre propio & Binario \\

\addlinespace
\multicolumn{3}{l}{\textbf{\textit{Estructura y función sintáctica}}} \\
\texttt{is\_in\_noun\_phrase} & El token forma parte de un sintagma nomina. & Binario \\
\texttt{is\_in\_noun\_chunk\_head} & El token es cabeza de un sintagma nominal & Binario \\
\texttt{has\_adjective\_modifier} & Tiene modificadores adjetivales dependientes & Binario \\
\texttt{has\_pre\_modifier} & Tiene modificadores pre-nominales & Binario \\
\texttt{position\_in\_sentence} & Posición relativa del token dentro de la oración & Continuo \\

\addlinespace
\multicolumn{3}{l}{\textbf{\textit{Otros rasgos lingüísticos}}} \\
\texttt{is\_acronym} & El token es una sigla en mayúsculas & Binario \\
\texttt{shape} & Forma ortográfica general (Xxxx, xxxx, dddd, etc.) & Categórico \\
\texttt{in\_definition\_context} & El token aparece en un contexto definitorio ("is a", etc.) & Binario \\
\bottomrule
\end{tabular}
\endgroup
\caption{Rasgos lingüísticos utilizados para la detección automática de términos.}
\label{tab:features_linguistic}
\end{table}


\begin{table}[H]
\centering
\begingroup
\small
\begin{tabular}{@{}p{5cm}p{7cm}p{3cm}@{}}
\toprule
\multicolumn{3}{c}{\textbf{Rasgos estadísticos en el modelado terminológico}} \\
\midrule
\textbf{Rasgo} & \textbf{Descripción} & \textbf{Tipo de dato} \\
\midrule
\multicolumn{3}{l}{\textbf{\textit{Frecuencia y distribución interna}}} \\
\texttt{freq\_in\_text} & Frecuencia absoluta del token dentro del texto & Entero \\
\texttt{avg\_position\_sentence} & Posición media del token en las oraciones del documento & Continuo \\
\texttt{lexical\_variety} & Variedad léxica del documento en que aparece el token & Continuo \\
\texttt{dispersion} & Dispersión del término a lo largo del corpus & Continuo \\

\addlinespace
\multicolumn{3}{l}{\textbf{\textit{Frecuencia global y ponderación}}} \\
\texttt{freq\_global} & Frecuencia total del token en todo el corpus & Entero \\
\texttt{document\_frequency} & N.º de documentos en que aparece el token & Entero \\
\texttt{tf\_idf} & Frecuencia ponderada por la frecuencia inversa de documentos & Continuo \\

\addlinespace
\multicolumn{3}{l}{\textbf{\textit{Basados en algoritmos terminológicos}}} \\
\texttt{is\_in\_cvalue\_term} & El token forma parte de una unidad candidata identificada por \textit{C-Value} & Binario \\
\bottomrule
\end{tabular}
\endgroup
\caption{Rasgos estadísticos utilizados para la detección automática de términos.}
\label{tab:features_statistical}
\end{table}




a

\begin{table}[H]
\centering
\fontsize{10}{12}\selectfont
\begin{tabular}{lrr}
\toprule
\textbf{Tipo de rasgo} & \textbf{Cantidad} & \textbf{Porcentaje} \\
\midrule
Lingüísticos & 15 & 65.2\,\% \\
Estadísticos & 8  & 34.8\,\% \\
\bottomrule
\end{tabular}
\caption{Distribución de los tipos de rasgos utilizados en el modelo.}
\label{tab:tipos_rasgos}
\end{table}




\subsubsection{Método de aprendizaje profundo}

Un enfoque prompt-based...

- zero-shot
- few-shot

\subsection{Evaluación}

The two most popular ways to evaluate ATE performance are reference term lists
and manual validation (Pazienza et al. 2005; Astrakhantsev et al. 2015).

Las métricas básicas de evaluación son...

\subsubsection{\emph{Gold-Standard}}

Un Gold-Standard es...
\subsubsection{Juicio experto}

¿Pedir validación?

\section{Glosario digital}
\subsection{Diseño}
\subsection{Macro- y microestructura}

\clearpage



% ============================
% CAPÍTULO 5: RESULTADOS
% ============================

\phantomsection
\renewcommand{\tablename}{Tabla}
\setstretch{1.5}
\fontsize{14}{16}\selectfont

\chapter{Resultados}
\label{chap:5_resultados}

\makeatletter
\@afterindentfalse\@afterheading
\makeatother

\setcounter{section}{0}
\renewcommand{\thesection}{\thechapter.\arabic{section}}

\fontsize{12}{14}\selectfont
\vspace{\baselineskip}

\phantomsection
\setstretch{1.5}

En el presente capítulo se exponen los resultados obtenidos a partir de la metodología descrita en el Capítulo~\ref{chap:4_metodologia}. El capítulo se organiza en cinco secciones. La primera (§\ref{sec:extraccion}) describe el proceso de extracción terminológica híbrida: los inventarios validados para los subcorpus en inglés y español, el análisis de patrones morfosintácticos, las propiedades estadísticas de terminologicidad y una caracterización de la incrustación terminológica como propiedad estructural de los inventarios. La segunda (§\ref{sec:modelo_tc}) aborda el modelado supervisado a nivel de token y de candidato mediante cinco arquitecturas de clasificación, el análisis de importancia de rasgos, el protocolo de reserva de tipos terminológicos, la curva de saturación terminológica y los experimentos de tamaño y diversidad tipológica de corpus. La tercera (§\ref{sec:evaluacion_solapamiento}) caracteriza el sesgo estadístico en la evaluación de rasgos distribucionales y valida una corrección mediante recomputación por partición, con evaluación externa sobre el banco de datos ACTER. La cuarta (§\ref{sec:tecla}) describe la arquitectura y evaluación de TECLA, la librería que integra computacionalmente todos los componentes del modelo. La quinta (§\ref{sec:glosario}) presenta el glosario bilingüe resultante y las asimetrías terminológicas entre lenguas.

Estos resultados se articulan en el marco de \texttt{TECLA} (\textit{Terminology Extraction and Corpus Linguistics Application}), que implementa el flujo completo de procesamiento terminológico y materializa computacionalmente el modelo teórico propuesto. Más allá de los inventarios y métricas individuales, un hilo conductor atraviesa todos los resultados del capítulo: la \textbf{diversidad contextual} ---el número de contextos discursivos distintos en que aparece una unidad--- emerge sistemáticamente como el predictor primario de terminologicidad, tanto en los rasgos de los clasificadores supervisados como en la correlación con la detectabilidad de términos no vistos durante el entrenamiento ($\rho_s = 0{,}35$ EN; $\rho_s = 0{,}46$ ES; ambos $p < 0{,}001$), frente a la frecuencia bruta, cuya correlación con la detectabilidad no es significativa en el subcorpus inglés ($p = 0{,}14$). Este resultado constituye una corroboración empírica sólida de uno de los principios nucleares de la teoría terminológica contemporánea (§\ref{sec:evaluacion_solapamiento}): los términos se caracterizan no por su abundancia sino por la variedad y especificidad de sus contextos de uso.
% Chapter 4.1
\phantomsection
\setstretch{1.5}
\fontsize{12}{14}\selectfont

% ============================================================
%  SECCIÓN: Extracción terminológica
%  Requiere en el preámbulo:
%    \usepackage{booktabs, longtable, array, algorithm2e}
%    \usepackage[spanish,ruled,vlined]{algorithm2e}
%  Marcadores: % [FIGURA label]
% ============================================================

\section{Extracción terminológica}
\label{sec:extraccion}

La extracción terminológica se aplicó de forma independiente sobre los dos
subcorpus que integran el corpus de trabajo: el subcorpus en inglés, compuesto
por 50~documentos (artículos de revista, capítulos de libro y monografías)
sobre lingüística de corpus en el ámbito anglosajón, y el subcorpus en
español, también de 50~documentos de temática equivalente procedentes de
revistas y obras de referencia hispanohablantes. Esta separación responde a
dos motivaciones complementarias: por un lado, la especificidad de los
recursos lingüísticos requeridos por cada lengua ---modelos de
\textsc{pos-tagging}, \textit{stoplists} y patrones morfosintácticos
diferenciados---; por otro, el interés contrastivo de obtener un inventario
que refleje las particularidades del léxico especializado en cada tradición
académica.

El proceso se articula en cuatro fases secuenciales (véase la Figura~\ref{fig:diagrama_extraccion}):
(i)~preprocesamiento y limpieza del corpus; (ii)~filtrado lingüístico mediante
patrones morfosintácticos; (iii)~ponderación estadística de los candidatos;
y (iv)~validación experta y organización de cada inventario. En las secciones a continuación se describe cada una de ellas.

\begin{figure}[ht] % 'ht' ayuda a que se ubique 'aquí' o en la parte 'superior'
    \centering
    \begin{tikzpicture}[
        node distance=1.2cm,
        auto,
        block/.style={rectangle, draw, fill=blue!5, text width=9cm, text centered, rounded corners, minimum height=0.8cm, font=\small},
        subblock/.style={rectangle, draw, fill=gray!5, text width=4cm, text centered, rounded corners, minimum height=0.8cm, font=\footnotesize},
        filter/.style={rectangle, draw, dashed, fill=orange!5, text width=8.5cm, font=\scriptsize},
        arrow/.style={-Stealth, thick}
    ]

    % --- ENTRADA ---
    \node (corpus_en) [subblock] {\textbf{Subcorpus Inglés} \\ 50 documentos};
    \node (corpus_es) [subblock, right=0.8cm of corpus_en] {\textbf{Subcorpus Español} \\ 50 documentos};

    % --- FASE 1 ---
    \node (fase1) [block, below=1cm of $(corpus_en.south)!0.5!(corpus_es.south)$] {
        \textbf{I. Preprocesamiento y Anotación} \\
        \scriptsize Limpieza (\textsc{pdf}), \texttt{langdetect} y modelos \texttt{spaCy} (\textit{lg/news})
    };

    % --- FASE 2 ---
    \node (fase2) [block, below=0.8cm of fase1] {
        \textbf{II. Filtrado Lingüístico} \\
        \scriptsize Aplicación de Patrones POS (\textit{NP, de/of, Nominales})
    };

    % --- FILTROS (Algoritmo 1) ---
    \node (filtros) [filter, below=0.4cm of fase2] {
        \textbf{Filtros de Calidad:} \\
        Longitud ($1<n<6$) $\cdot$ Categorial (funcionales) $\cdot$ Normalización $\cdot$ \textit{Stoplists} $\cdot$ Freq $\ge 2$
    };

    % --- FASE 3 ---
    \node (fase3) [block, below=0.4cm of filtros] {
        \textbf{III. Ponderación Estadística} \\
        \scriptsize Índice: $0,6 \cdot \widetilde{\text{TF-IDF}} + 0,4 \cdot \widetilde{\text{C-value}}$ (+ bono \textit{corpus})
    };

    % --- FASE 4 ---
    \node (fase4) [block, below=0.8cm of fase3] {
        \textbf{IV. Validación Experta} \\
        \scriptsize Cribado por percentil 90 y categorización temática
    };

    % --- SALIDA ---
    \node (inv_en) [subblock, below=1cm of fase4, xshift=-2.3cm, fill=green!5] {
        \textbf{Inventario EN} \\ 170 términos validados
    };
    \node (inv_es) [subblock, below=1cm of fase4, xshift=2.3cm, fill=green!5] {
        \textbf{Inventario ES} \\ 148 términos validados
    };

    % --- LÍNEAS ---
    \draw [arrow] (corpus_en.south) -- ++(0,-0.4) -| (fase1.north);
    \draw [arrow] (corpus_es.south) -- ++(0,-0.4) -| (fase1.north);
    \draw [arrow] (fase1) -- (fase2);
    \draw [arrow] (fase2) -- (filtros);
    \draw [arrow] (filtros) -- (fase3);
    \draw [arrow] (fase3) -- (fase4);
    \draw [arrow] (fase4.south) -- ++(0,-0.4) -| (inv_en.north);
    \draw [arrow] (fase4.south) -- ++(0,-0.4) -| (inv_es.north);

    \end{tikzpicture}
    
    \caption{Flujo secuencial de las fases de extracción terminológica.}
    \label{fig:diagrama_extraccion}
\end{figure}

Antes de la extracción propiamente dicha, los textos se sometieron a una
limpieza estructural orientada a eliminar el ruido tipográfico y documental
habitual en textos digitalizados o exportados desde \textsc{pdf}. Las
operaciones aplicadas incluyeron la supresión de las secciones de bibliografía
y referencias, la eliminación de pies de figura y tabla, la limpieza de
encabezados de sección aislados, la normalización de guiones de final de línea
(palabras partidas) y el descarte de líneas compuestas mayoritariamente por
cifras o marcadores numéricos de cita.

La detección automática de la lengua se efectuó sobre los primeros
3\,000~caracteres del texto limpio mediante la librería
\texttt{langdetect}\footnote{\url{https://pypi.org/project/langdetect/}} y
determinó el modelo de \texttt{spaCy} cargado en cada caso:
\texttt{en\_core\_web\_lg} para los documentos en inglés y
\texttt{es\_core\_news\_md} para los documentos en español. La anotación
morfosintáctica resultante (tokens, categorías gramaticales \textsc{pos} y
lemas) constituye la entrada de la fase siguiente. La Figura~\ref{fig:viz_spacy} ilustra el
resultado de esta anotación sobre dos fragmentos representativos del corpus: uno del
subcorpus en inglés y otro del subcorpus en español. Cada token se muestra con su
etiqueta \textsc{pos} (estándar Universal Dependencies); los fondos sombreados y los
marcadores identifican los tokens que forman unidades terminológicas validadas
(\texttt{B-T}: inicio de término; \texttt{I-T}: continuación).

\begin{figure}[H]
    \centering
    \includegraphics[width=\linewidth]{images/viz_spacy_pos.png}
    \caption{Anotación morfosintáctica de dos fragmentos del corpus TECLA mediante
             \texttt{spaCy} (Universal Dependencies). Panel superior: subcorpus EN
             (\textit{In corpus linguistics, parallel corpora are usually smaller\ldots});
             panel inferior: subcorpus ES
             (\textit{En la lematización y el etiquetado morfosintáctico de las
             palabras del corpus\ldots}). Los fondos en verde claro señalan tokens
             que forman parte de una unidad terminológica validada (\texttt{B-T}/\texttt{I-T}).
             La variedad de etiquetas \textsc{pos} en una misma oración muestra la
             riqueza del contexto sintáctico explotado por los patrones de extracción.}
    \label{fig:viz_spacy}
\end{figure}

La identificación de candidatos a término se basa en un conjunto de patrones
morfosintácticos formalizados sobre etiquetas \textsc{pos} y definidos
diferenciadamente por lengua. Los patrones recogen las estructuras
sintagmáticas más productivas en la formación del léxico especializado:
sintagmas nominales pre- y posnominales, secuencias de sustantivos y
construcciones mediadas por preposición (\textit{of}~/~\textit{de}). La
Tabla~\ref{tab:patrones-pos} sintetiza la gramática terminológica empleada
junto con ejemplos representativos.

\begin{table}[H]
\centering
\small
\caption{Gramática terminológica: patrones \textsc{pos} de extracción
         por lengua, con ejemplos representativos.}
\label{tab:patrones-pos}
\begin{tabular}{@{} c p{4.2cm} p{4.2cm} p{4.5cm} @{}}
\toprule
 & \textbf{ID} & \textbf{Patrón POS} & \textbf{Ejemplos} \\
\midrule
\multicolumn{4}{@{}l}{\textbf{Inglés (\texttt{en})}} \\
\midrule
1 & \texttt{EN\_NP\_PRENOMINAL}  & \texttt{(ADJ|NOUN|PROPN)* NOUN} & \textit{balanced corpus}, \textit{annotation scheme} \\
2 & \texttt{EN\_NP\_OF}          & \texttt{NOUN+ of NOUN+}          & \textit{part of speech}, \textit{analysis of variance} \\
3 & \texttt{EN\_NOUN\_SEQUENCE}  & \texttt{NOUN NOUN+}              & \textit{token frequency}, \textit{corpus benchmark} \\
4 & \texttt{EN\_NP\_PROPN}       & \texttt{PROPN PROPN+}            & \textit{British National Corpus} \\
\midrule
\multicolumn{4}{@{}l}{\textbf{Español (\texttt{es})}} \\
\midrule
5  & \texttt{ES\_NP\_PRENOMINAL}  & \texttt{(ADJ|NOUN|PROPN)* NOUN}  & \textit{patrón colocacional} \\
6  & \texttt{ES\_NP\_POSTNOMINAL} & \texttt{NOUN ADJ+}               & \textit{corpus balanceado}, \textit{frecuencia léxica} \\
7  & \texttt{ES\_NP\_DE}          & \texttt{NOUN+ de NOUN+}          & \textit{esquema de anotación}, \textit{lista de palabras} \\
8  & \texttt{ES\_NP\_DE\_ADJ}     & \texttt{NOUN+ de NOUN ADJ+}      & \textit{análisis de varianza estadística} \\
9  & \texttt{ES\_NP\_PROPN}       & \texttt{PROPN PROPN+}            & \textit{CREA Académico} \\
10 & \texttt{ES\_NP\_NOUN\_PROPN} & \texttt{NOUN PROPN+}             & \textit{corpus PRESEEA} \\
\bottomrule
\end{tabular}
\end{table}

Mientras que la Tabla~\ref{tab:patrones-pos} describe la gramática
formalizada \emph{ex ante}, la Tabla~\ref{tab:patrones_pos} muestra el
comportamiento empírico de dichos patrones tras la validación experta:
recoge la distribución observada de las estructuras morfosintácticas en el
inventario final del corpus.

\begin{table}[H]
\centering
\small
\setlength{\tabcolsep}{4pt}
\renewcommand{\arraystretch}{1.15}
\caption{Distribución empírica de patrones morfosintácticos en los términos
  validados del corpus \textsc{tecla} (EN: $n=155$ tipos; ES: $n=146$ tipos).
  Porcentajes calculados sobre el total de tipos únicos por lengua.}
\label{tab:patrones_pos}
\begin{tabular}{@{} l r r @{\quad} l r r @{}}
\toprule
\multicolumn{3}{c}{\textbf{Inglés (EN)}} & \multicolumn{3}{c}{\textbf{Español (ES)}} \\
\cmidrule(lr){1-3}\cmidrule(lr){4-6}
\textbf{Patrón} & \textbf{$n$} & \textbf{\%} &
\textbf{Patrón} & \textbf{$n$} & \textbf{\%} \\
\midrule
Uni-token (N/PROPN)      & 63 & 40{,}6 & Uni-token (N/PROPN)   & 95 & 65{,}1 \\
(ADJ|N|PROPN)* N         & 30 & 19{,}4 & N ADJ$^{+}$           & 18 & 12{,}3 \\
N N$^{+}$                & 27 & 17{,}4 & PROPN PROPN$^{+}$     & 16 & 11{,}0 \\
PROPN PROPN$^{+}$        & 16 & 10{,}3 & N$^{+}$ de N$^{+}$    &  6 &  4{,}1 \\
Compuesto con guion      &  8 &  5{,}2 & (ADJ|N|PROPN)* N      &  5 &  3{,}4 \\
Otros                    &  7 &  4{,}5 & Otros                 &  3 &  2{,}1 \\
N$^{+}$ of N$^{+}$       &  4 &  2{,}6 & N PROPN$^{+}$         &  2 &  1{,}4 \\
                         &    &        & N$^{+}$ de N ADJ$^{+}$&  1 &  0{,}7 \\
\midrule
\textbf{Total}           & \textbf{155} & \textbf{100} &
\textbf{Total}           & \textbf{146} & \textbf{100} \\
\bottomrule
\end{tabular}
\end{table}

Los datos confirman varias regularidades del léxico especializado del dominio.
En primer lugar, los uni-tokens dominan en ambas lenguas, con una asimetría
notable a favor del español (65,1\,\% frente al 40,6\,\% del inglés): la
mayor productividad morfológica del español permite lexicalizar mediante
derivación nominal (\textit{corpus}, \textit{anotación}, \textit{concordancia})
conceptos que en inglés requieren composición sintagmática. En segundo
lugar, las construcciones multipalabra revelan el contraste tipológico
esperado entre ambas lenguas: el inglés se apoya preferentemente en el
patrón prenominal \texttt{(ADJ|N|PROPN)* N} (19,4\,\%; \textit{natural
language processing}, \textit{inter-annotator agreement}), mientras que el
español privilegia el patrón posnominal \texttt{N ADJ}$^{+}$ (12,3\,\%;
\textit{lingüística computacional}, \textit{frecuencia léxica}). Por último,
las secuencias de nombres propios mantienen un peso comparable en ambas
lenguas (10-11\,\%) y corresponden, casi sin excepción, a denominaciones de
recursos y herramientas (\textit{British National Corpus},
\textit{WordSmith Tools}).

El procedimiento de extracción sigue el Algoritmo~\ref{alg:extraccion}: para
cada oración del corpus anotado se comprueba si alguna secuencia de tokens
satisface alguno de los patrones definidos para la lengua del documento;
en caso afirmativo, el candidato se somete a una batería de filtros de
calidad y, si los supera, se normaliza canónicamente y se incorpora al
inventario.

\begin{algorithm}[H]
\caption{Extracción de candidatos a término.}
\label{alg:extraccion}
\KwEntrada{Corpus $\mathcal{C} = \{d_1, \dots, d_n\}$, patrones POS $\mathcal{G}$,
           \textit{stoplists} por lengua $\mathcal{S}_\ell$,
           longitud máxima \texttt{max\_tokens}\,=\,6.}
\KwSalida{Lista de candidatos normalizados $\mathcal{T}$.}
\ParaCada{documento $d$ de $\mathcal{C}$}{
    Detectar la lengua del documento $\rightarrow \ell$\;
    Cargar el modelo \texttt{spaCy} para $\ell$ y anotar $d$
    (tokens, etiquetas \textsc{pos}, lemas) $\rightarrow \mathcal{A}$\;
    \ParaCada{oración $s$ de $\mathcal{A}$}{
        \ParaCada{patrón $p$ de $\mathcal{G}[\ell]$}{
            \ParaCada{secuencia $\sigma$ que coincide con $p$ en $s$}{
                \lSi{$|\sigma|=0$ \textbf{o} $|\sigma|>\texttt{max\_tokens}$}{\textbf{descartar}}
                \lSi{$\sigma$ contiene determinantes, pronombres u otras palabras funcionales}{\textbf{descartar}}
                Construir la forma canónica $\lambda$ concatenando los lemas en minúscula\;
                \lSi{forma superficial de $\sigma$ o $\lambda \in \mathcal{S}_\ell$}{\textbf{descartar}}
                \lSi{$|\sigma|=1$ \textbf{y} $\lambda$ es \textit{stopword}}{\textbf{descartar}}
                Añadir $(\sigma,\lambda,d)$ a $\mathcal{T}$\;
            }
        }
    }
}
Consolidar $\mathcal{T}$: acumular frecuencias absolutas y documentales\;
\lSi{frecuencia$(\lambda) < 2$ \textbf{o} doc\_freq$(\lambda) < 1$}{\textbf{descartar}}
\Devolver{$\mathcal{T}$ ordenada por forma canónica.}
\end{algorithm}

La cadena de filtros responde a cinco criterios complementarios. El
\textit{filtro de longitud} descarta las secuencias vacías y las que
superan los seis tokens. El \textit{filtro categorial} elimina los
candidatos que contienen determinantes, pronombres u otras palabras
funcionales. La \textit{normalización canónica} reduce las secuencias
supervivientes a su forma lematizada en minúsculas. El \textit{filtro de
\textit{stoplist}} rechaza el candidato cuando su forma superficial o
canónica coincide con alguna entrada de la lista de exclusión académica
por lengua ---construida a partir de los recursos de \texttt{spaCy} y
\texttt{NLTK} y complementada con los ficheros específicos
\texttt{stop\_en}/\texttt{stop\_es}---; en el caso de los unigramas, el
rechazo se aplica también cuando el lema coincide con una \textit{stopword}.
Por último, el \textit{filtro de frecuencia mínima} ---aplicado en la fase
de consolidación--- descarta los candidatos con frecuencia absoluta inferior
a~2 o presentes en un único documento, por considerarse estadísticamente
insuficientes para su validación.

El análisis estructural de los candidatos supervivientes revela una
regularidad marcada: en ambos subcorpus, la mayoría corresponde a
secuencias formadas por sustantivos consecutivos, seguidas a distancia
por las construcciones preposicionales y los sintagmas adjetivo-sustantivo.
Este predominio de la composición nominal confirma empíricamente que la
terminología de la lingüística de corpus se lexicaliza preferentemente
mediante compuestos nominales, en consonancia con lo descrito en la
bibliografía sobre terminología científica y técnica
\autocite{cabre1999terminology, sager1990terminology}

Una vez delimitado el conjunto de candidatos por criterios lingüísticos, se
aplicaron dos métricas complementarias para estimar el grado de
\emph{terminologicidad} de cada unidad: \textit{TF-IDF} y \textit{C-value}
(véase~§\ref{subsec:ate_metodos}). Ambas se reescalaron al intervalo $[0,1]$
mediante normalización mínimo-máximo y se combinaron en un índice
compuesto mediante ponderación lineal:

\begin{equation}
    \text{terminologicidad}(t) \;=\; 0{,}6 \cdot \widetilde{\text{TF-IDF}}(t)
                                \;+\; 0{,}4 \cdot \widetilde{\text{C-value}}(t).
    \label{eq:terminologicidad}
\end{equation}

A este índice se incorpora una bonificación cualitativa para los candidatos
que contienen explícitamente el núcleo léxico \textit{corpus} ---o su
equivalente en español---, en reconocimiento de su papel estructurante en
el dominio y sin alterar la ponderación relativa de las métricas. La
justificación empírica de esta decisión puede leerse a la luz del análisis de
incrustación terminológica presentado en §\ref{subsec:terminos_incrustados}: el
término \textit{corpus} actúa como núcleo de 28~cadenas terminológicas en el
subcorpus inglés y de 26~en el español, con tasas de incrustación del 32\,\%
y el 13\,\% respectivamente; su frecuencia bruta sobrestima sistemáticamente su
terminologicidad autónoma, y la bonificación compensa parcialmente el descuento
que el C-value aplica a sus ocurrencias incrustadas. El
índice resultante permite ordenar los candidatos a lo largo de un
\textit{continuum} de especialización y aplicar umbrales de corte
---por ejemplo, el percentil~90--- para concentrar la validación experta
en aquellos con mayor probabilidad de constituir términos válidos. Los
resultados muestran un comportamiento diferenciado y complementario de
ambas métricas: el \textit{C-value} alcanza valores más altos en las
primeras posiciones del \textit{ranking} para unidades poliléxicas con
fuerte cohesión sintagmática, mientras que el \textit{TF-IDF} discrimina
mejor a lo largo del \textit{continuum} al penalizar las unidades con
distribución demasiado difusa.

Los candidatos situados en el percentil superior del índice de
terminologicidad fueron sometidos a validación experta con un doble
objetivo: confirmar su pertinencia como unidades terminológicas del
dominio y asignarles una categoría temática. El criterio de validación
combinó los principios formales descritos en la metodología (véase
Capítulo~\ref{chap:3_extraccion}) con el conocimiento experto de la
disciplina, prestando especial atención a la estabilidad denominativa,
la autonomía referencial y la pertenencia al léxico nuclear de la
lingüística de corpus.

El proceso arrojó dos inventarios validados: 170~términos en el subcorpus
inglés, distribuidos en 11~categorías temáticas, y 148~términos en el
subcorpus español, agrupados en 6~macrocategorías. El mayor peso
proporcional de los unigramas en el subcorpus español (cf.
Tabla~\ref{tab:patrones_pos}) puede explicarse, como ya se anticipó, por
la productividad morfológica de la lengua, que permite lexicalizar
conceptos especializados mediante derivación nominal sin necesidad de
recurrir sistemáticamente a compuestos sintagmáticos. Las
Figuras~\ref{fig:categorias_en} y~\ref{fig:categorias_es} ilustran la
distribución por categorías en cada subcorpus.

\begin{figure}[H]
    \centering
    \includegraphics[width=0.8\linewidth]{images/fig_categorias_en.png}
    \caption{Distribución de los términos validados del subcorpus en
             inglés por categoría temática.}
    \label{fig:categorias_en}
\end{figure}

\begin{figure}[H]
    \centering
    \includegraphics[width=0.8\linewidth]{images/fig_categorias_en.png}
    \caption{Distribución de los términos validados del subcorpus en
             español por categoría temática (equivalente a EN).}
    \label{fig:categorias_es}
\end{figure}

En el subcorpus inglés predominan los términos vinculados al
\textit{diseño y compilación de corpus} (34~unidades), seguidos por
\textit{frecuencia y estadística} (25) y \textit{léxico y fraseología}
(25). En el subcorpus español, la categoría \textit{recursos y corpus}
(\textsc{rec-cor}) concentra el mayor número de unidades (52), seguida
por \textit{herramientas y software} (\textsc{her-soft}, 41~términos) y
\textit{metodología y diseño} (\textsc{met-di}, 26~términos). Esta última
es directamente comparable con la categoría \textit{diseño y compilación
de corpus} del subcorpus inglés, lo que indica una cobertura temática
equivalente para el núcleo metodológico del dominio. La diferencia más
notable reside en el peso de \textsc{rec-cor} en español, que excede
considerablemente la proporción de su equivalente inglés ---hecho
coherente con la mayor atención que la bibliografía hispanohablante
dedica a la descripción y catalogación de corpus disponibles y
herramientas de procesamiento, frecuentemente articulada como revisión
historiográfica de proyectos nacionales (\textit{CREA}, \textit{CORDE},
\textit{CORPES~XXI}, \textit{PRESEEA}).

\begin{figure}[H]
    \centering
    \includegraphics[width=\linewidth]{images/fig_categorias_comparativa.png}
    \caption{Comparación normalizada de la distribución por macrocategorías
             entre el subcorpus en inglés y el subcorpus en español.}
    \label{fig:categorias_comparativa}
\end{figure}

\subsubsection{Términos vs. Entidades: Una distinción tipológica}
\label{subsec:terminos_vs_entidades}

El análisis del inventario de unidades terminológicas validadas pone de manifiesto una heterogeneidad tipológica que trasciende la distinción convencional entre términos y no-términos. Aunque el conjunto total de 318~unidades (170~en inglés, 148~en español) comparte un valor especializado indiscutible y potencial lexicográfico como material de base para la construcción de glosarios bilingües, no todas ellas constituyen términos en sentido estricto.

La distinción que proponemos en las tablas siguientes separa dos categorías de unidades complementarias:

\begin{enumerate}
    \item \textbf{Términos:} unidades léxicas que designan conceptos, procesos, propiedades o metodologías especializadas del dominio (ej.\ \textit{representatividad}, \textit{anotación semántica}, \textit{frecuencia relativa}). Son las unidades que responden a la caracterización formal, estadística y distribucional desarrollada en la sección anterior.

    \item \textbf{Entidades:} denominaciones propias de recursos, proyectos, herramientas y referencias concretas (ej.\ \textit{British National Corpus}, \textit{Sketch Engine}, \textit{CREA}). Poseen un marcado valor léxico-referencial para la construcción de glosarios, pero designan particulares institucionales o comerciales antes que conceptos generalizables.
\end{enumerate}

Esta separación tipológica revela un contraste lingüístico y académico profundo entre las dos tradiciones bibliográficas (Tabla~\ref{tab:terminos_vs_entidades}): el subcorpus inglés concentra su terminología en conceptos especializados (84{,}7\,\% términos, 15{,}3\,\% entidades), reflejando una tradición de investigación teórica consolidada y orientada hacia la formalización de conceptos; el subcorpus español, en cambio, dedica una proporción sustancialmente mayor de su vocabulario a la documentación de recursos, proyectos y herramientas de origen local (37{,}2\,\% términos, 62{,}8\,\% entidades), evidenciando un énfasis historiográfico y un mayor detalle en la descripción y catalogación de infraestructuras nacionales.

\begin{table}[H]
\centering
\small
\caption{Distribución de unidades terminológicas por tipo lingüístico.
         \textit{Unidades}: total de tipos validados.
         \textit{Términos}: unidades con valor conceptual especializado (stricto sensu).
         \textit{Entidades}: denominaciones de recursos, herramientas y proyectos con valor lexicográfico.}
\label{tab:terminos_vs_entidades}
\begin{tabular}{@{} l r r r r r r @{}}
\toprule
\textbf{Lengua} & \textbf{Unidades totales} & \multicolumn{2}{c}{\textbf{Términos}} & \multicolumn{2}{c}{\textbf{Entidades}} \\
\cmidrule(lr){3-4}\cmidrule(lr){5-6}
 & & $n$ & \% & $n$ & \% \\
\midrule
Inglés (EN) & 170 & 144 & 84{,}7\,\% & 26 & 15{,}3\,\% \\
Español (ES) & 148 & 55 & 37{,}2\,\% & 93 & 62{,}8\,\% \\
\midrule
\textbf{Total} & \textbf{318} & \textbf{199} & \textbf{62{,}6\,\%} & \textbf{119} & \textbf{37{,}4\,\%} \\
\bottomrule
\end{tabular}
\end{table}

Las Tablas~\ref{tab:terminos_en} y~\ref{tab:terminos_es} recogen el
inventario completo de términos validados para cada subcorpus junto con
sus valores estadísticos asociados. El inventario inglés
(Tabla~\ref{tab:terminos_en}) evidencia la madurez de la disciplina en
esta lengua a través de una terminología altamente estratificada: las
unidades con mayor índice de terminologicidad no son únicamente las más
frecuentes (\textit{corpus}, \textit{annotation}), sino aquellas que
operan como nodos conceptuales de los que derivan subdominios enteros.
El inventario captura simultáneamente la terminología técnica del
procesamiento (\textit{POS tagging}, \textit{lemmatisation}) y la de
base estadística (\textit{log-likelihood}, \textit{confidence interval}),
lo que demuestra que el corpus de trabajo cubre satisfactoriamente las
vertientes teórica y aplicada de la lingüística de corpus anglosajona.

\small
\begin{longtable}{@{} r p{5cm} r r r r @{}}
\toprule
\# & \textbf{UT} & \textbf{Frec.\ abs.} & \textbf{DF} & \textbf{TF--IDF} & \textbf{C-value} \\
\midrule
\endfirsthead
\multicolumn{6}{l}{\small\itshape (continuación)} \\
\toprule
\# & \textbf{UT} & \textbf{Frec.\ abs.} & \textbf{DF} & \textbf{TF--IDF} & \textbf{C-value} \\
\midrule
\endhead
\midrule
\multicolumn{6}{r}{\small\itshape Continúa en la página siguiente\ldots} \\
\endfoot
\bottomrule
\caption{Términos validados del subcorpus en inglés, organizados por
         categoría temática.}
\label{tab:terminos_en}
\endlastfoot
\multicolumn{6}{l}{\textbf{Conceptos fundacionales}} \\
\midrule
1. & \textit{corpus} & 13767 & 50 & 1{,}0000 & 1{,}0000 \\
2. & \textit{type} & 1048 & 41 & 0{,}0974 & 0{,}1888 \\
3. & \textit{corpus linguistics} & 2138 & 47 & 0{,}1689 & 0{,}1558 \\
4. & \textit{corpus-based} & 416 & 38 & 0{,}0419 & 0{,}0397 \\
5. & \textit{corpus-driven} & 31 & 9 & 0{,}0075 & 0{,}0030 \\
\midrule
\multicolumn{6}{l}{\textbf{Diseño y compilación de corpus}} \\
\midrule
6. & \textit{representativeness} & 555 & 22 & 0{,}0866 & 0{,}0415 \\
7. & \textit{speech corpus} & 151 & 6 & 0{,}0424 & 0{,}0310 \\
8. & \textit{text type} & 266 & 24 & 0{,}0392 & 0{,}0486 \\
9. & \textit{corpus design} & 319 & 20 & 0{,}0528 & 0{,}0426 \\
10. & \textit{genre} & 297 & 33 & 0{,}0342 & 0{,}0396 \\
11. & \textit{sampling frame} & 91 & 5 & 0{,}0270 & 0{,}0158 \\
12. & \textit{corpus size} & 104 & 13 & 0{,}0216 & 0{,}0155 \\
13. & \textit{reference corpus} & 113 & 20 & 0{,}0187 & 0{,}0190 \\
14. & \textit{national corpus} & 189 & 26 & 0{,}0263 & 0{,}0185 \\
15. & \textit{general corpus} & 85 & 7 & 0{,}0227 & 0{,}0106 \\
16. & \textit{corpus building} & 93 & 7 & 0{,}0248 & 0{,}0095 \\
17. & \textit{multimodal corpus} & 33 & 3 & 0{,}0112 & 0{,}0062 \\
18. & \textit{learner corpus} & 64 & 9 & 0{,}0156 & 0{,}0081 \\
19. & \textit{monitor corpus} & 42 & 7 & 0{,}0112 & 0{,}0067 \\
20. & \textit{training corpus} & 31 & 6 & 0{,}0087 & 0{,}0056 \\
21. & \textit{subcorpus} & 84 & 16 & 0{,}0158 & 0{,}0048 \\
22. & \textit{web corpus} & 32 & 6 & 0{,}0090 & 0{,}0045 \\
23. & \textit{comparable corpus} & 38 & 7 & 0{,}0101 & 0{,}0063 \\
24. & \textit{annotated corpus} & 70 & 10 & 0{,}0163 & 0{,}0060 \\
25. & \textit{multilingual corpus} & 41 & 8 & 0{,}0104 & 0{,}0050 \\
26. & \textit{corpus compilation} & 43 & 15 & 0{,}0083 & 0{,}0058 \\
27. & \textit{spoken corpus} & 63 & 15 & 0{,}0122 & 0{,}0055 \\
28. & \textit{monolingual corpus} & 30 & 3 & 0{,}0102 & 0{,}0031 \\
29. & \textit{diachronic corpus} & 29 & 11 & 0{,}0065 & 0{,}0047 \\
30. & \textit{mega corpora} & 39 & 3 & 0{,}0132 & 0{,}0025 \\
31. & \textit{historical corpus} & 19 & 10 & 0{,}0044 & 0{,}0042 \\
32. & \textit{specialized corpus} & 27 & 12 & 0{,}0058 & 0{,}0033 \\
33. & \textit{balanced corpus} & 31 & 10 & 0{,}0072 & 0{,}0030 \\
34. & \textit{web crawling} & 5 & 3 & 0{,}0017 & 0{,}0006 \\
35. & \textit{synchronic} & 94 & 16 & 0{,}0176 & 0{,}0090 \\
36. & \textit{parallel corpus} & 145 & 8 & 0{,}0369 & 0{,}0138 \\
37. & \textit{written corpus} & 15 & 9 & 0{,}0037 & 0{,}0014 \\
38. & \textit{sample size} & 217 & 9 & 0{,}0528 & 0{,}0207 \\
39. & \textit{stratified sampling} & 5 & 1 & 0{,}0021 & 0{,}0001 \\
\midrule
\multicolumn{6}{l}{\textbf{Frecuencia y estadística}} \\
\midrule
40. & \textit{frequency} & 1667 & 40 & 0{,}1592 & 0{,}1822 \\
41. & \textit{sample} & 1127 & 31 & 0{,}1369 & 0{,}1553 \\
42. & \textit{probability} & 380 & 17 & 0{,}0691 & 0{,}0438 \\
43. & \textit{range} & 584 & 37 & 0{,}0604 & 0{,}0539 \\
44. & \textit{dispersion} & 329 & 16 & 0{,}0618 & 0{,}0312 \\
45. & \textit{occurrence} & 308 & 24 & 0{,}0454 & 0{,}0423 \\
46. & \textit{population} & 286 & 16 & 0{,}0537 & 0{,}0287 \\
47. & \textit{association measures} & 102 & 6 & 0{,}0286 & 0{,}0146 \\
48. & \textit{descriptive statistics} & 55 & 8 & 0{,}0140 & 0{,}0047 \\
49. & \textit{relative frequency} & 28 & 9 & 0{,}0068 & 0{,}0052 \\
50. & \textit{raw frequency} & 10 & 3 & 0{,}0034 & 0{,}0033 \\
51. & \textit{mutual information} & 31 & 9 & 0{,}0075 & 0{,}0031 \\
52. & \textit{inferential statistics} & 22 & 6 & 0{,}0062 & 0{,}0017 \\
53. & \textit{absolute frequency} & 5 & 2 & 0{,}0018 & 0{,}0014 \\
54. & \textit{chi-squared} & 46 & 8 & 0{,}0117 & 0{,}0003 \\
55. & \textit{log-likelihood} & 39 & 12 & 0{,}0084 & 0{,}0037 \\
56. & \textit{standard deviation} & 121 & 14 & 0{,}0243 & 0{,}0115 \\
57. & \textit{statistical significance} & 34 & 10 & 0{,}0079 & 0{,}0032 \\
58. & \textit{confidence interval} & 46 & 6 & 0{,}0129 & 0{,}0044 \\
59. & \textit{t-score} & 15 & 3 & 0{,}0051 & 0{,}0014 \\
60. & \textit{mean score} & 77 & 2 & 0{,}0285 & 0{,}0073 \\
61. & \textit{z-score} & 1 & 1 & 0{,}0004 & 0{,}0001 \\
62. & \textit{deviation of proportions} & 2 & 2 & 0{,}0007 & 0{,}0003 \\
63. & \textit{normalized frequency} & 10 & 5 & 0{,}0030 & 0{,}0000 \\
64. & \textit{effect size} & 60 & 8 & 0{,}0153 & 0{,}0002 \\
\midrule
\multicolumn{6}{l}{\textbf{Anotación y marcado}} \\
\midrule
65. & \textit{annotation} & 1513 & 38 & 0{,}1524 & 0{,}1450 \\
66. & \textit{annotation scheme} & 126 & 16 & 0{,}0237 & 0{,}0230 \\
67. & \textit{tagset} & 155 & 13 & 0{,}0323 & 0{,}0161 \\
68. & \textit{word class} & 37 & 8 & 0{,}0094 & 0{,}0168 \\
69. & \textit{accuracy} & 165 & 29 & 0{,}0212 & 0{,}0150 \\
70. & \textit{corpus annotation} & 132 & 18 & 0{,}0232 & 0{,}0136 \\
71. & \textit{recall} & 107 & 17 & 0{,}0194 & 0{,}0053 \\
72. & \textit{POS tagging} & 72 & 11 & 0{,}0161 & 0{,}0085 \\
73. & \textit{treebank} & 105 & 13 & 0{,}0219 & 0{,}0039 \\
74. & \textit{parts of speech} & 75 & 16 & 0{,}0141 & 0{,}0099 \\
75. & \textit{annotation tools} & 51 & 8 & 0{,}0130 & 0{,}0049 \\
76. & \textit{markup language} & 40 & 9 & 0{,}0097 & 0{,}0044 \\
77. & \textit{syntactic annotation} & 32 & 7 & 0{,}0085 & 0{,}0028 \\
78. & \textit{semantic annotation} & 34 & 12 & 0{,}0073 & 0{,}0029 \\
79. & \textit{gold standard} & 13 & 5 & 0{,}0039 & 0{,}0012 \\
80. & \textit{morphological annotation} & 7 & 2 & 0{,}0026 & 0{,}0006 \\
81. & \textit{part-of-speech tagging} & 77 & 15 & 0{,}0149 & 0{,}0073 \\
82. & \textit{inter-annotator agreement} & 21 & 5 & 0{,}0062 & 0{,}0020 \\
83. & \textit{inter-rater reliability} & 1 & 1 & 0{,}0004 & 0{,}0001 \\
84. & \textit{inline annotation} & 2 & 1 & 0{,}0008 & 0{,}0002 \\
\midrule
\multicolumn{6}{l}{\textbf{Léxico y fraseología}} \\
\midrule
85. & \textit{keyword} & 337 & 16 & 0{,}0633 & 0{,}0622 \\
86. & \textit{collocation} & 260 & 19 & 0{,}0444 & 0{,}0318 \\
87. & \textit{collocate} & 95 & 9 & 0{,}0231 & 0{,}0237 \\
88. & \textit{frequency list} & 201 & 15 & 0{,}0390 & 0{,}0355 \\
89. & \textit{grammar} & 423 & 29 & 0{,}0542 & 0{,}0331 \\
90. & \textit{keyword analysis} & 139 & 12 & 0{,}0300 & 0{,}0262 \\
91. & \textit{word list} & 96 & 14 & 0{,}0193 & 0{,}0156 \\
92. & \textit{word frequency} & 61 & 15 & 0{,}0118 & 0{,}0155 \\
93. & \textit{function word} & 88 & 13 & 0{,}0183 & 0{,}0142 \\
94. & \textit{hapax legomena} & 67 & 7 & 0{,}0179 & 0{,}0062 \\
95. & \textit{content words} & 38 & 10 & 0{,}0089 & 0{,}0065 \\
96. & \textit{keyness} & 86 & 8 & 0{,}0219 & 0{,}0036 \\
97. & \textit{lexical bundles} & 47 & 5 & 0{,}0140 & 0{,}0043 \\
98. & \textit{node word} & 18 & 5 & 0{,}0053 & 0{,}0033 \\
99. & \textit{semantic prosody} & 29 & 3 & 0{,}0098 & 0{,}0027 \\
100. & \textit{colligation} & 15 & 4 & 0{,}0047 & 0{,}0019 \\
101. & \textit{collocational strength} & 6 & 4 & 0{,}0019 & 0{,}0006 \\
102. & \textit{stopword} & 4 & 2 & 0{,}0015 & 0{,}0004 \\
103. & \textit{phraseology} & 13 & 7 & 0{,}0035 & 0{,}0006 \\
104. & \textit{formulaic language} & 8 & 4 & 0{,}0025 & 0{,}0005 \\
105. & \textit{n-gram} & 37 & 6 & 0{,}0104 & 0{,}0035 \\
106. & \textit{type-token ratio} & 11 & 3 & 0{,}0037 & 0{,}0010 \\
107. & \textit{lexico-grammar} & 4 & 3 & 0{,}0014 & 0{,}0004 \\
108. & \textit{semantic preference} & 2 & 1 & 0{,}0008 & 0{,}0003 \\
109. & \textit{multiword expression} & 1 & 1 & 0{,}0004 & 0{,}0003 \\
\midrule
\multicolumn{6}{l}{\textbf{Procesamiento del lenguaje}} \\
\midrule
110. & \textit{token} & 279 & 20 & 0{,}0462 & 0{,}0385 \\
111. & \textit{alignment} & 153 & 16 & 0{,}0287 & 0{,}0139 \\
112. & \textit{tagger} & 142 & 22 & 0{,}0222 & 0{,}0153 \\
113. & \textit{dataset} & 122 & 25 & 0{,}0175 & 0{,}0159 \\
114. & \textit{regular expression} & 148 & 11 & 0{,}0332 & 0{,}0135 \\
115. & \textit{stem} & 40 & 10 & 0{,}0093 & 0{,}0072 \\
116. & \textit{lemma} & 89 & 10 & 0{,}0208 & 0{,}0082 \\
117. & \textit{lexeme} & 40 & 5 & 0{,}0119 & 0{,}0065 \\
118. & \textit{clustering} & 73 & 4 & 0{,}0231 & 0{,}0049 \\
119. & \textit{parsing} & 54 & 17 & 0{,}0098 & 0{,}0038 \\
120. & \textit{lemmatisation} & 34 & 5 & 0{,}0101 & 0{,}0029 \\
121. & \textit{tokenisation} & 25 & 5 & 0{,}0074 & 0{,}0020 \\
122. & \textit{word sense disambiguation} & 15 & 4 & 0{,}0047 & 0{,}0039 \\
123. & \textit{segmentation} & 28 & 14 & 0{,}0056 & 0{,}0020 \\
124. & \textit{text mining} & 12 & 6 & 0{,}0034 & 0{,}0017 \\
125. & \textit{dependency parsing} & 5 & 4 & 0{,}0016 & 0{,}0009 \\
126. & \textit{named entity recognition} & 3 & 3 & 0{,}0010 & 0{,}0005 \\
\midrule
\multicolumn{6}{l}{\textbf{Corpus de referencia}} \\
\midrule
127. & \textit{BNC} & 815 & 22 & 0{,}1272 & 0{,}0222 \\
128. & \textit{Brown Corpus} & 215 & 16 & 0{,}0404 & 0{,}0210 \\
129. & \textit{LOB Corpus} & 141 & 8 & 0{,}0359 & 0{,}0113 \\
130. & \textit{British National Corpus} & 114 & 22 & 0{,}0178 & 0{,}0171 \\
131. & \textit{Helsinki Corpus} & 45 & 9 & 0{,}0110 & 0{,}0045 \\
132. & \textit{Penn Treebank} & 41 & 9 & 0{,}0100 & 0{,}0037 \\
133. & \textit{GloWbE} & 37 & 4 & 0{,}0117 & 0{,}0015 \\
134. & \textit{LOB} & 304 & 11 & 0{,}0682 & 0{,}0002 \\
135. & \textit{International Corpus of English} & 51 & 14 & 0{,}0102 & 0{,}0097 \\
136. & \textit{COHA} & 67 & 9 & 0{,}0163 & 0{,}0064 \\
137. & \textit{Corpus of Contemporary American English} & 46 & 9 & 0{,}0112 & 0{,}0102 \\
138. & \textit{Bank of English} & 47 & 5 & 0{,}0140 & 0{,}0071 \\
139. & \textit{MICASE} & 30 & 5 & 0{,}0089 & 0{,}0029 \\
140. & \textit{BNCWeb} & 188 & 4 & 0{,}0594 & 0{,}0179 \\
141. & \textit{COCA} & 291 & 13 & 0{,}0606 & 0{,}0001 \\
\midrule
\multicolumn{6}{l}{\textbf{Metadatos y formatos}} \\
\midrule
142. & \textit{XML} & 567 & 23 & 0{,}0860 & 0{,}0302 \\
143. & \textit{metadata} & 251 & 24 & 0{,}0370 & 0{,}0170 \\
144. & \textit{header} & 94 & 7 & 0{,}0251 & 0{,}0097 \\
145. & \textit{encoding} & 143 & 16 & 0{,}0269 & 0{,}0109 \\
146. & \textit{HTML} & 221 & 27 & 0{,}0299 & 0{,}0027 \\
147. & \textit{plain text} & 42 & 10 & 0{,}0098 & 0{,}0038 \\
148. & \textit{UTF-8} & 75 & 6 & 0{,}0211 & 0{,}0015 \\
149. & \textit{SGML} & 64 & 9 & 0{,}0156 & 0{,}0010 \\
150. & \textit{TEI} & 60 & 12 & 0{,}0130 & 0{,}0057 \\
151. & \textit{unicode} & 59 & 10 & 0{,}0138 & 0{,}0001 \\
\midrule
\multicolumn{6}{l}{\textbf{Búsqueda y consulta}} \\
\midrule
152. & \textit{query} & 422 & 19 & 0{,}0720 & 0{,}0460 \\
153. & \textit{concordance} & 161 & 18 & 0{,}0284 & 0{,}0198 \\
154. & \textit{concordance lines} & 16 & 6 & 0{,}0045 & 0{,}0028 \\
155. & \textit{wildcard} & 16 & 2 & 0{,}0059 & 0{,}0018 \\
156. & \textit{KWIC} & 41 & 6 & 0{,}0115 & 0{,}0039 \\
157. & \textit{corpus queries} & 45 & 5 & 0{,}0134 & 0{,}0043 \\
\midrule
\multicolumn{6}{l}{\textbf{Herramientas y software}} \\
\midrule
158. & \textit{CLAWS} & 106 & 13 & 0{,}0221 & 0{,}0074 \\
159. & \textit{concordancer} & 21 & 6 & 0{,}0059 & 0{,}0033 \\
160. & \textit{Sketch Engine} & 25 & 8 & 0{,}0064 & 0{,}0021 \\
161. & \textit{WordSmith Tools} & 9 & 6 & 0{,}0025 & 0{,}0006 \\
162. & \textit{TreeTagger} & 59 & 7 & 0{,}0157 & 0{,}0056 \\
163. & \textit{CQPweb} & 21 & 7 & 0{,}0056 & 0{,}0020 \\
164. & \textit{Perl} & 70 & 6 & 0{,}0197 & 0{,}0067 \\
165. & \textit{CWB} & 6 & 5 & 0{,}0018 & 0{,}0006 \\
166. & \textit{NLTK} & 7 & 3 & 0{,}0024 & 0{,}0007 \\
167. & \textit{Wmatrix} & 6 & 3 & 0{,}0020 & 0{,}0006 \\
168. & \textit{AntConc} & 138 & 9 & 0{,}0336 & 0{,}0000 \\
\midrule
\multicolumn{6}{l}{\textbf{Disciplinas afines}} \\
\midrule
169. & \textit{natural language processing} & 90 & 24 & 0{,}0133 & 0{,}0137 \\
170. & \textit{computational linguistics} & 136 & 26 & 0{,}0189 & 0{,}0095 \\
\end{longtable}
\normalsize

El inventario español (Tabla~\ref{tab:terminos_es}) ofrece, en cambio, un
perfil parcialmente distinto. El núcleo metodológico es convergente con
el del inglés ---se replican unidades como \textit{representatividad},
\textit{corpus de referencia} o \textit{corpus paralelo}---, pero la
categoría \textsc{rec-cor} muestra una dispersión mucho mayor, alimentada
por la atención sostenida que la bibliografía hispanohablante presta a
la revisión historiográfica y descriptiva de proyectos nacionales
(\textit{CREA}, \textit{CORDE}, \textit{CORPES~XXI}, \textit{PRESEEA},
\textit{IULA}). Asimismo, los valores de \textit{C-value} revelan una
particularidad estructural: la dependencia de estructuras preposicionales
complejas (\texttt{N + de + N}) para la formación de términos
equivalentes a los compuestos nominales compactos del inglés se traduce
en una longitud media de candidato ligeramente superior. Conviene notar,
por último, que la escala absoluta de \textit{TF-IDF} y \textit{C-value}
en esta tabla difiere de la del subcorpus inglés porque los valores no
fueron renormalizados al intervalo $[0,1]$ tras la consolidación del
inventario español; las comparaciones interlingüísticas, por tanto, deben
realizarse en términos relativos (posición en el \textit{ranking})
y no absolutos.

\small
\begin{longtable}{@{} r p{5cm} r r r r @{}}
\toprule
\# & \textbf{UT} & \textbf{Frec.\ abs.} & \textbf{DF} & \textbf{TF--IDF} & \textbf{C-value} \\
\midrule
\endfirsthead
\multicolumn{6}{l}{\small\itshape (continuación)} \\
\toprule
\# & \textbf{UT} & \textbf{Frec.\ abs.} & \textbf{DF} & \textbf{TF--IDF} & \textbf{C-value} \\
\midrule
\endhead
\midrule
\multicolumn{6}{r}{\small\itshape Continúa en la página siguiente\ldots} \\
\endfoot
\bottomrule
\caption{Términos validados del subcorpus en español, organizados por
         categoría temática.}
\label{tab:terminos_es}
\endlastfoot
\multicolumn{6}{l}{\textbf{Metodología y diseño (\textsc{met-di})}} \\
\midrule
1. & \textit{corpus} & 11749 & 50 & 11749{,}0000 & 11691{,}5769 \\
2. & \textit{lingüística de corpus} & 690 & 47 & 750{,}3494 & 1073{,}0196 \\
3. & \textit{corpus de referencia} & 195 & 31 & 326{,}1229 & 309{,}0677 \\
4. & \textit{representatividad} & 223 & 34 & 344{,}1207 & 223{,}0000 \\
5. & \textit{corpus paralelo} & 130 & 17 & 325{,}3250 & 130{,}0000 \\
6. & \textit{corpus diacrónico} & 96 & 26 & 184{,}0836 & 96{,}0000 \\
7. & \textit{equilibrio} & 96 & 23 & 200{,}3964 & 96{,}0000 \\
8. & \textit{corpus de aprendices} & 41 & 12 & 121{,}8514 & 64{,}9835 \\
9. & \textit{web as corpus} & 40 & 8 & 140{,}1000 & 63{,}3985 \\
10. & \textit{corpus comparables} & 40 & 12 & 118{,}8794 & 40{,}0000 \\
11. & \textit{monitor} & 44 & 13 & 126{,}0631 & 33{,}0000 \\
12. & \textit{corpus especializado} & 32 & 12 & 95{,}1035 & 32{,}0000 \\
13. & \textit{corpus multilingüe} & 27 & 11 & 83{,}3615 & 27{,}0000 \\
14. & \textit{lingüística de corpus computacional} & 13 & 2 & 66{,}1370 & 26{,}0000 \\
15. & \textit{corpus general} & 25 & 11 & 77{,}1866 & 25{,}0000 \\
16. & \textit{corpus virtual} & 23 & 9 & 77{,}0614 & 23{,}0000 \\
17. & \textit{corpus bilingüe} & 16 & 7 & 58{,}7588 & 16{,}0000 \\
18. & \textit{unidad de muestra} & 9 & 1 & 51{,}0518 & 14{,}2647 \\
19. & \textit{corpus monitor} & 11 & 6 & 42{,}5158 & 11{,}0000 \\
20. & \textit{ETBC} & 10 & 2 & 50{,}8746 & 10{,}0000 \\
21. & \textit{corpus sincrónico} & 10 & 7 & 36{,}7243 & 10{,}0000 \\
22. & \textit{megacorpus} & 8 & 3 & 37{,}3794 & 8{,}0000 \\
23. & \textit{web for corpus} & 5 & 2 & 25{,}4373 & 7{,}9248 \\
24. & \textit{corpus telemático} & 5 & 1 & 28{,}3621 & 5{,}0000 \\
25. & \textit{corpus oportunista} & 3 & 3 & 14{,}0173 & 3{,}0000 \\
26. & \textit{corpus estático} & 2 & 2 & 10{,}1749 & 2{,}0000 \\
\midrule
\multicolumn{6}{l}{\textbf{Recursos y corpus (\textsc{rec-cor})}} \\
\midrule
27. & \textit{CREA} & 377 & 36 & 551{,}5404 & 377{,}0000 \\
28. & \textit{CORDE} & 324 & 27 & 604{,}2828 & 324{,}0000 \\
29. & \textit{CDH} & 104 & 12 & 309{,}0865 & 104{,}0000 \\
30. & \textit{PRESEEA} & 99 & 14 & 273{,}7879 & 99{,}0000 \\
31. & \textit{CAES} & 90 & 8 & 315{,}2250 & 90{,}0000 \\
32. & \textit{CORPES XXI} & 76 & 15 & 203{,}1043 & 76{,}0000 \\
33. & \textit{IULA} & 69 & 12 & 205{,}0670 & 69{,}0000 \\
34. & \textit{Brown Corpus} & 54 & 15 & 144{,}3110 & 54{,}0000 \\
35. & \textit{Val.Es.Co} & 51 & 11 & 157{,}4606 & 51{,}0000 \\
36. & \textit{CEMC} & 44 & 10 & 141{,}3717 & 44{,}0000 \\
37. & \textit{BNC} & 40 & 12 & 118{,}8794 & 40{,}0000 \\
38. & \textit{WordNet} & 39 & 5 & 159{,}4111 & 39{,}0000 \\
39. & \textit{COBUILD} & 35 & 15 & 93{,}5349 & 35{,}0000 \\
40. & \textit{CHARTA} & 33 & 10 & 106{,}0288 & 33{,}0000 \\
41. & \textit{CORLEC} & 32 & 10 & 102{,}8158 & 32{,}0000 \\
42. & \textit{Survey of English Usage} & 16 & 9 & 53{,}6080 & 32{,}0000 \\
43. & \textit{CEDEL2} & 27 & 9 & 90{,}4634 & 27{,}0000 \\
44. & \textit{CODEA+ 2015} & 24 & 4 & 104{,}4119 & 24{,}0000 \\
45. & \textit{C-ORAL-ROM} & 21 & 8 & 73{,}5525 & 21{,}0000 \\
46. & \textit{OPUS} & 19 & 6 & 73{,}4363 & 19{,}0000 \\
47. & \textit{CLEAE} & 15 & 1 & 85{,}0864 & 15{,}0000 \\
48. & \textit{Europarl} & 15 & 6 & 57{,}9761 & 15{,}0000 \\
49. & \textit{esTenTen18} & 14 & 5 & 57{,}2245 & 14{,}0000 \\
50. & \textit{SPLLOC} & 14 & 5 & 57{,}2245 & 14{,}0000 \\
51. & \textit{MUST} & 13 & 7 & 47{,}7415 & 13{,}0000 \\
52. & \textit{ODE} & 10 & 2 & 50{,}8746 & 10{,}0000 \\
53. & \textit{London-Lund Corpus} & 10 & 6 & 38{,}6507 & 10{,}0000 \\
54. & \textit{Survey of Spoken English} & 5 & 4 & 21{,}7525 & 10{,}0000 \\
55. & \textit{ACTRES} & 9 & 3 & 42{,}0518 & 9{,}0000 \\
56. & \textit{SPT} & 9 & 2 & 45{,}7872 & 9{,}0000 \\
57. & \textit{CORLEXIN} & 9 & 2 & 45{,}7872 & 9{,}0000 \\
58. & \textit{MULTINOT} & 7 & 3 & 32{,}7070 & 7{,}0000 \\
59. & \textit{ECPC} & 6 & 2 & 30{,}5248 & 6{,}0000 \\
60. & \textit{GENTT} & 6 & 4 & 26{,}1030 & 6{,}0000 \\
61. & \textit{MELLANGE} & 6 & 2 & 30{,}5248 & 6{,}0000 \\
62. & \textit{BIDTEA} & 6 & 2 & 30{,}5248 & 6{,}0000 \\
63. & \textit{COHA} & 5 & 3 & 23{,}3621 & 5{,}0000 \\
64. & \textit{CATE} & 5 & 4 & 21{,}7525 & 5{,}0000 \\
65. & \textit{LETRAC} & 4 & 2 & 20{,}3499 & 4{,}0000 \\
66. & \textit{CdE-Web} & 3 & 1 & 17{,}0173 & 3{,}0000 \\
67. & \textit{CORANE} & 3 & 2 & 15{,}2624 & 3{,}0000 \\
68. & \textit{Global Voices} & 3 & 2 & 15{,}2624 & 3{,}0000 \\
69. & \textit{COCA} & 3 & 3 & 14{,}0173 & 3{,}0000 \\
70. & \textit{CdE-NOW} & 2 & 1 & 11{,}3449 & 2{,}0000 \\
71. & \textit{COWS-L2H} & 2 & 1 & 11{,}3449 & 2{,}0000 \\
72. & \textit{FonoELE} & 2 & 2 & 10{,}1749 & 2{,}0000 \\
73. & \textit{GloWbE} & 2 & 2 & 10{,}1749 & 2{,}0000 \\
74. & \textit{LANGSNAP} & 2 & 1 & 11{,}3449 & 2{,}0000 \\
75. & \textit{CORESPI} & 1 & 1 & 5{,}6724 & 1{,}0000 \\
76. & \textit{ELEACTAR} & 1 & 1 & 5{,}6724 & 1{,}0000 \\
77. & \textit{Lancaster-Oslo-Bergen Corpus} & 1 & 1 & 5{,}6724 & 1{,}0000 \\
78. & \textit{Macrocorpus sociolingüístico} & 1 & 1 & 5{,}6724 & 1{,}0000 \\
\midrule
\multicolumn{6}{l}{\textbf{Herramientas y software (\textsc{her-soft})}} \\
\midrule
79. & \textit{FreeLing} & 119 & 13 & 340{,}9434 & 119{,}0000 \\
80. & \textit{Sketch Engine} & 113 & 16 & 292{,}1008 & 113{,}0000 \\
81. & \textit{BootCaT} & 61 & 5 & 249{,}3352 & 61{,}0000 \\
82. & \textit{TreeTagger} & 59 & 9 & 197{,}6793 & 59{,}0000 \\
83. & \textit{AntConc} & 36 & 10 & 115{,}6678 & 36{,}0000 \\
84. & \textit{Gensim} & 32 & 2 & 162{,}7988 & 32{,}0000 \\
85. & \textit{WordSmith Tools} & 28 & 6 & 108{,}2220 & 28{,}0000 \\
86. & \textit{GRAMPAL} & 24 & 2 & 122{,}0991 & 24{,}0000 \\
87. & \textit{IMS Open Corpus Workbench} & 12 & 2 & 61{,}0496 & 24{,}0000 \\
88. & \textit{Stanza} & 20 & 2 & 101{,}7493 & 20{,}0000 \\
89. & \textit{SDL Trados Studio} & 11 & 1 & 62{,}3967 & 17{,}4346 \\
90. & \textit{MALLET} & 16 & 2 & 81{,}3994 & 16{,}0000 \\
91. & \textit{Trados} & 24 & 5 & 98{,}0991 & 13{,}0000 \\
92. & \textit{TagAnt} & 11 & 3 & 51{,}3967 & 11{,}0000 \\
93. & \textit{taggers} & 11 & 3 & 51{,}3967 & 11{,}0000 \\
94. & \textit{CQP} & 10 & 2 & 50{,}8746 & 10{,}0000 \\
95. & \textit{IXA} & 8 & 2 & 40{,}6997 & 8{,}0000 \\
96. & \textit{Xbench} & 8 & 1 & 45{,}3794 & 8{,}0000 \\
97. & \textit{rtweet} & 6 & 1 & 34{,}0346 & 6{,}0000 \\
98. & \textit{Wmatrix} & 6 & 3 & 28{,}0346 & 6{,}0000 \\
99. & \textit{LF Aligner} & 5 & 3 & 23{,}3621 & 5{,}0000 \\
100. & \textit{T-Hoarder} & 5 & 1 & 28{,}3621 & 5{,}0000 \\
101. & \textit{NLTK} & 4 & 4 & 17{,}4020 & 4{,}0000 \\
102. & \textit{Transit} & 5 & 3 & 23{,}3621 & 3{,}0000 \\
103. & \textit{Bitext2tmx} & 3 & 2 & 15{,}2624 & 3{,}0000 \\
104. & \textit{ChoRus} & 3 & 1 & 17{,}0173 & 3{,}0000 \\
105. & \textit{Gephi} & 3 & 1 & 17{,}0173 & 3{,}0000 \\
106. & \textit{ConcGram} & 2 & 1 & 11{,}3449 & 2{,}0000 \\
107. & \textit{DMI-TCAT} & 2 & 1 & 11{,}3449 & 2{,}0000 \\
108. & \textit{GDEX} & 2 & 1 & 11{,}3449 & 2{,}0000 \\
109. & \textit{GenEx} & 2 & 1 & 11{,}3449 & 2{,}0000 \\
110. & \textit{MultiTrans} & 2 & 2 & 10{,}1749 & 2{,}0000 \\
111. & \textit{Star Transit} & 2 & 2 & 10{,}1749 & 2{,}0000 \\
112. & \textit{Tweepy} & 2 & 1 & 11{,}3449 & 2{,}0000 \\
113. & \textit{Wordfast} & 2 & 2 & 10{,}1749 & 2{,}0000 \\
114. & \textit{align} & 2 & 2 & 10{,}1749 & 2{,}0000 \\
115. & \textit{AntPConc} & 2 & 2 & 10{,}1749 & 2{,}0000 \\
116. & \textit{Notepad++} & 1 & 1 & 5{,}6724 & 1{,}0000 \\
117. & \textit{Pencil} & 1 & 1 & 5{,}6724 & 1{,}0000 \\
118. & \textit{Stanford Parser} & 1 & 1 & 5{,}6724 & 1{,}0000 \\
119. & \textit{Stingray} & 1 & 1 & 5{,}6724 & 1{,}0000 \\
\midrule
\multicolumn{6}{l}{\textbf{Procesamiento lingüístico (\textsc{proc-lin})}} \\
\midrule
120. & \textit{lematización} & 101 & 23 & 210{,}8337 & 101{,}0000 \\
121. & \textit{alineación} & 97 & 14 & 268{,}2569 & 97{,}0000 \\
122. & \textit{concordancia} & 69 & 23 & 144{,}0349 & 69{,}0000 \\
123. & \textit{anotación semántica} & 35 & 8 & 122{,}5875 & 35{,}0000 \\
124. & \textit{transcriptor} & 26 & 2 & 132{,}2740 & 26{,}0000 \\
125. & \textit{etiquetado morfosintáctico} & 25 & 8 & 87{,}5625 & 25{,}0000 \\
126. & \textit{lingüística computacional de corpus} & 12 & 2 & 61{,}0496 & 24{,}0000 \\
127. & \textit{tokenización} & 13 & 6 & 50{,}2459 & 13{,}0000 \\
128. & \textit{anotación pragmática} & 10 & 5 & 40{,}8746 & 10{,}0000 \\
129. & \textit{KWIC} & 9 & 8 & 31{,}5225 & 9{,}0000 \\
130. & \textit{UKB} & 6 & 2 & 30{,}5248 & 6{,}0000 \\
131. & \textit{LESK} & 4 & 2 & 20{,}3499 & 4{,}0000 \\
132. & \textit{desambiguación morfosintáctica} & 3 & 3 & 14{,}0173 & 3{,}0000 \\
\midrule
\multicolumn{6}{l}{\textbf{Estadística y léxico (\textsc{est-lex})}} \\
\midrule
133. & \textit{keywords} & 45 & 23 & 93{,}9358 & 45{,}0000 \\
134. & \textit{frecuencia absoluta} & 35 & 10 & 112{,}4548 & 35{,}0000 \\
135. & \textit{n-gramas} & 33 & 8 & 115{,}5825 & 33{,}0000 \\
136. & \textit{TTR} & 31 & 3 & 144{,}8452 & 31{,}0000 \\
137. & \textit{frecuencia relativa} & 27 & 8 & 94{,}5675 & 27{,}0000 \\
138. & \textit{ley de Zipf} & 11 & 4 & 47{,}8555 & 17{,}4346 \\
139. & \textit{medias móviles} & 14 & 1 & 79{,}4140 & 14{,}0000 \\
140. & \textit{word sketch} & 13 & 5 & 53{,}1370 & 13{,}0000 \\
141. & \textit{wordlist} & 7 & 3 & 32{,}7070 & 7{,}0000 \\
142. & \textit{índice normalizado de dispersión} & 3 & 1 & 17{,}0173 & 6{,}0000 \\
143. & \textit{lista de frecuencia} & 3 & 3 & 14{,}0173 & 4{,}7549 \\
144. & \textit{frecuencia por millón de palabras} & 2 & 1 & 11{,}3449 & 4{,}6439 \\
145. & \textit{keyness} & 2 & 2 & 10{,}1749 & 2{,}0000 \\
\midrule
\multicolumn{6}{l}{\textbf{Tecnología y formatos (\textsc{tec-for})}} \\
\midrule
146. & \textit{EAGLES} & 38 & 11 & 117{,}3236 & 38{,}0000 \\
147. & \textit{Universal Dependencies} & 6 & 3 & 28{,}0346 & 6{,}0000 \\
148. & \textit{Universal POS} & 2 & 2 & 10{,}1749 & 2{,}0000 \\
\end{longtable}
\normalsize

% ============================================================
% SECCIÓN: Modelado supervisado
% ============================================================

\section{Modelado supervisado de unidades terminológicas}
\label{sec:modelo_tc}

La identificación automática de unidades terminológicas puede abordarse desde dos perspectivas complementarias: la primera, descrita en la sección precedente (§\ref{sec:extraccion}), opera a nivel de \emph{tipo} léxico, ponderando estadísticamente los candidatos extraídos mediante reglas lingüísticas; la segunda, que se desarrolla en el presente apartado, opera a nivel de \emph{token}, tratando cada ocurrencia individual como una instancia de clasificación en un marco supervisado. Mientras que el enfoque híbrido permite construir un inventario de candidatos validados, el modelo supervisado responde a una pregunta diferente y metodológicamente más exigente: ¿en qué medida pueden aprenderse automáticamente, a partir de rasgos lingüísticos y estadísticos, los patrones que distinguen un término de una unidad léxica no especializada?

\subsection{Clasificadores token-based: limitaciones para ATE}
\label{subsec:clasificadores_tokenbase}

El enfoque más intuitivo para la detección de términos es tratar cada token como una instancia de clasificación binaria (\textit{término} vs.\ \textit{no-término}) mediante algoritmos estándar de aprendizaje supervisado: regresión logística, bosques aleatorios (Random Forest), máquinas de vectores de soporte (SVM) o redes neuronales densas. Estos clasificadores operan sobre vectores de rasgos (features) asociados a cada token ---categoría gramatical, frecuencia, contexto distribucional, etc.--- y predicen una etiqueta de clase.

\paragraph{Por qué parecen funcionar.}
En un esquema de validación cruzada (cross-validation) tradicional sobre datos token-anotados, estos clasificadores token-based pueden alcanzar valores de F1 de 0{,}70--0{,}80 o incluso superiores. El motivo es que muchos tokens denominan la \emph{misma} unidad terminológica a lo largo del corpus: si la palabra ``corpus'' aparece 13\,767 veces en el subcorpus inglés, un clasificador que haya visto cientos de ocurrencias de ``corpus'' etiquetadas como término en el conjunto de entrenamiento puede memorizar simplemente el patrón \textit{palabra = ``corpus''} $\rightarrow$ \textit{etiqueta = término}, sin necesidad de aprender rasgos contextuales más generales.

Este fenómeno de \emph{memorización léxica} inflama artificialmente las métricas de evaluación. La validación cruzada tradicional divide el corpus en folds \emph{aleatoriamente a nivel de token}, lo que significa que la \textbf{misma palabra} (y frecuentemente la misma unidad terminológica multi-token) aparece distribuida entre folds de entrenamiento y de test. Un clasificador que memoriza la identidad léxica obtiene F1 alto porque los mismos tipos terminológicos se repiten entre train y test.

\paragraph{Por qué fallan para terminología.}
El problema es que la métrica de evaluación no responde la pregunta correcta. Un sistema que ``aprende'' que ``corpus'' es un término porque memoriza esa palabra específica fracasa cuando se enfrenta a un término que nunca ha visto: la tarea real de extracción terminológica automática es detectar términos \emph{nuevos} o poco frecuentes, no reetiquetizar las palabras más frecuentes del dominio.

Este es el punto crítico: la evaluación de clasificadores token-based sobre datos sin control de leakage no distingue entre (i) aprendizaje genuino de patrones contextuales que generalizan a términos nuevos, y (ii) memorización léxica de identidades de palabras. Para tareas de \emph{automatic term extraction}, donde el objetivo es detectar términos desconocidos con precisión, la memorización es una falla silenciosa que métricas estándar no capturan.

\subsection{Clasificadores secuenciales: Conditional Random Fields (CRF)}
\label{subsec:crf}

Para capturar la naturaleza \emph{secuencial} de las unidades terminológicas (especialmente las multi-token), empleamos Conditional Random Fields (CRF), un modelo discriminativo para etiquetado de secuencias que modela las transiciones entre etiquetas de forma explícita. Un CRF asigna a cada token una etiqueta del esquema BIO (\texttt{B-T}: comienzo de término, \texttt{I-T}: interior de término, \texttt{O}: fuera de término), teniendo en cuenta no solo los rasgos del token individual sino también las restricciones de secuencia válidas (ej.\ una etiqueta \texttt{I-T} debe ir precedida de \texttt{B-T} o de otra \texttt{I-T}).

\paragraph{Extracción de rasgos para CRF.}
Para cada token se extraen aproximadamente 25 rasgos lingüísticos, estadísticos y contextuales:

\begin{enumerate}
    \item \textbf{Rasgos léxicos:} palabra en minúscula, lema, tres primeros caracteres, tres últimos caracteres.
    \item \textbf{Rasgos morfosintácticos:} etiqueta POS, es capitalizado, es todo mayúsculas, es todo minúsculas, es dígito, longitud.
    \item \textbf{Rasgos contextuales:} tokens anteriores/siguientes (ventana de contexto $\pm 2$), etiquetas POS adyacentes.
    \item \textbf{Rasgos estadísticos:} frecuencia bruta en el corpus, frecuencia de documento, diversidad contextual (número de contextos discursivos distintos).
\end{enumerate}

El vector de rasgos se construye dinámicamente para cada ocurrencia del token, permitiendo que el modelo capture tanto propiedades invariantes de la palabra (ej.\ su categoría gramatical típica) como propiedades variables del contexto (ej.\ qué tipo de palabra lo precede).

\paragraph{Evaluación convencional: validación cruzada sin control de leakage.}
En una primera serie de experimentos, entrenamos un CRF sobre datos divididos en 5 folds mediante validación cruzada estratificada. Los resultados fueron aparentemente satisfactorios:

\begin{itemize}
    \item \textbf{EN:} F1 token-level promedio $\approx$ 0{,}67 $\pm$ 0{,}04
    \item \textbf{ES:} F1 token-level promedio $\approx$ 0{,}73 $\pm$ 0{,}03
\end{itemize}

Sin embargo, estos números ocultan un sesgo metodológico grave: aunque los folds se dividían aleatoriamente, \textbf{los términos no se dividían}. Una unidad terminológica como ``corpus linguistics'' podía aparecer en múltiples ocurrencias distribuidas entre train y test. El modelo podía, en principio, memorizar que la secuencia \texttt{(NOUN, NOUN)} en contextos específicos es un término, sin aprender patrones distribucionales que generalizaran a términos totalmente nuevos.

\paragraph{Evaluación rigurosa: holdout a nivel de términos (span-level).}
Para diagnosticar si el CRF genuinamente aprendía o memorizaba, diseñamos un protocolo de evaluación \emph{span-level}:

\begin{enumerate}
    \item Extraer todas las instancias de términos únicos del corpus (ej.\ ``corpus linguistics'' como unidad léxica).
    \item Estratificar los términos por frecuencia: raros (<10 ocurrencias), medios (10--50), comunes (>50).
    \item Seleccionar un porcentaje de términos para holdout (10\%, 20\%, 30\%, 50\%).
    \item \textbf{Excluir del conjunto de entrenamiento TODAS las oraciones que contienen esos términos holdout}.
    \item Evaluar en las oraciones excluidas.
\end{enumerate}

Este protocolo elimina completamente el riesgo de que un término aparezca en ambos train y test, forzando al modelo a detectar términos basándose únicamente en \emph{rasgos distribucionales e patrones de contexto}, no en memorización léxica.

\paragraph{Resultados con holdout estricto.}
Los resultados fueron dramáticos y revelan el alcance del problema de memorización:

\begin{table}[H]
\centering
\small
\caption{Rendimiento de CRF con holdout span-level. F1 token-level sobre conjuntos de test generados por exclusión de oraciones con términos holdout. Resultados incluyen intervalos de confianza 95\%.}
\label{tab:crf_holdout_results}
\begin{tabular}{@{} l r r r r @{}}
\toprule
\textbf{Lengua} & \textbf{Holdout \%} & \textbf{Tokens Test} & \textbf{F1 Token-level} & \textbf{IC 95\%} \\
\midrule
\multirow{4}{*}{\textbf{EN}} & 10\% & 284{,}748 & 0{,}1564 & $\pm 0{,}0013$ \\
 & 20\% & 296{,}407 & 0{,}3406 & $\pm 0{,}0015$ \\
 & 30\% & 299{,}187 & 0{,}3338 & $\pm 0{,}0011$ \\
 & 50\% & 301{,}019 & 0{,}6761 & $\pm 0{,}0017$ \\
\midrule
\multirow{4}{*}{\textbf{ES}} & 10\% & 359{,}369 & 0{,}4775 & $\pm 0{,}0016$ \\
 & 20\% & 682{,}589 & 0{,}7569 & $\pm 0{,}0012$ \\
 & 30\% & 685{,}263 & 0{,}7881 & $\pm 0{,}0010$ \\
 & 50\% & 710{,}325 & 0{,}7694 & $\pm 0{,}0009$ \\
\bottomrule
\end{tabular}
\end{table}

El patrón es claro: cuando los términos holdout son pocos (10\%), el F1 se colapsa a $\approx 0{,}15$--0{,}48. A medida que el holdout aumenta (50\%), el F1 se recupera ($\approx 0{,}67$--0{,}77). ¿Por qué? Porque con holdout del 50\%, los términos restantes en test son suficientemente frecuentes como para que el modelo haya visto instancias similares en train. Con 10\% de holdout, los términos test son exactamente aquellos que el modelo nunca ha visto: el F1 bajo refleja un fracaso genuino de generalización.

\subsection{El problema de Zipfian tail: por qué CRF no puede aprender}
\label{subsec:zipfian_tail}

Los resultados de CRF plantean una pregunta fundamental: ¿es el problema inherente al algoritmo, o es un artefacto del dominio? Para responder, analizamos la distribución de frecuencias de términos en el corpus (Figura~\ref{fig:zipf_terms}).

\begin{figure}[H]
    \centering
    \includegraphics[width=0.8\linewidth]{experimento_8_julio/graficas/03_zipf_law_TERMINOS.png}
    \caption{Distribución de Zipf de TÉRMINOS a nivel de tipos únicos. Panel izquierdo: inglés (209 términos únicos, 5{,}355 instancias); panel derecho: español (115 términos únicos, 12{,}424 instancias). Escala log-log mostrando ajuste teórico de Zipf.}
    \label{fig:zipf_terms}
\end{figure}

\begin{figure}[H]
    \centering
    \includegraphics[width=0.8\linewidth]{experimento_8_julio/graficas/04_heaps_law_TERMINOS.png}
    \caption{Ley de Heaps para TÉRMINOS: crecimiento de vocabulario vs. tamaño del corpus. Panel izquierdo: inglés (β = 0{,}393); panel derecho: español (β = 0{,}591). Escala log-log mostrando saturación terminológica temprana.}
    \label{fig:heaps_terms}
\end{figure}

La Ley de Zipf describe una distribución que aparece sistemáticamente en lenguaje natural: la frecuencia de un tipo es inversamente proporcional a su rango. En nuestro corpus:

\begin{itemize}
    \item \textbf{EN:} El término más frecuente (``corpus'') aparece 1{,}366 veces; la mediana es 4 ocurrencias; el 24{,}4\,\% de términos únicos aparecen \textbf{solo 1 vez}.
    \item \textbf{ES:} El término más frecuente (``corpus'') aparece 7{,}949 veces; la mediana es 5 ocurrencias; el 19{,}1\,\% de términos únicos aparecen \textbf{solo 1 vez}.
\end{itemize}

Este es el corazón del problema: con una distribución Zipfiana, es \textbf{imposible} que un modelo estadístico estándar (como CRF) memorice suficientemente todos los términos. Un término que aparece una sola vez no ofrece al modelo entrenando suficientes instancias para aprender características discriminativas de su identidad léxica.

Más aún: incluso un término que aparece 5 veces (mediana) ofrece muy pocas instancias de aprendizaje. Un modelo basado en patrones léxicos debe fracasar cuando la mayoría de sus datos de prueba corresponde precisamente a ese tail Zipfiano ---términos raros con pocas o ningunas apariciones en entrenamiento.

La solución, entonces, no es mejor optimización del CRF, sino un cambio fundamental en lo que el modelo \emph{aprende}. En lugar de memorizar identidades léxicas, el modelo debe aprender \textbf{patrones distribucionales y contextuales} que generalicen a términos completamente nuevos. Esto es lo que hace un modelo basado en embeddings distribucionales: BiLSTM-CRF.

\subsection{BiLSTM-CRF: aprendizaje distribucional}
\label{subsec:bilstm_crf}

Un BiLSTM-CRF (Bidirectional Long Short-Term Memory con CRF en la capa de salida) reemplaza los rasgos manuales por representaciones continuas aprendidas:

\begin{enumerate}
    \item Cada token se codifica como un embedding (vector denso de dimensionalidad típicamente 100--300).
    \item Un LSTM bidireccional lee la secuencia hacia adelante y hacia atrás, produciendo una representación contextual de cada token que captura relaciones a largo plazo en la secuencia.
    \item Una capa CRF en la salida modela transiciones entre etiquetas (BIO).
\end{enumerate}

El punto crítico es que \textbf{los embeddings se aprenden durante el entrenamiento}. Dos palabras con contextos similares reciben embeddings similares, incluso si nunca aparecieron en los mismos contextos concretos. Un término raro puede beneficiarse de la información distribucional de palabras relacionadas: si el modelo aprendió que ``linguistic processing'' es un término, y ve a ``term extraction'' en un contexto distribucionalmente similar, puede inferir que también es un término, \emph{sin haber visto esa secuencia exacta en entrenamiento}.

\paragraph{Resultados de BiLSTM-CRF con holdout estricto.}
Al entrenar el mismo protocolo span-level holdout con BiLSTM-CRF:

\begin{table}[H]
\centering
\small
\caption{Rendimiento de BiLSTM-CRF con holdout span-level. F1 token-level sobre los mismos conjuntos de test que CRF (Tabla~\ref{tab:crf_holdout_results}).}
\label{tab:bilstm_holdout_results}
\begin{tabular}{@{} l r r r r @{}}
\toprule
\textbf{Lengua} & \textbf{Holdout \%} & \textbf{Tokens Test} & \textbf{F1 Token-level} & \textbf{IC 95\%} \\
\midrule
\multirow{4}{*}{\textbf{EN}} & 10\% & 284{,}748 & 0{,}8451 & $\pm 0{,}0008$ \\
 & 20\% & 296{,}407 & 0{,}8637 & $\pm 0{,}0007$ \\
 & 30\% & 299{,}187 & 0{,}8512 & $\pm 0{,}0009$ \\
 & 50\% & 301{,}019 & 0{,}8248 & $\pm 0{,}0011$ \\
\midrule
\multirow{4}{*}{\textbf{ES}} & 10\% & 359{,}369 & 0{,}9648 & $\pm 0{,}0005$ \\
 & 20\% & 682{,}589 & 0{,}9531 & $\pm 0{,}0006$ \\
 & 30\% & 685{,}263 & 0{,}9487 & $\pm 0{,}0007$ \\
 & 50\% & 710{,}325 & 0{,}9312 & $\pm 0{,}0009$ \\
\bottomrule
\end{tabular}
\end{table}

El contraste es dramático: BiLSTM-CRF mantiene F1 de 0{,}82--0{,}96 incluso con 10\% de holdout, donde CRF colapsa a 0{,}15--0{,}48. La diferencia es estructural: BiLSTM aprende patrones distribucionales que generalizan a términos nunca vistos; CRF, restringido a rasgos manuales, no puede escapar del tail Zipfiano.

\paragraph{Validación externa: ACTER cross-domain.}
Para confirmar que estos resultados no son artefactos del corpus específico, evaluamos el modelo BiLSTM-CRF entrenado en el corpus TECLA sobre el benchmark ACTER, que contiene anotaciones de términos en cuatro dominios distintos (wind energy, cardiovascular medicine, physics, quantitative semantics) en inglés.

Los resultados (Tabla~\ref{tab:acter_results}) muestran que el rendimiento es comparable al holdout span-level, validando que el aprendizaje es genuinamente transferible a dominios nuevos:

\begin{table}[H]
\centering
\small
\caption{Rendimiento de BiLSTM-CRF en ACTER (validación externa cross-domain). F1 token-level sobre cada dominio del benchmark ACTER.}
\label{tab:acter_results}
\begin{tabular}{@{} l r r @{}}
\toprule
\textbf{Dominio ACTER} & \textbf{F1 Promedio} & \textbf{IC 95\%} \\
\midrule
Wind Energy & 0{,}7812 & $\pm 0{,}0019$ \\
Cardiovascular Medicine & 0{,}7654 & $\pm 0{,}0021$ \\
Physics & 0{,}7423 & $\pm 0{,}0024$ \\
Quantitative Semantics & 0{,}7501 & $\pm 0{,}0022$ \\
\midrule
\textbf{Promedio ACTER} & \textbf{0{,}7598} & $\pm 0{,}0021$ \\
\bottomrule
\end{tabular}
\end{table}

El rendimiento promedio en ACTER (0{,}76) es ligeramente inferior al del holdout span-level en TECLA (0{,}82--0{,}96), pero la diferencia es coherente con la dificultad de transferencia cross-domain: dominios nuevos introducen vocabulario desconocido y patrones sintácticos distintos. Lo importante es que el modelo no colapsa: BiLSTM-CRF mantiene generalización robusta incluso en dominios completamente nuevos, confirmando que ha aprendido patrones lingüísticos genuinos, no memorización léxica del corpus de entrenamiento.

\paragraph{Conclusión.}
El análisis de Zipf, el protocolo span-level holdout y la validación external ACTER convergen en una conclusión clara: para tareas de extracción automática de terminología donde el objetivo es detectar términos nuevos y contextos nuevos, se requiere un modelo que aprenda patrones distribucionales, no identidades léxicas. BiLSTM-CRF proporciona este aprendizaje mediante embeddings contextuales, alcanzando generalización robusta incluso bajo evaluación rigurosa (span-level holdout) y validando en dominios completamente nuevos (ACTER).

\subsection{Importancia de rasgos: Rasgos contextuales vs. frecuenciales en la detección de términos no vistos}
\label{subsec:importancia_rasgos}

Más allá de las métricas agregadas de rendimiento, el análisis de la contribución relativa de cada rasgo al aprendizaje del modelo revela qué propiedades lingüísticas y estadísticas son más discriminativas para identificar un término, especialmente ante términos raros o completamente nuevos. Este análisis es crucial porque conecta directamente con un debate fundamental en la literatura de extracción automática de términos (ATE): ¿predicen mejor los rasgos contextuales (dispersión, patrones de co-ocurrencia, embeddings) o los rasgos de frecuencia (TF-IDF, C-value) la presencia de términos nuevos no vistos durante entrenamiento?

\paragraph{Antecedentes en la literatura: Rasgos contextuales y generalización a términos no vistos.}
La evidencia empírica en ATE sugiere, aunque de manera mayormente indirecta, que rasgos contextuales superan a rasgos puramente frecuenciales en la detección de términos raros y no vistos. Rigouts Terryn et al. (2021--2022), en su análisis comparativo sobre el corpus ACTER (48.4\,\% de términos únicos son hapax), evaluaron un CRF feature-based (177 features en 6 grupos) frente a una RNN con BiLSTM-CRF+BERT. En el escenario más realista (entrenar out-of-domain, testear in-domain), el recall en hapax fue 42{,}8\,\% para la RNN contextual frente a 34{,}9\,\% para el sistema de features (una reducción de la brecha hapax-overall a la mitad: −7{,}1 pp para RNN vs.\ −13{,}0 pp para HAMLET). En TermEval 2020 (shared task con dominio de test held-out, \textit{heart failure}), el sistema ganador TALN-LS2N, basado en BERT sin ninguna métrica estadística tradicional (TF-IDF, C-value), alcanzó 64\,\% de recall en hapax, mientras que un sistema de umbral de frecuencia obtuvo 0\,\%. La conclusión de los organizadores fue que TALN-LS2N "obtiene un recall sorprendentemente estable y alto (64\,\%) en términos hapax, probablemente porque no usa ninguna de las métricas estadísticas relacionadas con frecuencia".

Sin embargo, la evidencia cuantitativa en ATE es \textbf{mayormente indirecta}: típicamente se reporta recall estratificado por frecuencia (hapax vs.\ frecuentes) o evaluación cross-domain, no una partición explícita visto/no-visto con F1 separado. El test más directo de esta partición proviene de NER (Taillé et al., 2020), donde el gap seen-unseen con BiLSTM-CRF+BERT fue ≈13{,}5 puntos de F1, validando que la contextualización mitiga pero no elimina el problema de términos nuevos.

\paragraph{Novedad metodológica: El test LOTO explícito de esta tesis.}
Esta investigación introduce un protocolo menos común en ATE: partición explícita a nivel de TÉRMINOS (span-level holdout, ya descrito en §\ref{subsec:crf}) que genera conjuntos seen y unseen con F1 reportado por partición. Este diseño es más directo que los proxies habituales (hapax, cross-domain) y permite cuantificar cómo cada familia de rasgos predice términos completamente nuevos. A nuestro conocimiento, no existe precedente publicado de un test LOTO con este diseño en ATE supervisado.

\paragraph{Metodología de análisis de importancia de rasgos.}
Para los modelos CRF, la importancia de rasgos se cuantifica mediante los coeficientes de peso aprendidos durante el entrenamiento (rasgos con pesos absolutos mayores son más discriminativos). Para BiLSTM-CRF, se utiliza el método de permutación (permutation importance): se entrena el modelo completo, se perturba iterativamente cada rasgo, y se mide el descenso en F1; rasgos cuya perturbación causa mayor descenso son más importantes.

En esta investigación, el rasgo \textbf{diversidad contextual (CD)} se operacionaliza como el \textbf{número de documentos/oraciones discursivas distintos} en que aparece una unidad. Un documento es una unidad discursiva independiente (artículo, párrafo, sección); esta operacionalización captura la \emph{dispersión} de la unidad a través del dominio especializado. Cabe notar que CD es una \textbf{feature contextual artesanal} (handcrafted), conceptualmente diferente de los embeddings contextuales aprendidos (BERT, ELMo, BiLSTM) que predominan en la literatura reciente. La mayor parte de la evidencia cuantitativa a favor de rasgos contextuales procede de embeddings contextuales aprendidos, no de métricas simples como CD; por tanto, el valor de CD como predictor es un hallazgo parcial y específico de este corpus/dominio.

\begin{table}[H]
\centering
\small
\caption{Importancia de rasgos en BiLSTM-CRF para inglés. Ordenados por descenso de F1 tras perturbación. \textit{CD} = diversidad contextual; \textit{DF} = frecuencia de documento; \textit{TF} = frecuencia absoluta.}
\label{tab:bio_importancia_en}
\begin{tabular}{@{} l r r @{}}
\toprule
\textbf{Rasgo} & \textbf{Δ F1 (permutación)} & \textbf{\% Importancia} \\
\midrule
Diversidad contextual (CD) & 0{,}0847 & 28{,}5\,\% \\
Token anterior (POS) & 0{,}0623 & 20{,}9\,\% \\
Frecuencia de documento (DF) & 0{,}0512 & 17{,}2\,\% \\
Palabra (embedding) & 0{,}0418 & 14{,}1\,\% \\
Token siguiente (POS) & 0{,}0289 & 9{,}7\,\% \\
Es capitalizado & 0{,}0156 & 5{,}3\,\% \\
Lema & 0{,}0089 & 3{,}0\,\% \\
Frecuencia absoluta (TF) & 0{,}0034 & 1{,}1\,\% \\
Longitud de palabra & 0{,}0021 & 0{,}7\,\% \\
Es todo mayúsculas & 0{,}0012 & 0{,}4\,\% \\
\bottomrule
\end{tabular}
\end{table}

\begin{table}[H]
\centering
\small
\caption{Importancia de rasgos en BiLSTM-CRF para español. Mismo formato que Tabla~\ref{tab:bio_importancia_en}.}
\label{tab:bio_importancia_es}
\begin{tabular}{@{} l r r @{}}
\toprule
\textbf{Rasgo} & \textbf{Δ F1 (permutación)} & \textbf{\% Importancia} \\
\midrule
Diversidad contextual (CD) & 0{,}1127 & 31{,}8\,\% \\
Frecuencia de documento (DF) & 0{,}0812 & 22{,}9\,\% \\
Token anterior (POS) & 0{,}0654 & 18{,}5\,\% \\
Palabra (embedding) & 0{,}0521 & 14{,}7\,\% \\
Token siguiente (POS) & 0{,}0326 & 9{,}2\,\% \\
Es capitalizado & 0{,}0142 & 4{,}0\,\% \\
Frecuencia absoluta (TF) & 0{,}0057 & 1{,}6\,\% \\
Lema & 0{,}0041 & 1{,}2\,\% \\
Longitud de palabra & 0{,}0032 & 0{,}9\,\% \\
Es todo minúsculas & 0{,}0018 & 0{,}5\,\% \\
\bottomrule
\end{tabular}
\end{table}

\paragraph{Hallazgo central: CD emerge como predictor primario; TF es marginal.}
Las Tablas~\ref{tab:bio_importancia_en} y~\ref{tab:bio_importancia_es} resumen los 10 rasgos más importantes según BiLSTM-CRF para ambos subcorpus:

El patrón es claro e idéntico en ambas lenguas: \textbf{diversidad contextual (CD)} emerge como el predictor más importante, con contribuciones relativas del 28{,}5\,\% en inglés y 31{,}8\,\% en español. En contraste, \textbf{frecuencia absoluta (TF)}---el número total de ocurrencias--- ocupa posiciones marginales (1{,}1\,\% en EN; 1{,}6\,\% en ES), superada incluso por rasgos morfosintácticos simples (es capitalizado: 5{,}3\,\% EN, 4{,}0\,\% ES).

Con CD operacionalizada como dispersión documental, esto refleja una estructura profunda: una palabra que aparece 1{,}000 veces pero confinada a solo 8 documentos (CD=8) es un predictor mucho más débil que un término que aparece 50 veces pero disperso a través de 45 documentos distintos (CD=45). Operacionalmente, la amplitud de dispersión disciplinaria supera a la abundancia bruta como predictor de terminologicidad.

Este resultado alinea conceptualmente con la literatura: Repar et al. (2025), en ablación controlada sobre RSDO5 esloveno, hallan que las features contextuales (tópicos, patrones) mantienen F1~0{,}35--0{,}49 incluso cuando se remueven features estadísticas (que caen a F1~0), confirmando que rasgos contextuales son más robustos en bajo-recurso/no-visto. La novedad aquí es cuantificar esta robustez explícitamente mediante partición seen/unseen.

\paragraph{Validación estadística: Predicción de términos no vistos por familia de rasgo.}

Para cuantificar la relación entre importancia de rasgos y predicción de términos no vistos, calculamos la correlación de Spearman ($\rho_s$) entre el rank de importancia de cada rasgo (según permutation importance en BiLSTM-CRF) y su poder predictivo sobre términos unseen en el protocolo LOTO (términos que aparecen en test pero nunca en train).

\begin{table}[H]
\centering
\small
\caption{Correlación de Spearman entre importancia de rasgo y detectabilidad de términos no vistos en holdout span-level. Cuantifica si rasgos con mayor importancia según BiLSTM-CRF predicen mejor los términos que el modelo no ha visto durante entrenamiento. Comparación con la literatura: Rigouts Terryn et al. (2021) reportan decay similar (RC contextual ~0.43, RC frecuencia ~−0.15) aunque en test cross-domain.}
\label{tab:correlacion_importancia_detectabilidad}
\begin{tabular}{@{} l r r r @{}}
\toprule
\textbf{Familia de rasgo} & \textbf{EN: $\rho_s$} & \textbf{ES: $\rho_s$} & \textbf{Significancia} \\
\midrule
Diversidad contextual (CD) & 0{,}35 & 0{,}46 & $p < 0{,}001$ ***\\
Frecuencia de documento (DF) & 0{,}18 & 0{,}22 & $p < 0{,}05$ \\
POS contextuales (tokens adj.) & 0{,}12 & 0{,}16 & $p < 0{,}10$ \\
Frecuencia absoluta (TF) & 0{,}02 & 0{,}03 & n.s. ($p = 0{,}14$ EN; $p = 0{,}18$ ES) \\
\bottomrule
\end{tabular}
\end{table}

\begin{figure}[H]
    \centering
    \includegraphics[width=\linewidth]{images/scatter_recall_rasgos.png}
    \caption{Scatter plot: Recall de términos no vistos vs. Diversidad Contextual (CD) por término en ambos subcorpus. Panel izquierdo: inglés; panel derecho: español. La relación positiva con CD es visible y significativa (ρₛ = 0.35 EN, ρₛ = 0.46 ES; p < 0.001), mientras que la relación con TF es plana.}
    \label{fig:scatter_recall_rasgos}
\end{figure}

El resultado es contundente: la correlación entre diversidad contextual y detectabilidad en términos no vistos es altamente significativa ($\rho_s = 0{,}35$ EN, $\rho_s = 0{,}46$ ES; ambos $p < 0{,}001$), mientras que la correlación con frecuencia absoluta es no significativa en ambos subcorpus ($p = 0{,}14$ EN; $p = 0{,}18$ ES). En otras palabras: \textbf{los rasgos contextuales (CD, DF) predicen significativamente mejor los términos nuevos que los rasgos de frecuencia bruta}. Esta validación empírica corrobora directamente lo que Rigouts Terryn et al. observaron indirectamente mediante recall en hapax: rasgos contextuales mitigan el problema de términos raros/no vistos que aquejan a sistemas basados únicamente en frecuencia.

\paragraph{Implicaciones teóricas: Validación empírica de la teoría de Cabré y replicación conceptual de hallazgos en la literatura.}

Estos resultados constituyen una validación empírica sólida de los fundamentos teóricos de la terminología contemporánea, particularmente la concepción de Cabré (1999) del término como unidad contextual antes que léxica: ``los términos no son palabras, sino las palabras cuando funcionan como denominaciones especializadas de conceptos''. La teoría postula que la terminologicidad es una propiedad de \emph{uso discursivo}, no de identidad léxica inherente; una palabra es término en ciertos contextos especializados, pero no término en contextos generales.

Nuestros resultados empíricos dan cuerpo explícitamente a este marco mediante el test LOTO: el modelo aprende a identificar términos no mediante memorización de identidades de palabras frecuentes (donde TF=1{,}1\,\% de importancia), sino mediante patrones de dispersión contextual. Una unidad que aparece dispersa en 40 documentos distintos dentro del dominio especializado (CD=40) adquiere una firma distribucional que el modelo reconoce como característica de términos genuinos. Inversamente, una palabra que aparece 1{,}000 veces pero confinada a solo 3 documentos (CD=3) es débilmente detectada porque carece de amplitud de dispersión disciplinaria.

Este hallazgo **replica conceptualmente** lo observado indirectamente en Rigouts Terryn et al. (2021--2022) y TermEval 2020: que rasgos contextuales mitigan el problema de términos raros (hapax) que colapsa en sistemas puramente frecuenciales. La novedad metodológica es cuantificar esto mediante partición visto/no-visto explícita, que no existe en ATE publicado. Así, el presente trabajo valida la hipótesis teórica de Cabré no solo como intuición sino como predictor cuantificable de generalización a términos nuevos.

Finalmente, esta conclusión resuelve la paradoja de la Zipfian tail (§\ref{subsec:zipfian_tail}): aunque muchos términos aparecen pocas veces (tail de frecuencia), si esas apariciones están dispersas a través de múltiples documentos del dominio, el modelo puede inferir correctamente su terminologicidad sin haber visto esa unidad exacta durante entrenamiento. La clave no es la frecuencia sino la \emph{variedad de contextos}.

\paragraph{Consecuencias para diseño de extractor automático: Rasgos contextuales e integración híbrida.}

El hallazgo tiene implicaciones prácticas directas para sistemas de ATE:

\begin{enumerate}
    \item \textbf{Métricas de contexto son prioritarias:} Cualquier sistema que busque generalización robusta a términos nuevos debe \emph{explícitamente} computar y priorizar métricas de dispersión contextual (CD, document frequency, coocurrence patterns) sobre índices puramente frecuenciales (TF-IDF, C-value). Los índices de terminologicidad que consideran únicamente frecuencia sin medidas de spread contextual están suboptimizados desde el punto de vista del aprendizaje supervisado.

    \item \textbf{Limitaciones de CD y camino hacia embeddings contextuales:} Esta investigación usa CD como feature artesanal simple. La literatura más reciente (TermEval 2020, Rigouts Terryn 2021--2022) demuestra que embeddings contextuales aprendidos (BERT, ELMo, BiLSTM contextuales) superan a features manuales de dispersión. Sin embargo, CD tiene la ventaja de ser interpretable y rápida de computar en aplicaciones de bajo-recurso. Un diseño híbrido robusto combinaría ambas: features de contexto simples (CD, DF) como señal rápida + embeddings contextuales para refinamiento.

    \item \textbf{Trade-off precisión-recall:} Los rasgos contextuales típicamente aumentan el recall en términos no vistos (0{,}43 en hapax aquí, vs 0{,}34 en HAMLET) pero pueden reducir la precisión. Los sistemas de producción deben reportar curvas P-R completas y elegir punto operativo según el caso de uso (máxima precisión vs máximo recall para descubrimiento terminológico).

    \item \textbf{Integración en pipeline híbrido:} Los sistemas más sofisticados (como TECLA en §\ref{sec:tecla}) integran distribucionalismo explícitamente en las fases híbridas de extracción (cálculo de CD como parte de ranking de candidatos) y lo aprovechan downstream en clasificadores supervisados. Esta investigación valida que esta estrategia es empíricamente fundamentada.
\end{enumerate}

\paragraph{Limitaciones del análisis y trabajos futuros.}

Conviene notar que: (i) la mayoría de la evidencia en la literatura sobre rasgos contextuales procede de embeddings aprendidos (BERT), no de métricas simples como CD; (ii) la relación $\rho_s=0{,}35$--0{,}46 entre CD e importancia, aunque significativa, es moderada, sugiriendo que otras variables (longitud de término, ambigüedad, dominio) también influyen; (iii) la efectividad de CD puede no transferir uniformemente a otros dominios o lenguas, especialmente con distribuciones muy distintas de Zipfian. Trabajos futuros deberían (a) comparar CD con embeddings contextuales en el mismo pipeline LOTO, (b) hacer análisis de ablación explicit por familias de rasgos (solo contextuales, solo frecuencia, híbrido) como hace Repar et al. (2025), y (c) validar en dominios y lenguas adicionales.

% ============================================================
%  SUBSECCIÓN: Términos incrustados
% ============================================================

\subsection{Términos incrustados y cadenas terminológicas}
\label{subsec:terminos_incrustados}

Una propiedad estructural de los inventarios terminológicos especializados que el análisis de tipos y frecuencias no revela directamente es la tendencia de ciertas unidades a actuar como \textbf{núcleos léxicos productivos}: términos que, además de funcionar de manera autónoma, aparecen incrustados como componentes de cadenas terminológicas más largas dentro del propio inventario. Esta propiedad ---denominada aquí \textit{incrustación terminológica}--- tiene consecuencias tanto para la evaluación de la representatividad del inventario como para la ponderación estadística de los candidatos, y constituye precisamente el fenómeno que la métrica \textit{C-value} está diseñada para formalizar.

Para el conjunto de términos validados de ambos subcorpus se ha calculado, para cada unidad $t$ del inventario, cuántos términos más largos del mismo inventario la contienen como subsecuencia contigua de tokens (sus \emph{formas contenedoras}), qué fracción de las ocurrencias de $t$ en el corpus corresponde a apariciones autónomas (\emph{standalone}) y qué fracción corresponde a apariciones incrustadas en una cadena más larga (\emph{nested}).

\paragraph{Inventario en inglés.}
De los 134 términos del subcorpus inglés con valores estadísticos completos, 10~unidades (7,5\,\%) presentan al menos una forma contenedora en el inventario; las 124~restantes (92,5\,\%) no participan como componente de ningún otro término validado. Sin embargo, estas 10~unidades concentran el 16,4\,\% del total de ocurrencias terminológicas del subcorpus, por lo que su impacto en los cómputos de frecuencia es desproporcionadamente alto. La Tabla~\ref{tab:terminos_incrustados_en} recoge las unidades con al menos una forma contenedora, ordenadas por número de formas contenedoras.

\begin{table}[H]
\centering
\small
\setlength{\tabcolsep}{4pt}
\renewcommand{\arraystretch}{1.15}
\caption{Términos del subcorpus inglés con al menos una forma contenedora en el inventario validado. \textit{Frec.\ total} = ocurrencias absolutas en el corpus; \textit{Frec.\ autónoma} = ocurrencias en que el término aparece de forma independiente; \textit{\% incrustada} = proporción de ocurrencias incrustadas en una cadena más larga; \textit{FC} = número de formas contenedoras en el inventario.}
\label{tab:terminos_incrustados_en}
\begin{tabular}{@{} l r r r r r @{}}
\toprule
\textbf{Término} & \textbf{Frec.\ total} & \textbf{Frec.\ autónoma} & \textbf{\% incrustada} & \textbf{FC} & \textbf{C-value} \\
\midrule
\textit{corpus}       & 13\,767 &  9\,335 & 32\,\% & 28 & 10\,480{,}0 \\
\textit{annotation}   &  1\,513 &  1\,129 & 25\,\% &  7 &  1\,520{,}1 \\
\textit{frequency}    &  1\,667 &  1\,352 & 19\,\% &  6 &  1\,909{,}9 \\
\textit{type}         &  1\,048 &    782  & 25\,\% &  1 &  1\,978{,}6 \\
\textit{keyword}      &    337  &    198  & 41\,\% &  1 &    651{,}5 \\
\textit{national corpus} & 189 &     75  & 60\,\% &  1 &    194{,}2 \\
\textit{treebank}     &    105  &     64  & 39\,\% &  1 &     40{,}9 \\
\textit{concordance}  &    161  &    145  & 10\,\% &  1 &    207{,}4 \\
\textit{lob}          &    304  &    163  & 46\,\% &  1 &      2{,}4 \\
\textit{parsing}      &     54  &     49  &  9\,\% &  1 &     40{,}1 \\
\bottomrule
\end{tabular}
\end{table}

Los datos ponen de manifiesto patrones de diferente naturaleza. En un extremo, \textit{corpus} actúa como el núcleo más productivo del dominio con diferencia: 28~de los términos validados lo contienen como componente (desde \textit{reference corpus} o \textit{spoken corpus} hasta \textit{British National Corpus} o \textit{diachronic corpus}), y el 32\,\% de sus 13\,767~ocurrencias son incrustadas ---es decir, corresponden a contextos en que el término forma parte de una cadena denominativa más larga---. En un extremo opuesto, \textit{lob} ilustra con especial claridad el efecto de la incrustación sobre la estimación de la terminologicidad: con 304~ocurrencias absolutas y frecuencia de documento de~11, su frecuencia bruta lo situaría entre los candidatos de peso medio del inventario; sin embargo, el 46\,\% de sus apariciones corresponden a la forma contenedora \textit{LOB Corpus}, y su C-value resultante es de~2{,}4, el más bajo entre todos los términos con formas contenedoras. La desconexión entre frecuencia absoluta y C-value en este caso no es un artefacto del sistema, sino la consecuencia directa y correcta del descuento por incrustación: \textit{lob} es terminológicamente poco autónomo porque, en la mayoría de sus contextos de uso, su referencia queda absorbida por una unidad denominativa más precisa.

\paragraph{Inventario en español.}
En el subcorpus español, la incrustación terminológica es considerablemente menos frecuente: solo 5~de los 157~términos validados (3,2\,\%) presentan al menos una forma contenedora, y el conjunto de sus ocurrencias incrustadas representa el 9,0\,\% del total de ocurrencias terminológicas del subcorpus. La Tabla~\ref{tab:terminos_incrustados_es} recoge estas unidades.

\begin{table}[H]
\centering
\small
\setlength{\tabcolsep}{4pt}
\renewcommand{\arraystretch}{1.15}
\caption{Términos del subcorpus español con al menos una forma contenedora en el inventario validado. Mismas definiciones de columna que la Tabla~\ref{tab:terminos_incrustados_en}.}
\label{tab:terminos_incrustados_es}
\begin{tabular}{@{} l r r r r r @{}}
\toprule
\textbf{Término} & \textbf{Frec.\ total} & \textbf{Frec.\ autónoma} & \textbf{\% incrustada} & \textbf{FC} & \textbf{C-value} \\
\midrule
\textit{corpus}               & 11\,749 & 10\,256 & 13\,\% & 26 & 11\,691{,}6 \\
\textit{lingüística de corpus}&    690  &    677  &  2\,\% &  1 &  1\,073{,}0 \\
\textit{monitor}              &     44  &     33  & 25\,\% &  1 &     33{,}0 \\
\textit{Trados}               &     24  &     13  & 46\,\% &  1 &     13{,}0 \\
\textit{Transit}              &      5  &      3  & 40\,\% &  1 &      3{,}0 \\
\bottomrule
\end{tabular}
\end{table}

La asimetría entre ambos subcorpus es lingüísticamente significativa. El inglés forma términos especializados de manera preferente mediante composición sintagmática ---secuencias del tipo \texttt{N\,N} o \texttt{(ADJ|N)*\,N}--- que generan cadenas terminológicas de extensión variable sobre un núcleo nominal compartido; el resultado natural son constelaciones de términos que comparten componentes (todas las variantes de \textit{X corpus} o \textit{corpus X}). El español, en cambio, lexicaliza los mismos conceptos en buena medida mediante derivación morfológica (\textit{lematización}, \textit{tokenización}, \textit{representatividad}) o mediante sintagmas preposicionales (\textit{etiquetado de corpus}, \textit{frecuencia de documento}) que no generan la misma relación de incrustación entre términos del inventario. Así, aunque \textit{corpus} es también el núcleo más productivo del subcorpus español con 26~formas contenedoras, su tasa de incrustación es del 13\,\% ---frente al 32\,\% en inglés---, lo que refleja que la diversificación denominativa en español se apoya menos en la composición nominal y más en la expansión preposicional y la derivación.

\paragraph{Incrustación y C-value.}
La relación entre incrustación terminológica y C-value no es accidental: el C-value \autocite{frantzi1999c} fue diseñado precisamente para abordar este problema. Su fórmula descuenta de la frecuencia bruta de un candidato las ocurrencias que corresponden a su aparición como subcadena de términos más largos, ponderadas por el número de esas formas contenedoras:

\begin{equation}
C\text{-value}(t) =
\begin{cases}
\log_2 |t| \cdot f(t) & \text{si } t \text{ no tiene formas contenedoras,} \\[4pt]
\log_2 |t| \cdot \left( f(t) - \dfrac{1}{|\mathcal{C}(t)|}\displaystyle\sum_{b \in \mathcal{C}(t)} f(b) \right) & \text{en caso contrario,}
\end{cases}
\label{eq:cvalue_incrustacion}
\end{equation}

donde $|t|$ es la longitud en tokens, $f(t)$ la frecuencia bruta, $\mathcal{C}(t)$ el conjunto de formas contenedoras y $f(b)$ la frecuencia de cada una de ellas. El análisis de incrustación presentado en esta sección es, por tanto, la manifestación empírica de lo que el C-value computa internamente: un término con alta frecuencia bruta pero bajo C-value es, casi sin excepción, una unidad que aparece mayoritariamente como componente de términos más largos, como ilustra el caso de \textit{lob} en inglés (frecuencia bruta~304, C-value~2{,}4) o de \textit{Trados} en español (frecuencia bruta~24, C-value~13{,}0, absorbida en un 46\,\% por \textit{SDL Trados Studio}).

En suma, la incrustación terminológica constituye una propiedad discreta pero cuantitativamente relevante de los inventarios especializados: no afecta al mayor número de tipos del vocabulario, pero concentra su efecto en los núcleos conceptuales más productivos del dominio ---precisamente aquellos con mayor capacidad generativa para la formación de cadenas denominativas nuevas---. Su formalización mediante el C-value permite que la puntuación de terminologicidad final (§\ref{sec:extraccion}) refleje la autonomía referencial de cada unidad y no solo su peso distribucional bruto.

% ── Modelo formal del término (OE1.3) ──────────────────────────────────────

\paragraph{El modelo teórico-computacional del término.}
Los resultados de la extracción permiten articular de forma explícita el \textbf{modelo teórico-computacional del término} al que apunta el objetivo específico~OE1.3: una representación formal que integra los rasgos identificados en una abstracción operacionalizable. En el marco de esta investigación, una unidad léxica $u$ se caracteriza como término a través de tres dimensiones complementarias:

\begin{enumerate}
    \item \textbf{Dimensión formal — \textit{unithood}.} $u$ satisface al menos uno de los patrones morfosintácticos de la gramática terminológica $\mathcal{G}$ (Tabla~\ref{tab:patrones-pos}), garantizando que posee la estructura sintagmática adecuada para actuar como unidad denominativa:
    \begin{equation}
        \mathrm{unithood}(u) = \mathbf{1}\!\left[\exists\, p \in \mathcal{G} : u \models p\right]
        \label{eq:unithood}
    \end{equation}

    \item \textbf{Dimensión estadística — \textit{termhood}.} La terminologicidad de $u$ en el corpus se cuantifica mediante el índice compuesto (ec.~\ref{eq:terminologicidad}), que combina el peso distribucional TF-IDF con la autonomía sintagmática capturada por el C-value:
    \begin{equation}
        T(u) = 0{,}6\cdot\widetilde{\mathrm{TF\text{-}IDF}}(u) + 0{,}4\cdot\widetilde{C\text{-value}}(u)
        \label{eq:termhood}
    \end{equation}
    siendo $T(u) > \theta_{90}$ ---el percentil~90 de la distribución--- la condición de umbral para la validación experta.

    \item \textbf{Dimensión distribucional — gradiente de terminologicidad.} Los rasgos distribucionales del lema de $u$ en el corpus ---en especial su diversidad contextual ($\mathrm{CD}$), su categoría gramatical y su frecuencia de documento--- conforman un vector de propiedades que el clasificador supervisado aprende a asociar con la etiqueta de término independientemente de la identidad léxica. El análisis de importancia (Tablas~\ref{tab:bio_importancia_en} y~\ref{tab:bio_importancia_es}) muestra que $\mathrm{CD}(u)$ es el predictor primario en ambas lenguas, resultado que se confirma también en el holdout ($\rho_s \geq 0{,}35$, $p < 0{,}001$): los términos con mayor diversidad de contextos discursivos son sistemáticamente más detectables como unidades terminológicas incluso cuando nunca han sido etiquetados durante el entrenamiento.
\end{enumerate}

La combinación de las tres dimensiones define una caracterización del término como unidad que (i)~cumple restricciones formales estructurales, (ii)~posee un peso distribucional estadísticamente marcado en el corpus y (iii)~aparece en una variedad suficientemente heterogénea de contextos discursivos. Esta triple condición articula computacionalmente la concepción terminológica contemporánea del término como unidad \emph{formal}, \emph{funcional} y \emph{contextual} \cite{cabre1999terminology}, y constituye el modelo teórico-computacional que sustenta tanto la herramienta TECLA como los experimentos de extracción supervisada del capítulo.

La sección siguiente traslada el foco desde la caracterización del inventario terminológico hacia su detección automática: ¿en qué medida pueden los clasificadores supervisados aprender, a partir de rasgos lingüísticos y estadísticos, qué propiedades hacen que una unidad léxica sea terminológica, y cuánto de ese aprendizaje es genuinamente transferible a términos que el modelo nunca ha visto etiquetados?

\section{TECLA: \emph{Term Extraction and Corpus Linguistics Application}}
\label{sec:tecla}

La librería TECLA (\textit{Term Extraction and Corpus Linguistics Application}) se articula como una estación de trabajo integral para lingüistas y especialistas en terminología que deseen abordar de forma sistemática el análisis terminológico a partir de la explotación de corpus. Su diseño responde a la necesidad de contar con un entorno unificado que reúna en un único sistema todas las fases del trabajo terminológico automatizado: desde la compilación y descripción del corpus hasta la extracción de candidatos, la evaluación de la representatividad terminológica y la detección supervisada de unidades terminológicas. El objetivo es proporcionar una herramienta flexible, transparente y accesible que combine metodologías lingüísticas y computacionales sin renunciar a los principios de interpretabilidad, trazabilidad y control por parte del investigador. Frente a estrategias fragmentadas o centradas exclusivamente en criterios estadísticos, TECLA integra tanto enfoques basados en reglas lingüísticas como métricas distribucionales, permitiendo articular flujos completos de trabajo terminológico con salida replicable y personalizable.

\subsection{Arquitectura}
\label{subsec:arquitectura_tecla}

La arquitectura de TECLA se organiza en cuatro módulos funcionales interoperables, que pueden emplearse de forma independiente o encadenada dentro de un flujo de trabajo completo. A continuación se describe cada módulo con el nivel de detalle necesario para evaluar tanto sus capacidades como sus limitaciones metodológicas.

\subsubsection*{Módulo 1. Gestión del corpus}

Este módulo cubre las operaciones previas a la extracción terminológica propiamente dicha: la constitución, limpieza, enriquecimiento y descripción del corpus de trabajo. Acepta como entrada documentos en formato \texttt{PDF} o texto plano (\texttt{.txt}); en el primer caso, realiza la conversión a texto plano mediante \texttt{pdfminer.six} y deposita los ficheros resultantes en un directorio de trabajo estructurado (\texttt{\_tecla\_txt/raw/} para el texto sin procesar, \texttt{\_tecla\_txt/clean/} para el texto limpiado).

La \textbf{detección de lengua} puede configurarse manualmente (\texttt{es}/\texttt{en}) o delegarse al módulo de detección automática basado en \texttt{langid}, que identifica la lengua dominante de cada documento y selecciona el modelo spaCy correspondiente (\texttt{es\_core\_news\_md} para el español; \texttt{en\_core\_web\_lg} para el inglés).

El \textbf{módulo de limpieza} implementa más de una docena de operaciones configurables de forma independiente mediante banderas booleanas, lo que garantiza que nada se aplica implícitamente y que el usuario tiene control total sobre cada transformación:

\begin{itemize}
    \item \textit{Ruido estructural}: reunión de palabras partidas por guión al final de línea; eliminación de números de página y líneas muy cortas.
    \item \textit{Metadatos documentales}: eliminación de títulos de sección aislados (\textit{Abstract}, \textit{Introducción}, \textit{Methodology}…), líneas de palabras clave, pies de figura y tabla, líneas de URL/DOI/email, marcadores de copyright y licencias, datos de publicación de revista (volumen, número, páginas), marcadores de cita numérica en línea (\texttt{[1]}, \texttt{[2,3]}) y líneas mayoritariamente numéricas (tablas estadísticas).
    \item \textit{Sección de referencias}: eliminación automática del bloque bibliográfico al final del documento.
    \item \textit{Filtrado léxico}: eliminación de palabras vacías (combinando las listas de spaCy y NLTK con una stoplist adicional personalizable por el usuario) y lematización opcional mediante spaCy.
\end{itemize}

La \textbf{inferencia y edición de metadatos} permite asociar a cada documento información bibliográfica (autor, año, tema, fuente, notas libres) tanto por inferencia automática a partir del nombre de archivo como mediante edición manual en la interfaz gráfica. Los metadatos se almacenan en un fichero \texttt{metadata.csv} vinculado al directorio de trabajo.

Las \textbf{estadísticas descriptivas} se calculan tanto a nivel de documento como de corpus global e incluyen: número de tokens, número de tipos, relación tipo-token (TTR), proporción de hápax, densidad léxica, número de oraciones y longitud media de oración. Todos estos indicadores pueden exportarse a CSV para su análisis externo.

\subsubsection*{Módulo 2. Extracción de candidatos a término}

Este módulo implementa el núcleo metodológico del sistema de extracción terminológica no supervisada. Sigue una arquitectura híbrida en dos etapas secuenciales:

\paragraph{Filtro lingüístico.} La primera etapa aplica un conjunto de patrones morfosintácticos definidos en un fichero de gramática (\texttt{grammar.yaml}) sobre los documentos procesados por spaCy. Cada patrón especifica una secuencia de etiquetas Universal Dependencies (UD) que debe satisfacer un candidato para ser considerado estructuralmente plausible como término (\textit{unithood}). Los patrones son configurables: el sistema incluye una gramática de referencia para inglés y español, pero el usuario puede substituirla por una gramática propia adaptada a su dominio. La Tabla~\ref{tab:patrones_extraccion} recoge los patrones incorporados, con su descripción estructural y ejemplos.

\begin{table}[H]
\centering
\small
\begin{tabular}{@{} l l p{4.5cm} p{4cm} @{}}
\toprule
\textbf{ID} & \textbf{L} & \textbf{Patrón UD} & \textbf{Ejemplos} \\
\midrule
\texttt{EN\_NP\_PRENOMINAL}  & EN & \texttt{(ADJ|NOUN|PROPN)* NOUN}   & \textit{balanced corpus}, \textit{written corpus} \\
\texttt{EN\_NP\_OF}          & EN & \texttt{NOUN+ of NOUN+}           & \textit{part of speech}, \textit{analysis of variance} \\
\texttt{EN\_NP\_PROPN}       & EN & \texttt{PROPN PROPN+}             & \textit{Brown Corpus}, \textit{WordNet 3.0} \\
\texttt{EN\_NOUN\_SEQUENCE}  & EN & \texttt{NOUN NOUN+}               & \textit{token frequency}, \textit{corpus size} \\
\midrule
\texttt{ES\_NP\_POSTNOMINAL} & ES & \texttt{NOUN ADJ+}                & \textit{corpus paralelo}, \textit{frecuencia léxica} \\
\texttt{ES\_NP\_DE}          & ES & \texttt{NOUN+ de NOUN+}           & \textit{etiquetado de corpus}, \textit{lista de palabras} \\
\texttt{ES\_NP\_DE\_ADJ}     & ES & \texttt{NOUN+ de NOUN ADJ+}       & \textit{análisis de varianza estadística} \\
\texttt{ES\_NP\_PRENOMINAL}  & ES & \texttt{(ADJ|NOUN|PROPN)* NOUN}   & \textit{nueva lingüística computacional} \\
\texttt{ES\_NP\_PROPN}       & ES & \texttt{PROPN PROPN+}             & \textit{corpus CREA}, \textit{BNC} \\
\texttt{ES\_NP\_NOUN\_PROPN} & ES & \texttt{NOUN PROPN+}              & \textit{corpus CORDE} \\
\bottomrule
\end{tabular}
\caption{Patrones morfosintácticos incorporados en la gramática de TECLA para inglés (EN) y español (ES). Las etiquetas POS siguen el estándar Universal Dependencies. El operador \texttt{*} indica cero o más ocurrencias; \texttt{+} indica una o más.}
\label{tab:patrones_extraccion}
\end{table}

Adicionalmente, el filtro aplica dos mecanismos de depuración: un \textbf{filtro de palabras funcionales} que rechaza candidatos que contienen determinantes, pronombres, auxiliares, conjunciones o interjecciones en cualquier posición; y un \textbf{filtro de artefactos} que elimina segmentos producidos por errores de tokenización (palabras partidas por guión que no se reunificaron) y candidatos de alta frecuencia que aparecen en un único documento (ruido de índices o metadatos repetidos).

\paragraph{Filtro estadístico.} La segunda etapa enriquece la lista de candidatos ya depurada con métricas de terminologicidad (\textit{termhood}). Para cada candidato se calculan:

\begin{itemize}
    \item \textbf{Frecuencia absoluta} ($f$), \textbf{frecuencia de documento} (número de documentos en que aparece) y \textbf{frecuencia relativa} (ocurrencias por mil tokens).
    \item \textbf{TF-IDF} (frecuencia de término multiplicada por el logaritmo inverso de la frecuencia de documento): mide hasta qué punto una unidad es característica del corpus frente a ser omnipresente.
    \item \textbf{C-value} \autocite{frantzi1999c}: métrica diseñada específicamente para términos multipalabra que penaliza a los candidatos cortos que aparecen frecuentemente como subcadenas de candidatos más largos; formaliza la noción de que un candidato frecuente incrustado en muchos términos más largos tiene menos autonomía terminológica.
    \item \textbf{NC-value} \autocite{frantzi2000automatic}: extensión contextual del C-value que incorpora la información proporcionada por las palabras de contexto (±2 tokens) de cada ocurrencia del candidato. Cuanto más frecuentemente aparece una palabra de contenido como contexto de múltiples candidatos, mayor peso recibe, lo que permite discriminar candidatos terminológicos en función de su entorno discursivo habitual.
\end{itemize}

Los tres índices numéricos (TF-IDF, C-value, NC-value) se normalizan independientemente al intervalo $[0,1]$ mediante normalización min-max. La puntuación de \textbf{terminologicidad} final es un promedio ponderado de TF-IDF$_{\text{norm}}$ y NC-value$_{\text{norm}}$:

\begin{equation}
\text{terminologicidad} = \alpha \cdot \text{TF-IDF}_{\text{norm}} + (1-\alpha) \cdot \text{NC-value}_{\text{norm}}
\label{eq:terminologicidad_tecla}
\end{equation}

El parámetro $\alpha$ determina el perfil de puntuación y puede ajustarse mediante tres presets predefinidos: \textit{Frecuencia} ($\alpha=0{,}6$), orientado a términos con alta presencia en el corpus; \textit{Equilibrado} ($\alpha=0{,}5$), que pondera por igual ambas métricas; y \textit{Complejidad} ($\alpha=0{,}2$), que prioriza los términos multipalabra con estructuras complejas y bien delimitadas por el contexto.

\paragraph{Postprocesamiento.} Antes de presentar el inventario final al usuario, el sistema aplica tres operaciones adicionales:

\begin{itemize}
    \item \textbf{Detección de variantes}: las variantes ortográficas (diferencias de mayúsculas, guionización o espaciado, como \textit{machine learning} / \textit{machine-learning}) se agrupan bajo la forma canónica con mayor terminologicidad; las siglas cuyas iniciales coinciden con las palabras de un término más largo se vinculan bidireccionalmente (p.\,ej., \textit{NLP} ↔ \textit{natural language processing}).
    \item \textbf{Realce de núcleos}: los unigramas que actúan como núcleo sintáctico de varios términos multipalabra reciben un incremento proporcional de su puntuación de terminologicidad (hasta un máximo del 35\,\%), con el fin de compensar la pérdida de señal de los sustantivos nucleares de alta productividad terminológica.
    \item \textbf{Filtrado de términos genéricos}: una lista curada de expresiones académicas y metodológicas que superan los filtros morfosintácticos pero no denotan conceptos especializados (\textit{research method}, \textit{resultados obtenidos}, etc.) se elimina antes de presentar los resultados.
\end{itemize}

Los parámetros de extracción configurables por el usuario incluyen: longitud máxima del candidato (1--15 tokens), frecuencia mínima (1--50), umbral de percentil (\textit{p}25, \textit{p}50, \textit{p}75, \textit{p}90) y umbral de score mínimo. La salida puede exportarse en formato \textbf{CSV} (con todos los campos numéricos: término, lema, variantes, patrón, frecuencias, TF-IDF, C-value, NC-value, terminologicidad) o en formato \textbf{TBX-Basic} (ISO 30042), el estándar de intercambio terminológico utilizado por herramientas como SDL MultiTerm, OmegaT y Terminator, lo que facilita la integración del inventario en flujos de trabajo terminográficos profesionales.

\subsubsection*{Módulo 3. Evaluación de la representatividad terminológica}

Este módulo permite al investigador responder a una pregunta previa a cualquier tarea terminológica: \emph{¿es este corpus suficientemente representativo del dominio para extraer de él un inventario terminológico fiable?} Para ello implementa el Índice de Representatividad Terminológica (ITR), ya descrito en detalle en el Capítulo~\ref{chap:3_extraccion}, que sintetiza cuatro métricas complementarias en una puntuación única en la escala 0--100:

\begin{itemize}
    \item \textbf{ITC\textsubscript{norm}} (curva de cobertura incremental terminológica): mide la velocidad con que el corpus acumula nuevos tipos de término al añadir documentos sucesivos; un corpus saturado muestra una curva que se aplana rápidamente.
    \item \textbf{MATTR} (\textit{moving average type-token ratio}): estima la riqueza y variedad léxica del corpus mediante una ventana deslizante que evita el efecto de longitud del TTR clásico.
    \item \textbf{Exponente de Heaps} ($\beta$): parámetro de la ley de Heaps que describe la velocidad de crecimiento del vocabulario con el tamaño del corpus; valores bajos de $\beta$ indican alta densidad terminológica.
    \item \textbf{CMN} (cobertura media normalizada): estimación bootstrapped de qué proporción del inventario terminológico total del dominio cubre ya el corpus.
\end{itemize}

La escala de interpretación del ITR es la siguiente: $\geq80$ (muy representativo), $65$--$79$ (representativo), $50$--$64$ (moderadamente representativo), $35$--$49$ (poco representativo) y $<35$ (escasamente representativo).

Además del ITR global, el módulo proporciona dos análisis complementarios de particular interés metodológico:

\begin{itemize}
    \item \textbf{Dispersión de Gries (DP)}: para cada candidato a término se calcula su coeficiente de dispersión entre los documentos del corpus \autocite{gries2021dispersion}. Un DP próximo a 0 indica que el término aparece de forma uniforme en todos los documentos (señal distribucional robusta); un DP próximo a 1 indica concentración en uno o pocos documentos, lo que puede apuntar a un término de autor o a ruido. Los 50 candidatos con mayor terminologicidad se muestran en una tabla ordenada por DP, lo que permite identificar candidatos problemáticos antes de la validación manual.
    \item \textbf{Keyness G\textsuperscript{2}}: análisis de contraste léxico de los candidatos del corpus analizado frente a un corpus de referencia externo, cargado opcionalmente por el usuario. Para cada candidato se calcula el estadístico $G^2$ de log-verosimilitud \autocite{dunning1993accurate} y el log-ratio, lo que permite identificar qué unidades son genuinamente características del dominio especializado y cuáles podrían pertenecer al léxico general o a otro subdominio.
\end{itemize}

El módulo permite asimismo cargar una \textbf{lista de referencia terminológica} externa (p.\,ej., un glosario publicado, una base de datos de terminología como IATE, o el propio gold standard del investigador) y calcular qué proporción de esa lista aparece entre los candidatos extraídos, produciendo automáticamente métricas de precisión y recall respecto al estándar externo.

\subsubsection*{Módulo 4. Clasificación supervisada}

Este módulo permite entrenar clasificadores automáticos de detección terminológica a partir de un corpus anotado mediante la codificación BIO. El flujo metodológico incluye la extracción de rasgos lingüísticos y distribucionales para cada token del corpus, la organización de los datos de entrenamiento y evaluación mediante particiones cruzadas por documento, y la aplicación de algoritmos de aprendizaje supervisado. Los detalles del espacio de rasgos, los algoritmos implementados, la estrategia de evaluación y los resultados obtenidos en el corpus del presente trabajo se describen en detalle en la sección anterior (\S\ref{sec:modelo_tc}). En el contexto del flujo de trabajo de TECLA, este módulo ocupa la posición final de la cadena: el \textit{Silver Standard} producido por el módulo de extracción sirve como punto de partida para la anotación supervisada, y el clasificador resultante puede aplicarse como un filtro adicional de alta precisión sobre nuevos corpus del mismo dominio.

\paragraph{Disponibilidad en el release.} El módulo de clasificación supervisada, incluyendo la implementación de redes CRF y BiLSTM-CRF, forma parte integral del código fuente de TECLA y fundamenta los resultados de evaluación del presente trabajo. Sin embargo, su distribución en el release público se ha circunscrito al código fuente documentado sin incluir los pesos preentrenados de los modelos. Esta decisión responde a consideraciones prácticas: (i) los modelos neuronales BiLSTM entrenan sobre GPU y requieren dependencias adicionales (PyTorch, CRFsuite compilado) que aumentan significativamente el peso de instalación y la complejidad de configuración; (ii) el público objetivo de la herramienta (investigadores en lingüística, terminólogos, lingüistas computacionales sin experiencia en deep learning) puede extraer terminología con alta calidad mediante los módulos 1--3 sin necesidad de entrenar clasificadores; (iii) la calibración de estos modelos es dominio-dependiente y requiere anotación manual, por lo que un release con modelos preentrenados sería engañoso en su generalización a nuevos domios. El release estable de TECLA se enfoca, por tanto, en consolidar los módulos de extracción no supervisada (módulo 2), evaluación de representatividad (módulo 3) e interfaz de usuario, que ofrecen la máxima relación entre utilidad y accesibilidad. Los investigadores interesados en reproducir los experimentos de clasificación supervisada pueden consultar la documentación del repositorio para instrucciones de instalación de dependencias opcionales y reproducción del pipeline completo.

\subsubsection*{Interfaz gráfica y línea de comandos}

TECLA ofrece dos modos de interacción: una interfaz gráfica de escritorio y una interfaz de línea de comandos.

La \textbf{interfaz gráfica} (GUI) está construida con \texttt{customtkinter} y organiza el flujo de trabajo en cuatro pestañas:

\begin{enumerate}
    \item \textbf{Corpus}: muestra una tabla con todos los documentos del corpus (identificador, tokens, oraciones, tipos, TTR, autor, año, tema) y un panel de edición de metadatos por documento. Incluye exportación de estadísticas a CSV.
    \item \textbf{Extracción}: registra el proceso de extracción en tiempo real en un log y presenta la tabla de candidatos a término ordenable por cualquier métrica (score, TF-IDF, C-value, NC-value, frecuencia absoluta, frecuencia relativa, número de documentos). Las columnas visibles son configurables. Permite exportar la selección a CSV o TBX-Basic.
    \item \textbf{Evaluación}: muestra el score ITR con escala visual interpretada, la curva de crecimiento terminológico, la tabla de dispersión DP para los 50 candidatos más prominentes y el análisis de keyness G\textsuperscript{2} frente a un corpus de referencia externo (opcional). Calcula y muestra la cobertura de la extracción respecto a una lista de referencia terminológica externa (opcional).
    \item \textbf{Revisión}: módulo de validación manual que permite al investigador o terminógrafo clasificar cada candidato como \textit{aprobado}, \textit{rechazado} o \textit{pendiente}, añadir notas libres por término, filtrar la lista por estado de revisión y consultar las estadísticas completas de cada candidato (score, métricas individuales, variantes, patrón morfosintáctico). Esta pestaña implementa la interfaz de trabajo propia del terminógrafo, integrando la labor experta en el mismo entorno que el resto del análisis.
\end{enumerate}

La sesión se persiste automáticamente entre ejecuciones (fichero \texttt{.tecla\_session.json}), de modo que el investigador puede retomar el trabajo de extracción y revisión exactamente en el punto en que lo dejó.

La \textbf{interfaz de línea de comandos} (CLI), invocable mediante \texttt{tecla <subcomando>}, expone tres operaciones principales: \texttt{extract} (extracción completa con opciones de filtrado, exportación CSV y TBX), \texttt{stats} (estadísticas descriptivas del corpus) y \texttt{evaluate} (comparación de la extracción frente a un gold standard con cálculo de precisión, recall y F1). La CLI permite integrar TECLA en flujos de trabajo automatizados o en entornos sin interfaz gráfica.

La Tabla~\ref{tab:dependencias_tecla} recoge las dependencias del paquete.

\begin{table}[H]
\centering
\fontsize{10}{12}\selectfont
\begin{tabular}{@{}p{3.5cm}p{8cm}p{2cm}@{}}
\toprule
\textbf{Librería} & \textbf{Funcionalidad principal} & \textbf{Versión} \\
\midrule
\texttt{spaCy} & Tokenización, lematización y etiquetado morfosintáctico (POS) & $\geq$3.0 \\
\texttt{pdfminer.six} & Extracción de texto a partir de documentos PDF & -- \\
\texttt{langid} & Detección automática de la lengua del corpus & -- \\
\texttt{nltk} & Listas de palabras vacías complementarias & -- \\
\texttt{pyyaml} & Carga de la gramática POS desde fichero YAML & -- \\
\texttt{pandas} & Manipulación y exportación de datos tabulares & $\geq$2.0 \\
\texttt{numpy} & Operaciones numéricas y normalización de scores & $\geq$1.24 \\
\texttt{matplotlib} & Visualización de curvas de crecimiento y métricas & $\geq$3.7 \\
\texttt{tqdm} & Barra de progreso para el procesamiento de textos & $\geq$4.65 \\
\texttt{customtkinter} & Interfaz gráfica de escritorio (GUI) & -- \\
\bottomrule
\end{tabular}
\caption{Dependencias del paquete TECLA. Todas las bibliotecas son de código abierto e instalables mediante \texttt{pip}. Las bibliotecas sin versión mínima especificada son compatibles con cualquier versión estable reciente.}
\label{tab:dependencias_tecla}
\end{table}

\subsection{Evaluación}
\label{subsec:evaluacion_tecla}

La evaluación de TECLA se articula en dos planos complementarios: la verificación de componentes y la validación sobre datos reales del dominio.

A nivel de componente, la librería cuenta con un conjunto de pruebas unitarias que cubren las principales transformaciones de datos y cálculos numéricos de cada módulo: correcta conversión y limpieza de documentos, reproducibilidad de los índices estadísticos (TF-IDF, C-value, ITR), integridad del etiquetado BIO y estabilidad de los resultados de clasificación ante distintos entornos de ejecución. Estas pruebas garantizan que el comportamiento de la librería es predecible y replicable antes de su aplicación sobre cualquier corpus.

A nivel de sistema, la evaluación principal de TECLA coincide con el experimento de clasificación supervisada descrito en la sección anterior~(\S\ref{sec:modelo_tc}): los modelos de detección terminológica implementados en el módulo de clasificación automática se evalúan sobre los subcorpus EN y ES mediante validación cruzada con agrupación por documento. Los resultados confirman la adecuación del espacio de representación definido por TECLA: la regresión logística alcanza un F1 de término de 0,83 en inglés y 0,87 en español, y el BiLSTM-CRF supera el 0,96 en ambas lenguas. El protocolo de reserva de tipos terminológicos (\S\ref{tab:holdout_en}) añade una dimensión crítica a esta evaluación: demuestra que los clasificadores entrenados sin rasgos léxicos directos retienen cierta capacidad de generalización a términos nunca vistos, con un recall de entre 0,37 y 0,50 en inglés y entre 0,37 y 0,45 en español.

Por último, la librería ha sido sometida a una evaluación de usabilidad orientada al perfil de usuario objetivo ---investigadores en lingüística y terminología con conocimientos básicos de Python---. Los resultados confirmaron que la documentación generada con Sphinx y la \textsc{api} de la librería permiten reproducir los experimentos y adaptar los parámetros a corpus propios sin conocimientos especializados en aprendizaje automático.

\section{Glosario digital}
\label{sec:glosario}

Como resultado del proceso de extracción y validación descrito a lo largo de este capítulo, se ha elaborado un glosario terminológico bilingüe del dominio de la lingüística de corpus, que constituye uno de los productos finales del presente trabajo. El glosario reúne 155 tipos terminológicos en inglés y 146 en español, validados manualmente a partir de la lista de candidatos generada por TECLA y contrastados con fuentes de referencia especializadas.

Los términos se organizan en siete categorías temáticas que articulan los principales ámbitos conceptuales de la disciplina: \textit{Fundamentos y disciplinas} (conceptos constitutivos de la lingüística de corpus como campo), \textit{Diseño y compilación de corpus} (terminología asociada a la construcción y tipología de corpus), \textit{Recursos y corpus de referencia} (denominaciones de corpus y recursos de uso generalizado), \textit{Procesamiento y anotación} (terminología del análisis lingüístico computacional), \textit{Estadística, léxico y búsqueda} (métricas, índices y herramientas de consulta), \textit{Herramientas y software} (aplicaciones de análisis de corpus) y \textit{Tecnología y formatos} (estándares de codificación y formatos de intercambio).

El glosario incluye pares de equivalencia bilingüe cuando el término dispone de denominación consolidada en ambas lenguas: por ejemplo, \textit{corpus linguistics} / \textit{lingüística de corpus}, \textit{lemmatisation} / \textit{lematización} o \textit{concordance} / \textit{concordancia}. Cuando un concepto carece de denominación estabilizada en una de las dos lenguas, se registra como laguna terminológica. El glosario identifica 118 lagunas en el español ---términos del inglés sin equivalente consolidado en español, como \textit{corpus-based}, \textit{corpus-driven} o \textit{corpus size}--- y 97 lagunas en el inglés ---términos del español sin equivalente acuñado, como determinadas acepciones de \textit{corpus telemático} o \textit{corpus oportunista}---. Estas asimetrías reflejan el desequilibrio histórico en la producción bibliográfica de la disciplina: la mayor parte de los textos fundacionales y la terminología técnica de referencia fueron acuñados en inglés, de modo que numerosos conceptos metodológicos carecen de traducción estandarizada al español, lo que constituye en sí mismo un dato de interés para el estudio de la dinámica terminológica del campo.

El glosario digital se pone a disposición de la comunidad investigadora como recurso abierto en formato estructurado (\texttt{.csv}), e incluye para cada entrada la forma canónica del término, su categoría temática, los equivalentes bilingües cuando existen y la indicación de laguna cuando procede.

\input{chapters/conclusion}

% ---------- Bibliografía y apéndice ----------
\chapter{Referencias Bibliográficas}

\printbibliography[heading=none]

\input{End Matter Pages/appendix}

\end{document}
